\chapter{Differentiation}
\label{vol02:ch05}

\epigraph{The reckoning of motion is the science of comparing
quantities that change by quantities that vanish, and thereby
recovering ratios that remain.}{A working paraphrase of the
seventeenth-century calculus, in the engineer's vocabulary.}

\chapmeta{Half-life: foundational. Archetypes: optimisation (first
formal introduction), transport (gradient as flux density, informal
preview). Exercise target: 40. Project track: simulation with
analysis. Hazard class: none. Reader path default: standard, with
sections explicitly tagged. Named cases: Patriot missile system,
Dhahran (1991), revisited as a numerical-precision case in the
failure section.}

\noindent
The derivative is the rate at which one quantity changes with
respect to another. Volume I taught the reader to measure quantities
and report them with their uncertainty; this chapter teaches the
reader to compute how a quantity moves. We define the derivative as
the limit of a difference quotient, derive the working rules,
develop the local linearisation that the engineer uses for almost
every quick calculation, and confront the numerical instability
that meets the same idea on a finite-precision machine.

\section{The derivative as rate}\pathtag{core}

A quantity changes. Position changes with time, voltage changes
with current, temperature changes with depth into a wall, the bank
balance changes with the day, the population of a culture changes
with the hour. The change is sometimes fast, sometimes slow,
sometimes uniform, sometimes accelerating. Engineering begins with
measurement; once the measurement is in hand, the next question is
the rate.

\subsection{The motivating example: position and velocity}

A car drives along a straight road. Let $x(t)$ denote its position
at time $t$, with $x$ measured in metres and $t$ in seconds. Two
readings of the position taken at times $t_{1}$ and $t_{2}$, with
$t_{2} > t_{1}$, give a change in position
$\Delta x = x(t_{2}) - x(t_{1})$ over an elapsed time
$\Delta t = t_{2} - t_{1}$. The ratio

\[
\bar{v} = \frac{\Delta x}{\Delta t} = \frac{x(t_{2}) - x(t_{1})}
{t_{2} - t_{1}}
\]

has units of metres per second. We call it the average velocity
between $t_{1}$ and $t_{2}$. The average velocity is a property of
the interval, not of any single instant: a car that travels from
$\SI{0}{\meter}$ to $\SI{100}{\meter}$ between $t_{1} = 0$ and
$t_{2} = 10\,\si{\second}$ has average velocity
$10\,\si{\meter\per\second}$ over that interval, regardless of how
the trip was distributed over the ten seconds. The car may have
been stationary for nine seconds and then sprinted; it may have
travelled at exactly $10\,\si{\meter\per\second}$ throughout; it
may have accelerated and decelerated.

The instantaneous velocity at time $t$ is what the speedometer
reports. It is the rate of change of position at the instant $t$,
not over an interval. To recover it from positional measurements,
the engineer has to take the average velocity over an interval that
contains $t$ and shrink the interval. The shrinking is not a
metaphor; it is a numerical procedure. We measure the position at
$t$ and at $t + \Delta t$ for some small $\Delta t$, compute the
average velocity over $[t, t + \Delta t]$, and ask what this number
approaches as $\Delta t$ becomes small.

For an idealised model, a car whose position obeys
$x(t) = \tfrac{1}{2} a t^{2}$ with constant $a > 0$, the calculation
runs as follows. The average velocity between $t$ and $t + \Delta t$
is

\[
\bar{v} = \frac{x(t + \Delta t) - x(t)}{\Delta t} =
\frac{\tfrac{1}{2} a (t + \Delta t)^{2} - \tfrac{1}{2} a t^{2}}
{\Delta t} = \frac{a t \Delta t + \tfrac{1}{2} a (\Delta t)^{2}}
{\Delta t} = a t + \tfrac{1}{2} a \Delta t.
\]

For any $\Delta t > 0$, the average velocity exceeds $a t$ by
$\tfrac{1}{2} a \Delta t$. As $\Delta t$ shrinks toward zero, the
correction $\tfrac{1}{2} a \Delta t$ shrinks toward zero, and the
average velocity approaches $a t$. The instantaneous velocity at
time $t$ is $v(t) = a t$. We arrived at it by the limiting
procedure: difference, divide, shrink.

The procedure is general. The instantaneous rate of change of any
quantity $y$ with respect to any quantity $x$ is the limit of the
difference quotient $(y(x + \Delta x) - y(x)) / \Delta x$ as
$\Delta x \to 0$, when that limit exists.

\subsection{The formal definition}

\begin{definition}[Derivative]
Let $f$ be a real-valued function of a real variable, defined on
an interval containing the point $x$. The \engterm{derivative} of
$f$ at $x$, written $f'(x)$ or $\dv{f}{x}\big|_{x}$, is

\[
f'(x) = \lim_{h \to 0} \frac{f(x + h) - f(x)}{h},
\]

provided the limit exists. When the limit exists, $f$ is said to
be \engterm{differentiable} at $x$. When $f$ is differentiable at
every point of an interval, $f$ is said to be differentiable on
that interval, and the function $x \mapsto f'(x)$ is the
\engterm{derivative function} of $f$.
\end{definition}

The notation $\dv{f}{x}$, due to Leibniz, captures the structure of
the construction: a small change $\dd f$ in the function divided by
a small change $\dd x$ in the argument. The notation $f'(x)$, due
to Lagrange, is shorter and treats the derivative as a function of
$x$. We use both, choosing the one that reads more naturally in
context.

The phrase "provided the limit exists" carries weight. The
difference quotient is well-defined for any $h \neq 0$, but the
limit as $h \to 0$ may fail to exist for several reasons. The
function may have a jump at $x$, in which case the difference
quotient is unbounded as $h \to 0$ from one side. The function may
have a sharp corner at $x$ (a kink), in which case the limit from
the left and the limit from the right exist but differ. The
function may oscillate so rapidly near $x$ that the difference
quotient never settles. Section~5.8 takes up these failure modes;
for this section the reader assumes the functions in view are
smooth enough that the limit exists.

\begin{archetype}[Optimisation, first formal introduction]
The derivative is the engineer's first tool for choosing among
feasible alternatives. Given a function that maps a design
parameter to a cost, a yield, an efficiency, or any other figure
of merit, the points where the derivative vanishes are the
candidates for the maximum and minimum. Section~5.6 develops
this. The optimisation archetype recurs in every later volume
that asks the engineer to choose: thermal design under cost
constraint, structural design under weight constraint, network
routing under latency constraint, biological dose under toxicity
constraint. The pattern enters here, in one variable, with the
condition $f'(x) = 0$.
\end{archetype}

\subsection{Geometric reading: the slope of the tangent line}

The derivative has a geometric reading. Plot the graph of $y =
f(x)$ in the $xy$-plane. Two distinct points on the graph,
$(x, f(x))$ and $(x + h, f(x + h))$, define a secant line whose
slope is the difference quotient $(f(x + h) - f(x)) / h$. As $h$
shrinks, the secant rotates about the fixed point $(x, f(x))$, and
the second point slides along the graph toward the first. In the
limit $h \to 0$, the secant becomes the tangent line at $(x,
f(x))$, and the difference quotient becomes the slope of the
tangent. The derivative $f'(x)$ is the slope of the tangent line
to the graph at the point $x$.

\begin{figure}[ht]
\centering
\begin{tikzpicture}[x=1cm, y=1cm, font=\small]

% Axes
\draw[->, thick] (-0.3, 0) -- (8.2, 0) node[right] {$x$};
\draw[->, thick] (0, -0.3) -- (0, 5.5) node[above] {$y = f(x)$};

% Curve f(x) = 0.18 (x-1)^2 + 0.6
\draw[thick, engblue] plot[domain=0.3:7.6, samples=80]
  (\x, {0.18*(\x-1)*(\x-1) + 0.6});

% Point of tangency a = 3
\fill[engred] (3, {0.18*4 + 0.6}) circle (2pt);
\node[engred, anchor=north east] at (3, {0.18*4 + 0.6}) {$(a,\,f(a))$};

% Secant 1 (long h)
\draw[thick, engamber, dashed] (1.0, {0.18*0 + 0.6}) -- (7.4, {0.18*40.96 + 0.6});
\fill[engamber] (7.0, {0.18*36 + 0.6}) circle (1.6pt);
\node[engamber, anchor=south west, font=\scriptsize] at (7.0, {0.18*36 + 0.6 + 0.05}) {$h_{1}$ large};

% Secant 2 (medium h)
\draw[thick, enggray, dashed] (2.0, {0.18*1 + 0.6}) -- (6.0, {0.18*25 + 0.6});
\fill[enggray] (5.5, {0.18*20.25 + 0.6}) circle (1.4pt);
\node[enggray, anchor=south west, font=\scriptsize] at (5.5, {0.18*20.25 + 0.6 + 0.05}) {$h_{2}$};

% Tangent line at a = 3: slope = 0.36*(a-1) = 0.72
\draw[thick, engred] (0.5, {0.78 - 0.72*2.5}) -- (6.0, {0.78 + 0.72*3.0});
\node[engred, anchor=west] at (5.7, {0.78 + 0.72*2.7 + 0.2}) {tangent at $a$};

% Bracket showing h shrinking
\draw[<->, thick] (3, -0.5) -- (5.5, -0.5);
\node[anchor=north, font=\scriptsize] at (4.25, -0.5) {$h \to 0$: secants rotate into tangent};

\end{tikzpicture}
\caption[Tangent line as limit of secants]{The difference quotient
$(f(a + h) - f(a))/h$ is the slope of the secant line through
$(a, f(a))$ and $(a + h, f(a + h))$. As $h$ shrinks, the second
point slides along the graph toward $(a, f(a))$ and the secant
rotates about the fixed point. The derivative $f'(a)$ is the slope
of the limiting line, the tangent at $a$.}
\label{fig:vol02:ch05:tangent-secant}
\end{figure}


\Cref{fig:vol02:ch05:tangent-secant} shows the construction. The
secant for a large step misses the tangent's direction by a visible
amount; the secant for a smaller step is closer; the limit recovers
the tangent exactly. The construction is symmetric in the sense
that the same limit is obtained whether the second point approaches
from the right or from the left (when the limit exists). When the
two one-sided limits differ, the function has a kink at the point,
and the tangent line is not defined; section~5.8 takes up that
failure case.

The geometric reading is more than decoration. It tells the
engineer where the function is increasing (positive slope), where
it is decreasing (negative slope), where it is locally flat (zero
slope), and where it is steep (large absolute slope). It also
sets up the local linearisation of section~5.4: a smooth function
near a point looks, to first order, like its tangent line.

\subsection{Units and dimensions}

The derivative inherits its units from the ratio that defines it.
If $f$ has units $[f]$ and $x$ has units $[x]$, then $f'(x)$ has
units $[f]/[x]$. For position $x(t)$ in metres and time $t$ in
seconds, the derivative $\dv{x}{t}$ has units of
$\si{\meter\per\second}$, which is velocity. For temperature
$T(z)$ in degrees Celsius and depth $z$ in metres, the derivative
$\dv{T}{z}$ has units of $\si{\degreeCelsius\per\meter}$, which is
a thermal gradient. For voltage $V(I)$ in volts as a function of
current $I$ in amperes, the derivative $\dv{V}{I}$ has units of
$\si{\volt\per\ampere} = \si{\ohm}$, which is the differential
resistance. The unit-balance discipline of Volume I extends to
calculus: every derivative is a ratio of dimensioned quantities,
and its units must be tracked through every calculation.

\begin{estimation}
\textbf{Question.} A car accelerates from rest under constant
power. The kinetic energy at time $t$ obeys, approximately,
$E(t) = P t$ where $P$ is the power. If the car has mass $m$, the
velocity satisfies $v(t) = \sqrt{2 P t / m}$. Estimate the
acceleration at $t = 5\,\si{\second}$ for a $m = 1500\,\si{\kg}$
car with $P = 100\,\si{\kilo\watt}$, and compare the answer to the
acceleration at $t = 1\,\si{\second}$.

\textbf{Estimate.} At $t = 5\,\si{\second}$ the velocity is
$v = \sqrt{2 \cdot 10^{5} \cdot 5 / 1500} = \sqrt{667} \approx
26\,\si{\meter\per\second}$. The acceleration is
$\dv{v}{t} = \tfrac{1}{2}(2 P / m)^{1/2} t^{-1/2} = (P / (2 m
t))^{1/2}$, so at $t = 5\,\si{\second}$ we expect
$a \approx (10^{5} / (2 \cdot 1500 \cdot 5))^{1/2} = (10^{5} / 1.5
\times 10^{4})^{1/2} \approx 2.6\,\si{\meter\per\second\squared}$.
At $t = 1\,\si{\second}$ the same formula gives $a \approx
(10^{5} / 3000)^{1/2} \approx 5.8\,\si{\meter\per\second\squared}$.
The acceleration falls as $t^{-1/2}$ because the constant-power
constraint forces the car to spend more energy per unit
acceleration as the velocity rises. The estimate matches what a
working driver knows: a constant-power car accelerates fastest at
low speed.
\end{estimation}

\begin{estimation}
\textbf{Question.} A reader computes the central difference
$D_{0}(x; h) = (f(x + h) - f(x - h)) / (2 h)$ for a function whose
typical value is order unity, using double-precision arithmetic
($\varepsilon \approx 10^{-16}$ in each function evaluation).
Estimate the relative error in the derivative as $h$ shrinks from
$10^{-2}$ to $10^{-12}$, and predict the value of $h$ at which the
error is smallest.

\textbf{Estimate.} The two contributions to the error in $D_{0}$
are the truncation bias of order $|f'''| h^{2} / 6$ and the
round-off variance of order $\varepsilon / h$. Set $|f'''| \sim 1$
for an order-of-magnitude estimate. At $h = 10^{-2}$ the truncation
term is $\sim 1.7 \times 10^{-5}$ and the round-off term is
$\sim 10^{-14}$; truncation dominates. At $h = 10^{-12}$ the
truncation term is $\sim 1.7 \times 10^{-25}$ and the round-off
term is $\sim 10^{-4}$; round-off dominates. The minimum total
error occurs near $h^{*} = (\varepsilon)^{1/3} \approx 10^{-5}$
where the two contributions cross. The minimum error is of order
$\varepsilon^{2/3} \approx 2 \times 10^{-11}$. A reader who picks
$h = 10^{-8}$ (the naive choice, half the precision of a
double) lands on the round-off branch and loses about three
decades of accuracy compared with the optimum. The estimate is
borne out by direct computation; the chapter's code repository
contains \texttt{finite\_diff\_accuracy.py}, which tabulates the
sweep for $f(x) = e^{x}$ at $x = 1$.
\end{estimation}

\subsection{Examples computed from the definition}

We compute three derivatives directly from the limit definition.
The point of the exercise is to see the limit do real work; the
working rules of section~5.2 will then let the engineer compute
the same derivatives in two lines.

\engterm{Linear function}. Let $f(x) = a x + b$ with $a, b$
constants. Then $f(x + h) - f(x) = a(x + h) + b - (a x + b) = a h$,
so $(f(x + h) - f(x)) / h = a$. The limit as $h \to 0$ is $a$. The
derivative of a linear function is the slope of the line.

\engterm{Square function}. Let $f(x) = x^{2}$. Then $f(x + h) -
f(x) = (x + h)^{2} - x^{2} = 2 x h + h^{2}$, so $(f(x + h) -
f(x)) / h = 2 x + h$. The limit as $h \to 0$ is $2 x$. The
derivative of $x^{2}$ is $2 x$.

\engterm{Reciprocal}. Let $f(x) = 1/x$ for $x \neq 0$. Then
$f(x + h) - f(x) = 1/(x + h) - 1/x = -h/(x(x + h))$, so
$(f(x + h) - f(x)) / h = -1/(x(x + h))$. The limit as $h \to 0$
is $-1/x^{2}$. The derivative of $1/x$ is $-1/x^{2}$.

The pattern that emerges, called the \engterm{power rule}, is
that the derivative of $x^{n}$ is $n x^{n - 1}$ for integer $n$.
The proof from the definition uses the binomial expansion
$(x + h)^{n} = x^{n} + n x^{n-1} h + \tbinom{n}{2} x^{n-2} h^{2} +
\cdots$; only the linear-in-$h$ term survives the division by $h$
and the limit. The same formula extends to non-integer exponents
once the chain rule and the exponential function are in hand
(section~5.2).

\subsection{Derivatives of the elementary functions}

The reader will use the following derivatives without re-deriving
them in routine calculation. We list them here and prove the
non-obvious cases later in the chapter.

\begin{itemize}[leftmargin=*]
  \item $\dv{}{x}(x^{n}) = n x^{n - 1}$ for any real $n$, with
    $x > 0$ when $n$ is non-integer.
  \item $\dv{}{x}(\sin x) = \cos x$ and $\dv{}{x}(\cos x) = -\sin
    x$, with $x$ in radians.
  \item $\dv{}{x}(e^{x}) = e^{x}$. The exponential is, up to a
    multiplicative constant, the unique function whose derivative
    equals itself.
  \item $\dv{}{x}(\ln x) = 1/x$ for $x > 0$.
  \item $\dv{}{x}(a^{x}) = a^{x} \ln a$ for $a > 0$.
  \item $\dv{}{x}(\log_{a} x) = 1/(x \ln a)$ for $a > 0$, $a \neq
    1$, $x > 0$.
  \item $\dv{}{x}(\tan x) = \sec^{2} x = 1 / \cos^{2} x$.
  \item $\dv{}{x}(\arcsin x) = 1/\sqrt{1 - x^{2}}$ for $|x| < 1$.
  \item $\dv{}{x}(\arctan x) = 1/(1 + x^{2})$.
\end{itemize}

The trigonometric derivatives require $x$ in radians. A reader who
computes $\dv{}{x}(\sin x)$ with $x$ in degrees will recover an
extra factor of $\pi / 180$: in degrees, $\dv{}{x}(\sin x) =
(\pi / 180) \cos x$. Volume I reported angles in radians by
default, and Volume II continues the convention. The radian is the
SI-coherent unit of angle, and calculus on trigonometric functions
is dimensionally clean only when angles are measured in radians.

\subsection{Differentiability is stronger than continuity}

A function can be continuous at a point without being
differentiable there. The implication runs one way: a function
differentiable at $x$ is continuous at $x$, because the existence
of the limit $f'(x) = \lim_{h \to 0} (f(x + h) - f(x))/h$ forces
the numerator $f(x + h) - f(x)$ to vanish as $h \to 0$ (otherwise
the quotient would be unbounded), which is the definition of
continuity. The converse fails. The function $f(x) = |x|$ is
continuous everywhere and differentiable everywhere except at the
origin, where the difference quotient approaches $+1$ from the
right and $-1$ from the left. Continuity asks the function to have
no jumps; differentiability asks the function to have no corners,
which is a stronger demand. Section~5.8 catalogues the ways the
stronger demand fails on functions the engineer meets in practice.

The hierarchy has more levels than two. A function may be once
differentiable but not twice, twice but not three times, and so on.
The notation $C^{0}$ denotes continuous functions, $C^{1}$ denotes
functions with a continuous first derivative, $C^{k}$ denotes
functions with a continuous $k$-th derivative, and $C^{\infty}$
denotes infinitely differentiable (smooth) functions. The
smoothness class of a function controls which numerical methods
apply to it: a finite-difference scheme of order $p$ requires the
function to be in $C^{p+1}$ for its truncation-error analysis to
hold, and a spline interpolant is built to match a stated
smoothness class at its knots. The engineer who differentiates a
$C^{1}$ function with a method that assumes $C^{4}$ gets an answer
whose error bound does not apply.

\begin{remark}
A continuous function can fail to be differentiable not merely at
isolated points but everywhere. The Weierstrass function
$W(x) = \sum_{n=0}^{\infty} a^{n} \cos(b^{n} \pi x)$, with
$0 < a < 1$ and $b$ an odd integer satisfying $a b > 1 + 3\pi/2$,
is continuous on the whole real line and differentiable at no
point. The graph is a fractal: every magnification reveals the
same roughness. The engineer rarely meets a Weierstrass function
directly, but its qualitative behaviour, structure at every scale
with no well-defined slope, is the model for a class of real
signals: turbulent velocity fields, financial price series at fine
time resolution, and rough fracture surfaces. For such signals the
derivative is not a useful summary at any scale, and the engineer
reaches instead for a scaling exponent or a spectral description.
\end{remark}

\subsection{The derivative as a sensitivity}

The engineer reads $f'(x)$ as a sensitivity: the factor by which a
small change in the input $x$ is amplified into a change in the
output $f$. If $\Delta x$ is a small perturbation of the input,
then $\Delta f \approx f'(x) \, \Delta x$ to first order. A large
derivative means the output is sensitive to the input; a small
derivative means the output is insensitive. This reading is the
bridge between the calculus of this chapter and the
uncertainty-propagation discipline of Volume~I.

A worked instance. The period of a simple pendulum is $T(L) =
2\pi \sqrt{L/g}$, where $L$ is the length and $g$ is the local
gravitational acceleration. The sensitivity of the period to the
length is $\dv{T}{L} = \pi / \sqrt{g L} = T / (2 L)$. A relative
change in length produces half the relative change in period:
$\Delta T / T = \tfrac{1}{2} \, \Delta L / L$. A pendulum clock
whose rod lengthens by one part in $10^{4}$ from thermal expansion
gains or loses half a part in $10^{4}$ in its period, which over a
day of $86\,400 \,\si{\second}$ is about $4 \,\si{\second}$. The
derivative converts a dimensioned sensitivity into a working error
budget, and the half-power scaling is read directly off the
exponent in the formula. Volume~I deferred the general
several-variable form of this propagation to the present chapter;
section~5.5 redeems the promise.

\begin{mastery}
The difference-quotient definition is specific to a real function
of a real variable, but the idea it captures, best linear
approximation, generalises far beyond. For a function $f$ between
normed vector spaces, the \engterm{Gateaux derivative} at a point
$x$ in a direction $v$ is the one-sided rate $\lim_{t \to 0}
(f(x + tv) - f(x))/t$, the directional derivative familiar from
several-variable calculus reorganised to allow $x$ and $v$ to live
in a function space rather than in $\R^{n}$. It probes $f$ along one
direction at a time. The stronger \engterm{Frechet derivative}
asks for a single bounded linear operator $A$ such that $f(x + v) =
f(x) + A v + o(\lVert v \rVert)$ uniformly over all directions $v$,
the exact analogue of the tangent-line statement $f(a + h) = f(a) +
f'(a) h + o(h)$ with the scalar $f'(a)$ promoted to a linear map.
When the Frechet derivative exists it equals the Gateaux derivative,
but the converse can fail, just as in finite dimensions a function
can have all directional derivatives at a point without being
differentiable there. The engineering payoff arrives in the
calculus of variations and in optimal control, where the
unknown is a whole function (a trajectory, a control law, a shape)
and the condition for a minimum is that the Frechet derivative of a
cost functional vanishes. The Euler-Lagrange equation of Volume~VII
is exactly that vanishing-derivative condition, written out. The
one-variable derivative of this chapter is the ground floor of a
tower that reaches function spaces; the structure, value plus
linear correction plus higher-order remainder, is the same at every
level.
\end{mastery}

\section{Rules: product, chain, quotient, implicit}\pathtag{core}

The limit definition gives the derivative directly, but a working
engineer rarely returns to the limit. Instead, the engineer uses a
small set of algebraic rules that combine known derivatives into
the derivative of a more complex expression. We derive the rules
here and apply them.

\subsection{Linearity}

Differentiation is linear: if $f$ and $g$ are differentiable at
$x$ and $a, b$ are constants, then $(a f + b g)$ is differentiable
at $x$ and

\[
(a f + b g)'(x) = a f'(x) + b g'(x).
\]

The derivation is direct from the limit definition. The difference
quotient of $a f + b g$ is

\[
\frac{(a f + b g)(x + h) - (a f + b g)(x)}{h} = a \cdot
\frac{f(x + h) - f(x)}{h} + b \cdot \frac{g(x + h) - g(x)}{h}.
\]

Each term on the right has a limit as $h \to 0$, namely $a f'(x)$
and $b g'(x)$, so the sum has a limit, and that limit is $a f'(x)
+ b g'(x)$. Linearity says: differentiate term by term, factor out
constants.

\subsection{The product rule}

\begin{theorem}[Product rule]
If $f$ and $g$ are differentiable at $x$, then so is $f g$, and

\[
(f g)'(x) = f'(x) g(x) + f(x) g'(x).
\]
\end{theorem}

\begin{proof}
We compute the difference quotient of $f g$ at $x$:

\[
\frac{(f g)(x + h) - (f g)(x)}{h} = \frac{f(x + h) g(x + h) -
f(x) g(x)}{h}.
\]

Add and subtract $f(x + h) g(x)$ in the numerator:

\[
= \frac{f(x + h) g(x + h) - f(x + h) g(x) + f(x + h) g(x) -
f(x) g(x)}{h} =
f(x + h) \frac{g(x + h) - g(x)}{h} + g(x) \frac{f(x + h) - f(x)}{h}.
\]

As $h \to 0$, the first factor on the right approaches $f(x)$ (a
differentiable function is continuous, a fact taken from
Chapter~4), the second factor approaches $g'(x)$, and the third
group of terms converges to $g(x) f'(x)$. Adding,

\[
(f g)'(x) = f(x) g'(x) + f'(x) g(x).
\]
\end{proof}

The mnemonic "first times derivative of second, plus second times
derivative of first" captures the formula. The product rule is
not the rule a reader who has met only the linearity might guess;
the derivative of a product is not the product of the derivatives.

A worked check. Let $f(x) = x^{2}$ and $g(x) = x^{3}$. The
product is $f(x) g(x) = x^{5}$, whose derivative by the power
rule is $5 x^{4}$. The product rule gives $(x^{2})' x^{3} + x^{2}
(x^{3})' = 2 x \cdot x^{3} + x^{2} \cdot 3 x^{2} = 2 x^{4} + 3
x^{4} = 5 x^{4}$, agreeing.

\subsection{The chain rule}

\begin{theorem}[Chain rule]
If $g$ is differentiable at $x$ and $f$ is differentiable at
$g(x)$, then the composition $(f \circ g)$ is differentiable at
$x$, and

\[
(f \circ g)'(x) = f'(g(x)) \cdot g'(x).
\]
\end{theorem}

We derive the rule heuristically and indicate where the heuristic
needs care. Let $u = g(x)$ and write $f(g(x)) = f(u)$. A small
change $\Delta x$ in $x$ produces a small change $\Delta u =
g(x + \Delta x) - g(x)$ in $u$, and a small change $\Delta f =
f(u + \Delta u) - f(u)$ in $f$. We have

\[
\frac{\Delta f}{\Delta x} = \frac{\Delta f}{\Delta u} \cdot
\frac{\Delta u}{\Delta x},
\]

provided $\Delta u \neq 0$. As $\Delta x \to 0$, $\Delta u \to 0$
because $g$ is differentiable hence continuous, $\Delta f / \Delta
u \to f'(u) = f'(g(x))$, and $\Delta u / \Delta x \to g'(x)$. The
product on the right approaches $f'(g(x)) g'(x)$, which is the
claim.

The case $\Delta u = 0$ for some sequence of $\Delta x$ values,
which can happen if $g$ is locally constant, requires a more
careful argument; the standard fix replaces the troublesome ratio
$\Delta f / \Delta u$ by a continuous extension. The conclusion
holds in all cases the engineer encounters in practice.

In Leibniz notation the chain rule reads

\[
\dv{f}{x} = \dv{f}{u} \cdot \dv{u}{x},
\]

a form so suggestive that it can be remembered by treating
derivatives as fractions and "cancelling" $\dd u$. The cancellation
is heuristic; the formula it produces is correct.

A worked use. Let $h(x) = \sin(x^{2})$. Setting $u = x^{2}$ and
$f(u) = \sin u$, we have $f'(u) = \cos u$ and $u'(x) = 2 x$, so
$h'(x) = \cos(x^{2}) \cdot 2 x = 2 x \cos(x^{2})$.

Another. Let $h(x) = (3 x^{2} + 1)^{5}$. Setting $u = 3 x^{2} +
1$ and $f(u) = u^{5}$, we have $f'(u) = 5 u^{4}$ and $u'(x) = 6
x$, so $h'(x) = 5 (3 x^{2} + 1)^{4} \cdot 6 x = 30 x (3 x^{2} +
1)^{4}$.

A nested case. Let $h(x) = \sin(\cos(x))$. Two applications of
the chain rule give $h'(x) = \cos(\cos(x)) \cdot (-\sin(x)) =
-\sin(x) \cos(\cos(x))$.

\begin{figure}[p]
\centering
\begin{tikzpicture}[x=1cm, y=1cm, font=\small]

% Top panel: composition diagram
\begin{scope}[shift={(0, 8)}]
  \node[anchor=south] at (6.5, 4.0) {\textbf{Composition $f \circ g$: rates multiply}};

  % x axis box
  \draw[thick] (0, 0) rectangle (3, 1.2);
  \node at (1.5, 0.6) {$x$};
  \node[anchor=south, font=\scriptsize] at (1.5, 1.2) {input};

  % g maps x -> u
  \draw[->, very thick, engblue] (3.2, 0.6) -- (5.0, 0.6)
    node[midway, above, font=\scriptsize, engblue] {$g$}
    node[midway, below, font=\scriptsize, engblue] {rate $g'(x)$};

  % u box
  \draw[thick] (5.2, 0) rectangle (8.2, 1.2);
  \node at (6.7, 0.6) {$u = g(x)$};
  \node[anchor=south, font=\scriptsize] at (6.7, 1.2) {intermediate};

  % f maps u -> y
  \draw[->, very thick, engred] (8.4, 0.6) -- (10.2, 0.6)
    node[midway, above, font=\scriptsize, engred] {$f$}
    node[midway, below, font=\scriptsize, engred] {rate $f'(u)$};

  % y box
  \draw[thick] (10.4, 0) rectangle (13.0, 1.2);
  \node at (11.7, 0.6) {$y = f(u)$};
  \node[anchor=south, font=\scriptsize] at (11.7, 1.2) {output};

  % combined rate
  \node[align=center, font=\small] at (6.5, -1.0)
    {composite rate $\dv{y}{x} = \dv{y}{u} \cdot \dv{u}{x}
       = f'(g(x)) \cdot g'(x)$};

  % small increment arrows
  \draw[<->, engamber] (0.5, -0.4) -- (2.5, -0.4)
    node[midway, below, font=\scriptsize, engamber] {$\Delta x$};
  \draw[<->, engamber] (5.7, -0.4) -- (7.7, -0.4)
    node[midway, below, font=\scriptsize, engamber] {$\Delta u \approx g'(x)\Delta x$};
  \draw[<->, engamber] (10.7, -0.4) -- (12.7, -0.4)
    node[midway, below, font=\scriptsize, engamber] {$\Delta y \approx f'(u)\Delta u$};
\end{scope}

% Bottom panel: worked example sin(x^2)
\begin{scope}[shift={(0, 0)}]
  \node[anchor=south] at (6.5, 5.2) {\textbf{Worked example: $h(x) = \sin(x^{2})$}};

  % Axes
  \draw[->, thick] (-0.3, 1.5) -- (13.0, 1.5) node[right] {$x$};
  \draw[->, thick] (0, 1.0) -- (0, 5.0) node[above] {};

  % x axis tick labels
  \foreach \x/\lbl in {2/0.5, 4/1.0, 6/1.5, 8/2.0, 10/2.5, 12/3.0} {
    \draw (\x, 1.45) -- (\x, 1.55);
    \node[anchor=north, font=\scriptsize] at (\x, 1.45) {\lbl};
  }
  \node[anchor=east, font=\scriptsize] at (-0.1, 4.5) {$+1$};
  \draw (-0.1, 4.5) -- (0.1, 4.5);
  \node[anchor=east, font=\scriptsize] at (-0.1, 2.5) {$0$};
  \node[anchor=east, font=\scriptsize] at (-0.1, 0.5) {$-1$};

  % Plot of sin(x^2) for x in [0, 3], scaled
  % map x in [0,3] to plot x in [0, 12]: plotx = 4*x
  % map y in [-1, 1] to plot y in [0.5, 4.5]: ploty = 2.5 + 2*y
  \draw[thick, engblue] plot[domain=0:3, samples=200]
    ({4*\x}, {2.5 + 2*sin(deg(\x*\x))});

  % Plot of h'(x) = 2x cos(x^2), normalised to fit (divide by 6 to fit in [-1,1] approx)
  % Wait, magnitude at x=3 is 6, at x=2 is ~4 cos(4)*4. We'll cap and clip.
  % Just sample where |h'(x)| < ~2.5 visible region; use scale 1/3
  \draw[thick, engred, dashed] plot[domain=0:2.85, samples=200]
    ({4*\x}, {2.5 + 2*(2*\x*cos(deg(\x*\x)))/4});

  % Legend
  \node[engblue, anchor=west, font=\scriptsize] at (8.5, 4.6) {$h(x) = \sin(x^{2})$};
  \node[engred, anchor=west, font=\scriptsize] at (8.5, 4.2) {$h'(x) = 2x\cos(x^{2})$ (rescaled)};

  % Annotate construction at x = sqrt(pi/2) approx 1.253
  \fill[engblack] ({4*1.253}, {2.5 + 2*1}) circle (1.5pt);
  \node[engblack, anchor=south, font=\scriptsize] at ({4*1.253}, {2.5 + 2*1 + 0.05}) {peak of $h$};

\end{scope}

\end{tikzpicture}
\caption[Chain-rule decomposition and worked example]{Top: the chain
rule reads the composite $y = f(g(x))$ as a two-stage transformation.
A small input increment $\Delta x$ is mapped by $g$ to an intermediate
increment $\Delta u \approx g'(x) \Delta x$, which $f$ in turn maps
to $\Delta y \approx f'(u) \Delta u = f'(g(x)) g'(x) \Delta x$. The
composite rate is the product of the local rates at each stage.
Bottom: $h(x) = \sin(x^{2})$ and its derivative $h'(x) = 2x \cos(x^{2})$
illustrate the chain-rule mechanics. The derivative is large where
both the outer cosine and the inner factor $2x$ are large; the
derivative passes through zero at every peak and trough of $h$.}
\label{fig:vol02:ch05:chain-rule}
\end{figure}


\Cref{fig:vol02:ch05:chain-rule} pictures the composition as a
two-stage transformation and tracks the propagation of a small
increment through both stages. The bottom panel shows the worked
example $h(x) = \sin(x^{2})$ and its derivative; the derivative
passes through zero at every extremum of $h$, and the amplitude
of the derivative grows linearly in $x$ because the inner factor
$2x$ amplifies whatever the outer cosine returns.

\subsection{The quotient rule}

\begin{theorem}[Quotient rule]
If $f$ and $g$ are differentiable at $x$ and $g(x) \neq 0$, then
$f / g$ is differentiable at $x$, and

\[
\left(\frac{f}{g}\right)'(x) = \frac{f'(x) g(x) - f(x) g'(x)}
{g(x)^{2}}.
\]
\end{theorem}

Two derivations. The first uses the product and chain rules. Write
$f / g = f \cdot (1/g)$. The chain rule applied to $1/g$ with outer
function $u \mapsto 1/u$ (whose derivative is $-1/u^{2}$) gives
$(1/g)'(x) = -g'(x) / g(x)^{2}$. The product rule then gives

\[
\left(\frac{f}{g}\right)'(x) = f'(x) \cdot \frac{1}{g(x)} + f(x)
\cdot \frac{-g'(x)}{g(x)^{2}} = \frac{f'(x) g(x) - f(x) g'(x)}
{g(x)^{2}}.
\]

The second derivation is from the limit definition; the algebra is
the routine subtraction-of-fractions manipulation.

A worked use. Let $h(x) = (x^{2} + 1) / (x - 3)$ for $x \neq 3$.
With $f = x^{2} + 1$, $g = x - 3$, $f' = 2 x$, $g' = 1$,

\[
h'(x) = \frac{2 x (x - 3) - (x^{2} + 1) \cdot 1}{(x - 3)^{2}} =
\frac{2 x^{2} - 6 x - x^{2} - 1}{(x - 3)^{2}} = \frac{x^{2} -
6 x - 1}{(x - 3)^{2}}.
\]

\subsection{Implicit differentiation}

Some functional relationships are not given as $y = f(x)$ but as
$F(x, y) = 0$. The circle $x^{2} + y^{2} = 1$ is the canonical
example: it does not define $y$ as a single-valued function of
$x$, because for each $x \in (-1, 1)$ there are two values of $y$
satisfying the relation. The engineer who wants the slope of the
tangent to the circle at a particular point still wants
$\dv{y}{x}$ at that point. \engterm{Implicit differentiation}
recovers it without solving for $y$ explicitly.

The technique: differentiate both sides of $F(x, y) = 0$ with
respect to $x$, treating $y$ as a function of $x$, and apply the
chain rule wherever $y$ appears. For $x^{2} + y^{2} = 1$,

\[
2 x + 2 y \dv{y}{x} = 0 \implies \dv{y}{x} = -\frac{x}{y}.
\]

At the point $(x_{0}, y_{0}) = (3/5, 4/5)$ on the unit circle, the
slope of the tangent is $-3/4$. The technique works wherever the
relation defines $y$ as a differentiable function of $x$ in a
neighbourhood of the point of interest, which the implicit
function theorem (chapter~11 of this volume) makes precise.

A second example, where solving for $y$ explicitly would be
algebraically painful: $x \sin y + y \cos x = e^{x y}$.
Differentiating both sides with respect to $x$, applying the
product rule on each term and the chain rule on $\sin y$, $\cos
x$, and $e^{x y}$, gives

\[
\sin y + x \cos y \dv{y}{x} + \dv{y}{x} \cos x - y \sin x =
e^{x y} \left(y + x \dv{y}{x}\right).
\]

Solving for $\dv{y}{x}$:

\[
\dv{y}{x} (x \cos y + \cos x - x e^{x y}) = y e^{x y} + y \sin x
- \sin y \implies \dv{y}{x} = \frac{y e^{x y} + y \sin x - \sin y}
{x \cos y + \cos x - x e^{x y}}.
\]

The slope of the tangent at any specific point is recovered by
substituting the point's coordinates. The expression is unwieldy,
but it is also exact, and the engineer who needs the slope at a
single point reads it off in one substitution.

\subsection{Related rates}

Implicit differentiation with respect to time turns a geometric or
physical constraint into a relation between rates. If two
quantities are linked by an equation that holds at every instant,
then differentiating the equation with respect to time relates the
rate of change of one quantity to the rate of change of the other.
The technique is the workhorse of kinematics, fluid filling and
draining, and any setting where one rate is measured and another
is wanted.

\engterm{The sliding ladder}. A ladder of fixed length $L$ rests
against a vertical wall. Let $x(t)$ be the horizontal distance from
the wall to the foot of the ladder and $y(t)$ the height of the
top on the wall. The Pythagorean constraint holds at every instant:
$x(t)^{2} + y(t)^{2} = L^{2}$. Differentiating both sides with
respect to time,

\[
2 x \dv{x}{t} + 2 y \dv{y}{t} = 0 \implies \dv{y}{t} = -\frac{x}{y}
\dv{x}{t}.
\]

If the foot is pulled away from the wall at a constant
$\dv{x}{t} = 0.5 \,\si{\meter\per\second}$, then when the foot is
$x = 4 \,\si{\meter}$ from a wall with a $L = 5 \,\si{\meter}$
ladder, the top is at $y = 3 \,\si{\meter}$ and descends at
$\dv{y}{t} = -(4/3)(0.5) \approx -0.67 \,\si{\meter\per\second}$.

\begin{figure}[ht]
\centering
\begin{tikzpicture}[x=1cm, y=1cm, font=\small]

% Related rates: a ladder sliding down a wall. Right triangle with
% horizontal foot x, vertical height y, hypotenuse L fixed.
% Show the ladder at one instant and the velocity arrows.

% Wall and floor
\draw[very thick, enggray] (0, 0) -- (0, 5.2);   % wall
\draw[very thick, enggray] (0, 0) -- (6.2, 0);    % floor
\node[enggray, anchor=east, font=\scriptsize] at (-0.1, 4.8) {wall};
\node[enggray, anchor=north, font=\scriptsize] at (5.6, -0.1) {floor};

% Ladder: foot at (x0, 0), top at (0, y0), length L
\def\xo{3.6}
\def\yo{3.0}   % so L = 4.69
\draw[thick, engblue] (\xo, 0) -- (0, \yo);
\node[engblue, anchor=south west, font=\scriptsize] at (1.9, 1.7) {ladder, length $L$};

% Right-angle mark
\draw[enggray] (0.25,0) -- (0.25,0.25) -- (0,0.25);

% Labels x and y
\draw[<->, engamber] (0, -0.5) -- (\xo, -0.5);
\node[engamber, anchor=north, font=\scriptsize] at ({\xo/2}, -0.5) {$x(t)$};
\draw[<->, engamber] (-0.5, 0) -- (-0.5, \yo);
\node[engamber, anchor=east, font=\scriptsize, rotate=90] at (-0.55, {\yo/2}) {$y(t)$};

% Velocity arrows: foot moves right (dx/dt > 0), top moves down (dy/dt < 0)
\draw[->, very thick, engred] (\xo, 0) -- (\xo + 1.1, 0);
\node[engred, anchor=west, font=\scriptsize] at (\xo + 1.1, 0) {$\dot{x}>0$};
\draw[->, very thick, engred] (0, \yo) -- (0, \yo - 1.1);
\node[engred, anchor=south west, font=\scriptsize] at (0.05, \yo - 1.1) {$\dot{y}<0$};

% Constraint annotation
\node[engblack, anchor=west, font=\scriptsize, align=left] at (4.6, 3.6)
  {$x^{2}+y^{2}=L^{2}$\\[2pt] $x\dot{x}+y\dot{y}=0$};

\end{tikzpicture}
\caption[Related rates: the sliding ladder]{A ladder of fixed length
$L$ rests against a wall; the foot slides out at rate $\dot{x}$ and
the top slides down at rate $\dot{y}$. The geometric constraint
$x^{2}+y^{2}=L^{2}$ is differentiated implicitly with respect to
time to give $x\dot{x}+y\dot{y}=0$, which relates the two rates:
$\dot{y} = -(x/y)\,\dot{x}$. As the foot approaches the wall's base
($y\to 0$) the top accelerates downward without bound in the
idealised model, a kinematic singularity that real ladders avoid by
leaving the wall. Related-rates problems differentiate a constraint
to convert a known rate into an unknown one.}
\label{fig:vol02:ch05:related-rates}
\end{figure}


\Cref{fig:vol02:ch05:related-rates} pictures the constraint and the
two rates. The qualitative content is in the factor $-x/y$: as the
top nears the floor ($y \to 0$) the descent rate diverges in the
idealised model, a kinematic singularity that real ladders avoid by
separating from the wall before the limit. The accompanying script
\texttt{code/related\_rates.py} tabulates the descent rate as the
foot approaches the base.

\engterm{The draining tank}. A conical tank of half-angle $\alpha$
drains so that the water volume falls at a known rate $\dv{V}{t} =
-Q$. The volume of water at depth $h$ is $V = \tfrac{1}{3} \pi
(h \tan\alpha)^{2} h = \tfrac{1}{3} \pi \tan^{2}\alpha \, h^{3}$.
Differentiating with respect to time, $\dv{V}{t} = \pi \tan^{2}
\alpha \, h^{2} \dv{h}{t}$, so the surface descends at $\dv{h}{t} =
-Q / (\pi \tan^{2}\alpha \, h^{2})$. The descent accelerates as the
depth falls, because the cross-sectional area shrinks with depth in
a cone. An engineer sizing an overflow or a drain reads the worst
case off this relation: the fastest level change happens at the
shallowest depth.

\engterm{The inflating balloon}. A spherical balloon is inflated at
a constant volume rate $\dv{V}{t} = R$. With $V = \tfrac{4}{3} \pi
r^{3}$, differentiation gives $\dv{V}{t} = 4 \pi r^{2} \dv{r}{t}$,
so the radius grows at $\dv{r}{t} = R / (4 \pi r^{2})$. The radial
expansion slows as the balloon grows, because the same volume rate
is spread over a larger surface. The surface area $A = 4 \pi r^{2}$
itself grows at $\dv{A}{t} = 8 \pi r \dv{r}{t} = 2 R / r$. Three
rates (volume, radius, area) are linked by one geometric relation
and two differentiations.

\subsection{Inverse functions}

If $y = f(x)$ is differentiable and invertible near a point, with
inverse $x = f^{-1}(y)$, the derivative of the inverse is the
reciprocal of the derivative of the original, evaluated at the
matching point:

\[
\dv{}{y} f^{-1}(y) = \frac{1}{f'(x)}, \qquad x = f^{-1}(y),
\]

valid wherever $f'(x) \neq 0$. The result is the chain rule applied
to the identity $f(f^{-1}(y)) = y$: differentiating both sides with
respect to $y$ gives $f'(f^{-1}(y)) \cdot (f^{-1})'(y) = 1$. The
geometric reading is that the graph of $f^{-1}$ is the reflection
of the graph of $f$ across the line $y = x$, and reflection swaps
rise and run, inverting the slope.

A worked use derives the derivative of $\arcsin$. Let $y =
\arcsin x$, so $x = \sin y$ with $y \in (-\pi/2, \pi/2)$. Then
$\dv{x}{y} = \cos y$, and $\dv{y}{x} = 1/\cos y$. Since $\cos y =
\sqrt{1 - \sin^{2} y} = \sqrt{1 - x^{2}}$ (positive on the chosen
branch), $\dv{}{x} \arcsin x = 1/\sqrt{1 - x^{2}}$, the entry in
the table of section~5.1. The same route gives $\dv{}{x} \arctan x
= 1/(1 + x^{2})$ from $x = \tan y$, $\dv{x}{y} = \sec^{2} y = 1 +
\tan^{2} y = 1 + x^{2}$.

\begin{mastery}
Implicit differentiation and the inverse-function formula both presume
that the relation they differentiate actually defines a differentiable
function near the point of interest. Two theorems supply the licence and
say precisely when it is granted. The \engterm{inverse function
theorem} states that if $f$ is continuously differentiable near $a$ and
$f'(a) \neq 0$, then $f$ is invertible on a neighbourhood of $a$, the
inverse is continuously differentiable, and its derivative at $b = f(a)$
is $1/f'(a)$, the reciprocal formula of this section made rigorous. The
hypothesis $f'(a) \neq 0$ is exactly the condition that the tangent line
is not horizontal; where $f'(a) = 0$ the inverse has a vertical tangent
and the reciprocal formula divides by zero, the analytic signature of a
local failure of invertibility. The \engterm{implicit function theorem}
is the several-variable companion: if $F(x, y) = 0$ holds at a point and
$\partial F/\partial y \neq 0$ there, then the relation defines $y$ as a
differentiable function of $x$ near that point, with
$\dv{y}{x} = -(\partial F/\partial x)/(\partial F/\partial y)$, which is
the implicit-differentiation result of this section read as a ratio of
partial derivatives. The condition $\partial F/\partial y \neq 0$ is the
several-variable analogue of $f'(a) \neq 0$: it fails exactly at the
points where the level curve $F = 0$ has a vertical tangent and $y$
cannot be solved for as a single-valued function of $x$, the points on
the unit circle $x^{2} + y^{2} = 1$ where $y = 0$ and the slope $-x/y$
of this section blows up. Both theorems are the formal backstop of the
continuation method in the chapter's equation-of-state case study, which
tracks a root only as long as the relevant derivative stays nonzero
along the path; the moment it vanishes, the root path turns vertical and
the tracking must change variables. Chapter~11 develops the
several-variable statement; the present box names the theorems so the
reader knows that the techniques of this section rest on a guarantee with
a checkable hypothesis.
\end{mastery}

\subsection{Parametric differentiation}

A curve in the plane is often given parametrically, as $x = x(t)$
and $y = y(t)$ with $t$ a parameter (often time). The slope of the
curve $\dv{y}{x}$ is recovered from the two parametric derivatives
by the chain rule:

\[
\dv{y}{x} = \frac{\dv{y}{t}}{\dv{x}{t}}, \qquad \dv{x}{t} \neq 0.
\]

A worked use traces a projectile. With $x(t) = v_{0} \cos\theta
\cdot t$ and $y(t) = v_{0} \sin\theta \cdot t - \tfrac{1}{2} g
t^{2}$, the parametric derivatives are $\dv{x}{t} = v_{0}
\cos\theta$ and $\dv{y}{t} = v_{0} \sin\theta - g t$. The path
slope is $\dv{y}{x} = (v_{0} \sin\theta - g t)/(v_{0} \cos\theta)$,
which is zero at the top of the trajectory ($t = v_{0}
\sin\theta / g$) and negative on the descent. The second
derivative along the path, needed for curvature, follows from
differentiating $\dv{y}{x}$ with respect to $t$ and dividing again
by $\dv{x}{t}$:

\[
\frac{\dd^{2} y}{\dd x^{2}} = \frac{\dv{}{t}\!\left(\dv{y}{x}
\right)}{\dv{x}{t}}.
\]

For the projectile this gives $\frac{\dd^{2} y}{\dd x^{2}} = -g /
(v_{0} \cos\theta)^{2}$, a constant: the trajectory is a parabola,
whose second derivative in $x$ is the same everywhere, as the
algebra of $y$ as an explicit quadratic in $x$ confirms.

\subsection{Logarithmic differentiation}

Implicit differentiation has a useful special case for products,
quotients, and powers: take the logarithm first. If $y = f(x)
g(x) / h(x)$, then $\ln y = \ln f(x) + \ln g(x) - \ln h(x)$.
Differentiating both sides with respect to $x$ and applying the
chain rule on the left,

\[
\frac{1}{y} \dv{y}{x} = \frac{f'(x)}{f(x)} + \frac{g'(x)}{g(x)} -
\frac{h'(x)}{h(x)},
\]

so

\[
\dv{y}{x} = y \left(\frac{f'(x)}{f(x)} + \frac{g'(x)}{g(x)} -
\frac{h'(x)}{h(x)}\right) = \frac{f g}{h} \left(\frac{f'}{f} +
\frac{g'}{g} - \frac{h'}{h}\right).
\]

The technique converts a product-quotient mess into a sum of
logarithmic derivatives. It is the engineer's preferred way to
differentiate expressions like $y = x^{x}$, where neither the
power rule nor the exponential rule applies directly. With $\ln y
= x \ln x$, differentiating gives $(1/y) \dv{y}{x} = \ln x + 1$,
so $\dv{y}{x} = x^{x} (\ln x + 1)$.

\subsection{Working tactics}

A reader confronted with a derivative they do not recognise asks
three questions in order. First, is this a sum or scalar multiple
of simpler pieces? If yes, apply linearity and reduce to the
pieces. Second, is this a product, quotient, or composition? If
yes, apply the product, quotient, or chain rule. Third, is this
implicit, or does it have a structure that logarithmic
differentiation simplifies? If yes, take the logarithm or
differentiate implicitly.

After the structural decomposition, every leaf of the
decomposition is one of the elementary derivatives listed in
section~5.1. The composite derivative reassembles by applying the
combination rules backward.

\section{Higher-order derivatives}\pathtag{core}

The derivative of $f$ is itself a function of $x$, and one may
ask for its derivative.

\begin{definition}[Higher-order derivative]
The \engterm{second derivative} of $f$, written $f''(x)$ or
$\frac{\dd^{2} f}{\dd x^{2}}$, is the derivative of $f'(x)$:

\[
f''(x) = \dv{}{x} f'(x) = \lim_{h \to 0} \frac{f'(x + h) -
f'(x)}{h}.
\]

The \engterm{$n$-th derivative} $f^{(n)}(x)$ or $\frac{\dd^{n}
f}{\dd x^{n}}$ is defined recursively as the derivative of the
$(n - 1)$-th derivative.
\end{definition}

The first three derivatives of position $x(t)$ have engineering
names. The first is velocity $v(t) = \dv{x}{t}$. The second is
acceleration $a(t) = \dv{v}{t} = \frac{\dd^{2} x}{\dd t^{2}}$.
The third is sometimes called jerk, the rate of change of
acceleration, relevant in vehicle ride comfort and in aerospace
attitude control. Beyond the third, the named derivatives become
specialised.

For polynomials, repeated differentiation eventually produces
zero. The function $f(x) = a_{n} x^{n} + a_{n - 1} x^{n - 1} +
\cdots + a_{0}$ has $f^{(n)}(x) = n! a_{n}$, a constant, and
$f^{(k)}(x) = 0$ for $k > n$. For non-polynomial functions, the
higher derivatives may be nonzero at every order. The exponential
$f(x) = e^{x}$ has $f^{(k)}(x) = e^{x}$ for every $k$.
Trigonometric functions cycle: $\sin^{(0)} = \sin$,
$\sin^{(1)} = \cos$, $\sin^{(2)} = -\sin$, $\sin^{(3)} = -\cos$,
$\sin^{(4)} = \sin$, repeating with period four.

\subsection{The second derivative reads concavity}

The second derivative carries geometric information about the
graph of $f$. Where $f''(x) > 0$, the graph is curving upward (the
graph lies above its tangent line in a neighbourhood of $x$); we
say $f$ is \engterm{concave up} or \engterm{convex} at $x$. Where
$f''(x) < 0$, the graph is curving downward; we say $f$ is
\engterm{concave down} or \engterm{concave} at $x$. Where $f''(x)
= 0$ and $f''$ changes sign, the graph has an \engterm{inflection
point}.

\begin{figure}[ht]
\centering
\begin{tikzpicture}[x=1cm, y=1cm, font=\small]

% A cubic-like curve showing concave-up, inflection, concave-down,
% with the tangent lying below (convex) and above (concave) the curve.
% f(x) = (1/6) x^3 - (1/2) x  on roughly [-3, 3], shifted up.

\def\f(#1){(0.16*(#1)*(#1)*(#1) - 0.5*(#1) + 2.0)}

% Axes
\draw[->, thick] (-3.4, 0) -- (3.6, 0) node[right] {$x$};
\draw[->, thick] (0, 0) -- (0, 4.2) node[above] {$y = f(x)$};

% Curve
\draw[thick, engblue] plot[domain=-3.2:3.2, samples=100]
  (\x, {\f(\x)});

% Inflection at x = 0 for this cubic
\fill[engblack] (0, {\f(0)}) circle (1.8pt);
\node[engblack, anchor=south east, font=\scriptsize] at (-0.1, {\f(0)+0.05}) {inflection $f''=0$};

% Concave-down region (x < 0): tangent above curve. Pick xa = -2.
\def\xa{-2.0}
\pgfmathsetmacro{\fa}{0.16*\xa*\xa*\xa - 0.5*\xa + 2.0}
\pgfmathsetmacro{\fpa}{0.48*\xa*\xa - 0.5}
\draw[thick, engred] (\xa-1.0, {\fa - \fpa}) -- (\xa+1.0, {\fa + \fpa});
\fill[engred] (\xa, {\fa}) circle (1.6pt);
\node[engred, anchor=south, font=\scriptsize] at (\xa, {\fa+0.45}) {$f''<0$ (concave)};

% Concave-up region (x > 0): tangent below curve. Pick xb = 2.
\def\xb{2.0}
\pgfmathsetmacro{\fb}{0.16*\xb*\xb*\xb - 0.5*\xb + 2.0}
\pgfmathsetmacro{\fpb}{0.48*\xb*\xb - 0.5}
\draw[thick, engamber] (\xb-1.0, {\fb - \fpb}) -- (\xb+1.0, {\fb + \fpb});
\fill[engamber] (\xb, {\fb}) circle (1.6pt);
\node[engamber, anchor=north, font=\scriptsize] at (\xb, {\fb-0.4}) {$f''>0$ (convex)};

\end{tikzpicture}
\caption[The second derivative reads concavity]{Where $f'' > 0$ the
graph curves upward and lies above its tangent line (convex); where
$f'' < 0$ the graph curves downward and lies below its tangent line
(concave). The transition point, where $f''$ changes sign, is an
inflection point. The convexity reading is the working tool of
one-variable optimisation: a critical point in a convex region is a
minimum, a critical point in a concave region is a maximum, and the
second-derivative test of section~5.6 formalises the picture.}
\label{fig:vol02:ch05:concavity}
\end{figure}


\Cref{fig:vol02:ch05:concavity} draws the reading directly: in a
convex region the tangent line lies below the curve, in a concave
region it lies above, and the inflection point is where the two
regimes meet. The engineer who carries this picture reads the sign
of $f''$ off the shape of the graph without computing it, and reads
the shape of the graph off the sign of $f''$ when only the formula
is in hand.

The convexity reading is the engineer's working tool for
optimisation. A function whose second derivative is positive
everywhere has at most one local minimum, which is also its
global minimum, on a convex domain. A function whose second
derivative is negative everywhere has at most one local maximum,
which is also its global maximum. We return to this in
section~5.6 and treat it formally in chapter~15 of this volume,
where convex optimisation is the working framework for many
engineering design problems.

\subsection{Higher derivatives carry function shape}

Beyond the second, higher derivatives carry progressively finer
information about the function's shape near a point. The third
derivative tells us how the curvature is changing: a negative
third derivative at a point of zero second derivative is the
signature of a peak in curvature. The fourth derivative carries
information about how the curvature changes its rate of change.
For a function known to be smooth (infinitely differentiable),
the entire sequence of derivatives at a single point determines
the function in a neighbourhood of that point, via the Taylor
series of section~5.4.

For an engineering reader, the higher derivatives matter most in
three places. First, in numerical methods, where higher
derivatives appear in the truncation-error terms of finite
differences and quadrature rules, and where the rate of
convergence depends on which derivatives are bounded. Second, in
local linearisation and Taylor approximation, where the engineer
chooses how many terms to keep based on the desired accuracy.
Third, in the dynamics of physical systems, where conservation
laws and material constants relate the second derivative of a
field to forcing terms (Newton's second law, the wave equation,
the heat equation, the Schrödinger equation). The reader will
meet the higher derivatives most often through these three uses.

\subsection{Curvature and the radius of curvature}

The second derivative reads concavity through its sign. Its
magnitude, combined with the slope, reads how sharply the curve
bends. The \engterm{curvature} $\kappa$ of the graph of $y = f(x)$
at a point is the rate at which the tangent direction turns per unit
arc length travelled along the curve. A straight line never turns
and has zero curvature; a tight circle turns quickly and has large
curvature. The formula, derived by differentiating the tangent
angle $\theta(x) = \arctan f'(x)$ with respect to arc length
$s$, is

\[
\kappa(x) = \frac{|f''(x)|}{\left(1 + f'(x)^{2}\right)^{3/2}}.
\]

The derivation runs through the chain rule and the arc-length
element. The tangent angle satisfies $\tan\theta = f'(x)$, so
$\dv{\theta}{x} = f''(x)/(1 + f'(x)^{2})$ by the derivative of
$\arctan$. The arc length grows at $\dv{s}{x} = \sqrt{1 +
f'(x)^{2}}$, the Pythagorean combination of the horizontal step
$\dd x$ and the vertical step $f'(x)\,\dd x$. Curvature is the turn
per arc length, $\kappa = |\dv{\theta}{s}| = |\dv{\theta}{x}| /
\dv{s}{x}$, and dividing one expression by the other gives the
formula. The dimensionless-looking result carries units of inverse
length, because $\theta$ is dimensionless (an angle in radians) and
$s$ is a length: curvature has units of $\si{\per\meter}$.

The reciprocal of the curvature is the \engterm{radius of
curvature} $R = 1/\kappa$, the radius of the circle that best fits
the curve at the point. That circle, the \engterm{osculating
circle}, shares the curve's position, tangent, and second
derivative at the point of contact.

\begin{figure}[ht]
\centering
\begin{tikzpicture}[x=1cm, y=1cm, font=\small]

% Parabola y = 0.35 x^2 on x in [-2.6, 2.6], with osculating circles at
% the vertex (large radius implied by small curvature) and at a flank
% (smaller radius). Plot scale: x by 1.4, y by 1.4.

% Axes
\draw[->, thick] (-4.2, 0) -- (4.4, 0) node[right] {$x$};
\draw[->, thick] (0, -0.4) -- (0, 4.6) node[above] {$y$};

% Parabola: plotx = 1.4*x, ploty = 1.4*(0.35*x*x)
\draw[very thick, engblue] plot[domain=-2.6:2.6, samples=80]
  ({1.4*\x}, {1.4*0.35*\x*\x});

% Osculating circle at the vertex (x=0). For y = a x^2, kappa(0) = 2a,
% radius R = 1/(2a) = 1/0.7 approx 1.43 in data units; centre at (0, R).
% Scaled: centre (0, 1.4*1.43) = (0, 2.0), radius 1.4*1.43 = 2.0.
\draw[thick, engred, dashed] (0, 2.0) circle (2.0);
\fill[engred] (0, 2.0) circle (1.4pt);
\node[engred, anchor=west, font=\scriptsize] at (0.15, 2.0)
  {centre, vertex};
\fill[engblack] (0, 0) circle (1.6pt);
\node[engblack, anchor=north east, font=\scriptsize] at (-0.1, -0.05)
  {$R_{0} = 1/(2a)$};

% Osculating circle at a flank point x = 1.8. y = 0.35*1.8^2 = 1.134.
% y' = 0.7*1.8 = 1.26; y'' = 0.7. kappa = |y''|/(1+y'^2)^{3/2}
%   = 0.7/(1+1.5876)^{1.5} = 0.7/4.158 = 0.1684; R = 5.94 data units.
% This is large; to keep the figure compact we draw a representative
% arc near the point rather than the full circle.
% Tangent direction at x=1.8: slope 1.26, unit normal points up-left.
% Normal direction (data): (-1.26, 1)/|.| = (-0.783, 0.622).
% Centre (data) = (1.8, 1.134) + R*(-0.783, 0.622)
%   = (1.8 - 5.94*0.783, 1.134 + 5.94*0.622) = (-2.85, 4.83).
% Scaled centre: (1.4*-2.85, 1.4*4.83) = (-3.99, 6.76); radius 1.4*5.94=8.3.
% Draw only a short arc passing through the scaled flank point
% (1.4*1.8, 1.4*1.134) = (2.52, 1.588).
\fill[engblack] (2.52, 1.588) circle (1.6pt);
\node[engblack, anchor=west, font=\scriptsize] at (2.6, 1.45)
  {flank point};
\draw[thick, engorange, dashed]
  ([shift={(-3.99,6.76)}] 230:8.3) arc (230:255:8.3);
\node[engorange, anchor=south west, font=\scriptsize] at (2.7, 1.9)
  {$R_{1} > R_{0}$};

% Tangent at flank point
\draw[thin, enggray] (2.52 - 1.2*1.4, 1.588 - 1.2*1.26)
  -- (2.52 + 0.7*1.4, 1.588 + 0.7*1.26);

\node[engblue, anchor=west, font=\scriptsize] at (-4.0, 4.3)
  {$y = a x^{2}$};

\end{tikzpicture}
\caption[Radius of curvature and the osculating circle]{The
osculating circle at a point on $y = a x^{2}$ shares the curve's
position, tangent, and second derivative there. At the vertex the
curvature $\kappa = 2a$ is largest and the radius $R_{0} = 1/(2a)$
is smallest; the circle hugs the curve tightly. At the flank point
the slope $y'$ is no longer zero, the denominator
$(1 + (y')^{2})^{3/2}$ in the curvature formula grows, the curvature
drops, and the matching radius $R_{1} > R_{0}$ widens. A vehicle
tracing the curve at constant speed feels a lateral acceleration
$v^{2}/R$ that peaks where the radius is smallest.}
\label{fig:vol02:ch05:curvature}
\end{figure}


\Cref{fig:vol02:ch05:curvature}
draws the osculating circles at two points of a parabola: tight
where the parabola bends sharply, wide where it is nearly straight.

\engterm{Worked instance: a banked road}. A vehicle travelling at
speed $v$ along a curved path of radius of curvature $R$ feels a
lateral acceleration $a_{\text{lat}} = v^{2}/R = v^{2} \kappa$. The
curvature of the road, computed from its design profile $y(x)$,
sets the lateral acceleration directly. Consider a road whose
centreline near a bend is modelled by $y(x) = \beta x^{2}$ with
$\beta = 2 \times 10^{-3}\,\si{\per\meter}$, evaluated at the
vertex $x = 0$ where the road is steepest in curvature. There
$f'(0) = 0$ and $f''(0) = 2\beta = 4 \times 10^{-3}\,\si{\per
\meter}$, so $\kappa = 4 \times 10^{-3}\,\si{\per\meter}$ and $R =
250\,\si{\meter}$. A car at $v = 25\,\si{\meter\per\second}$
(about $90\,\si{\kilo\meter\per\hour}$) feels $a_{\text{lat}} =
625/250 = 2.5\,\si{\meter\per\second\squared}$, roughly a quarter
of $g$, a firm but manageable cornering load. The same calculation
sets the design speed of highway curves: the geometry fixes
$\kappa$, the comfort or friction limit fixes the maximum
tolerable $a_{\text{lat}}$, and the design speed follows from
$v_{\max} = \sqrt{a_{\text{lat,max}}/\kappa}$. The curvature is a
second-derivative quantity, and the road geometry that the
surveyor lays out is the integral, twice over, of the
acceleration the driver will feel.

\engterm{Worked instance: peak curvature of a Gaussian}. The
curvature of $f(x) = e^{-x^{2}/2}$ varies along the bell. At the
peak $x = 0$, $f'(0) = 0$ and $f''(0) = -1$, so $\kappa(0) = 1$
and the osculating circle has radius $1$. Moving outward, the
slope grows and the denominator $(1 + f'^{2})^{3/2}$ inflates,
while the second derivative $f''(x) = (x^{2} - 1)e^{-x^{2}/2}$
changes sign at $x = \pm 1$. At the inflection points $x = \pm 1$
the curvature passes through zero (the second derivative
vanishes), so the curve is locally straight there, and the
osculating circle degenerates to a line. A reader plotting the
curvature against $x$ finds a peak at the centre, a fall to zero
at the inflection points, and a small secondary structure in the
tails. The curvature, not the function value, is the quantity an
optical designer tracks when a Gaussian beam profile must match a
lens figure.

\subsection{The Leibniz rule for repeated products}

The product rule iterates. Differentiating a product $n$ times
produces a binomial sum, the \engterm{Leibniz rule}:

\[
(f g)^{(n)} = \sum_{k=0}^{n} \binom{n}{k} f^{(k)} g^{(n - k)}.
\]

The structure mirrors the binomial theorem, with derivatives in
place of powers. The first-order case $(f g)' = f' g + f g'$ is the
ordinary product rule; the second-order case is $(f g)'' = f'' g +
2 f' g' + f g''$, the coefficients $1, 2, 1$ being the second row
of Pascal's triangle. The rule earns its keep when one factor has a
short derivative sequence. For $f(x) = x^{2} e^{x}$, the factor
$x^{2}$ has only three nonzero derivatives ($x^{2}$, $2x$, $2$),
so the $n$-th derivative is a three-term sum: $(x^{2} e^{x})^{(n)}
= \binom{n}{0} x^{2} e^{x} + \binom{n}{1} 2x \, e^{x} +
\binom{n}{2} 2 \, e^{x} = e^{x}(x^{2} + 2 n x + n(n-1))$, since
every derivative of $e^{x}$ is $e^{x}$. A reader who needs the
fortieth derivative of $x^{2} e^{x}$ reads it off the formula
rather than differentiating forty times.

\begin{mastery}
The higher-derivative sequence at a single point encodes the
function's local behaviour to all orders. For an infinitely
differentiable function the Taylor series
$f(x) = \sum_{n=0}^{\infty} f^{(n)}(a) (x - a)^{n} / n!$ collects
the entire derivative sequence into one object. The series need not
converge to $f$ even when it converges: the function
$f(x) = e^{-1/x^{2}}$ (with $f(0) = 0$) has every derivative equal
to zero at the origin, so its Taylor series at $0$ is identically
zero, yet the function is positive for every $x \neq 0$. A
function whose Taylor series converges to it in a neighbourhood is
called \engterm{analytic}; the example shows that smooth and
analytic are different classes. The engineering consequence is that
Taylor-series accuracy claims require analyticity, not merely
smoothness, and that a numerical method calibrated on analytic test
functions may behave differently on a merely smooth signal. The
distinction returns in the convergence theory of chapter~16.
\end{mastery}

\section{Local linearisation}\pathtag{core}

The single most useful consequence of differentiability, for
engineering practice, is the local linearisation: a smooth
function near a working point looks, to first order, like its
tangent line. Many engineering calculations replace the function
by its tangent line over a small range of input, and the answer
is correct to within the chapter's stated accuracy. We make the
approximation precise.

\subsection{Taylor's theorem to first order}

\begin{theorem}[Taylor's theorem, first-order remainder form]
If $f$ is differentiable on an interval containing $a$, and $f''$
is continuous on that interval, then for any $x$ in the interval

\[
f(x) = f(a) + f'(a)(x - a) + R_{1}(x; a),
\]

where the remainder satisfies $R_{1}(x; a) = \tfrac{1}{2}
f''(\xi) (x - a)^{2}$ for some $\xi$ between $a$ and $x$.
\end{theorem}

The expression $L_{a}(x) = f(a) + f'(a)(x - a)$ is the
\engterm{linear approximation} or \engterm{tangent-line
approximation} of $f$ at $a$. It is the unique linear function
that agrees with $f$ at $a$ both in value and in slope. The
remainder $R_{1}$ measures the error of the approximation; it is
quadratic in the displacement $(x - a)$, with a coefficient set
by the second derivative at some intermediate point.

A reader who has been carrying error bars from Volume I notices
the structure: the linear approximation is a model, and the
remainder is the error of the model. The error scales as the
square of the displacement and as the magnitude of the second
derivative. A function that is gently curving (small $f''$) admits
a wide working range; a function that curves sharply (large
$f''$) admits a narrow working range.

\begin{figure}[ht]
\centering
\begin{tikzpicture}[x=1cm, y=1cm, font=\small]

% Plot e^x near x = 0, with linear and quadratic Taylor approximations
% Plot domain x in [-1.4, 1.4], scale x by 3, y by 1
% function exp(x) ranges [0.25, 4.06]
% map x in [-1.4, 1.4] to plot [-4.2, 4.2]
% map y in [0, 4] to plot y in [0, 4]

% Axes
\draw[->, thick] (-4.6, 0) -- (4.8, 0) node[right] {$x$};
\draw[->, thick] (0, -0.3) -- (0, 4.5) node[above] {$y$};

% x-axis tick labels
\foreach \x/\lbl in {-3/-1, -1.5/-0.5, 1.5/0.5, 3/1} {
  \draw (\x, 0) -- (\x, -0.1);
  \node[anchor=north, font=\scriptsize] at (\x, -0.1) {\lbl};
}
% y-axis ticks
\foreach \y in {1, 2, 3, 4} {
  \draw (0, \y) -- (-0.1, \y);
  \node[anchor=east, font=\scriptsize] at (-0.1, \y) {\y};
}

% True curve e^x: plotx = 3*x, ploty = exp(x)
\draw[very thick, engblue] plot[domain=-1.4:1.4, samples=80]
  ({3*\x}, {exp(\x)});

% Linear approximation 1 + x (tangent at 0)
\draw[thick, engred, dashed] plot[domain=-1.4:1.4, samples=20]
  ({3*\x}, {1 + \x});

% Quadratic approximation 1 + x + x^2/2
\draw[thick, engamber, dashdotted] plot[domain=-1.4:1.4, samples=80]
  ({3*\x}, {1 + \x + 0.5*\x*\x});

% Point of expansion
\fill[engblack] (0, 1) circle (2pt);
\node[engblack, anchor=south west, font=\scriptsize] at (0.1, 1.05) {$a = 0$};

% Legend
\node[engblue, anchor=west, font=\scriptsize] at (-4.4, 4.2) {$f(x) = e^{x}$};
\node[engred, anchor=west, font=\scriptsize] at (-4.4, 3.8) {linear: $1 + x$};
\node[engamber, anchor=west, font=\scriptsize] at (-4.4, 3.4) {quadratic: $1 + x + x^{2}/2$};

% Shaded error band illustration at x = 0.6
\draw[thick, enggray, dotted] (1.8, 0) -- (1.8, {exp(0.6)});
\fill[engred] (1.8, {1 + 0.6}) circle (1.4pt);
\fill[engamber] (1.8, {1 + 0.6 + 0.5*0.36}) circle (1.4pt);
\fill[engblue] (1.8, {exp(0.6)}) circle (1.4pt);
\draw[<->, engred, very thin] (1.95, {1.6}) -- (1.95, {exp(0.6)});
\node[engred, anchor=west, font=\tiny] at (2.0, 1.75) {linear gap};

\end{tikzpicture}
\caption[First and second-order Taylor approximation of $e^{x}$ at
$x = 0$]{The exponential $f(x) = e^{x}$ (blue) and its Taylor
approximations at $a = 0$. The linear approximation $L_{0}(x) = 1 + x$
(red dashed) agrees with the curve in value and slope at $x = 0$ but
diverges quadratically as $x$ moves away. The quadratic approximation
$Q_{0}(x) = 1 + x + x^{2}/2$ (amber dash-dotted) agrees in value,
slope, and curvature; its error scales cubically. At $x = 0.6$, the
linear approximation gives $1.6$ against the true $1.822$, an error
of $0.222$; the quadratic gives $1.78$, an error of $0.042$. Each
additional Taylor term reduces the error by a factor of order $h$.}
\label{fig:vol02:ch05:taylor-approximation}
\end{figure}


\Cref{fig:vol02:ch05:taylor-approximation} compares the first
two Taylor approximations of $e^{x}$ at $a = 0$ against the curve
itself. Each successive Taylor term reduces the error by a factor
of order $h$: the linear term is correct to $O(h)$, the quadratic
to $O(h^{2})$, the cubic to $O(h^{3})$. The engineer's choice of
order is set by the desired accuracy at the working range.

\subsection{Worked examples}

We linearise three functions at three working points and report
the range over which the linearisation is good to a stated
tolerance.

\engterm{Square root near unity}. Let $f(x) = \sqrt{x}$. We
linearise at $a = 1$. Then $f(1) = 1$, $f'(x) = 1/(2\sqrt{x})$,
$f'(1) = 1/2$. The linear approximation is $L_{1}(x) = 1 +
\tfrac{1}{2}(x - 1)$, equivalent to $\sqrt{1 + h} \approx 1 + h/2$
for $h = x - 1$ small. The second derivative is $f''(x) = -1/(4
x^{3/2})$, so $|f''| \leq 1/4$ for $x \geq 1$. The remainder bound
gives $|R_{1}| \leq \tfrac{1}{2} \cdot \tfrac{1}{4} \cdot h^{2} =
h^{2}/8$. For $h = 0.1$, the linear approximation gives
$\sqrt{1.1} \approx 1.05$, and the remainder bound is at most
$(0.1)^{2}/8 = 0.00125$. The actual value is $\sqrt{1.1} =
1.0488\ldots$, so the linear approximation overestimates by
$0.0012$, within the bound.

\engterm{Sine at zero}. Let $f(x) = \sin x$. We linearise at
$a = 0$. Then $f(0) = 0$, $f'(x) = \cos x$, $f'(0) = 1$. The
linear approximation is $L_{0}(x) = x$, the small-angle
approximation $\sin x \approx x$ from chapter~4 of this volume.
The second derivative is $f''(x) = -\sin x$, with $|f''| \leq 1$
for all $x$. The remainder bound gives $|R_{1}| \leq x^{2}/2$.
For $x = 0.1\,\si{\radian}$, the linear approximation gives
$\sin x \approx 0.1$ with remainder bound $0.005$; the actual
value is $0.0998\ldots$, so the linear approximation overestimates
by $0.002$, well within the bound. For $x = 0.5\,\si{\radian}$,
the bound is $0.125$, the linear approximation gives $0.5$, and
the actual value is $0.479\ldots$, so the actual error is $0.021$,
again within the bound.

\engterm{Logarithm at one}. Let $f(x) = \ln x$. We linearise at
$a = 1$. Then $f(1) = 0$, $f'(x) = 1/x$, $f'(1) = 1$. The linear
approximation is $L_{1}(x) = x - 1$, equivalent to $\ln(1 + h)
\approx h$ for $h$ small. The second derivative is $f''(x) =
-1/x^{2}$, so $|f''| \leq 1$ for $x \geq 1$. For $h = 0.05$, the
linear approximation gives $\ln(1.05) \approx 0.05$ with remainder
bound $0.00125$; the actual value is $0.0488\ldots$, error $0.001$.

\subsection{The working approximation}

The pattern is general. Whenever an engineer needs the value of a
smooth function at an input near a working point where the
function value and derivative are known, the linear approximation

\[
f(a + h) \approx f(a) + f'(a) h
\]

is correct to order $h^{2}$. The error scales with $h^{2}$ and
with the second derivative; the engineer who wants more accuracy
keeps the next term in the Taylor expansion, $\tfrac{1}{2}
f''(a) h^{2}$, and accepts an error of order $h^{3}$.

The local linearisation is the formal home of three Volume I
working tools: the small-angle approximation $\sin\theta \approx
\theta$ for $|\theta|$ small (Volume I chapter~7 used it
informally), the binomial approximation $(1 + h)^{n} \approx 1 +
n h$ for $|h|$ small (Volume I chapter~9), and the
algebraic-propagation rule for uncertainty $(u(y)/y)^{2} =
\sum_{i} p_{i}^{2} (u(x_{i})/x_{i})^{2}$ for products of powers
(Volume I chapter~4). Each is a Taylor approximation to first
order, and each acquires its remainder bound here.

\subsection{The second-order correction}

When first order is not accurate enough, the engineer keeps the
quadratic term. Taylor's theorem to second order reads

\[
f(a + h) = f(a) + f'(a) h + \tfrac{1}{2} f''(a) h^{2} +
\tfrac{1}{6} f'''(\xi) h^{3},
\]

for some $\xi$ between $a$ and $a + h$. The quadratic model
$Q_{a}(h) = f(a) + f'(a) h + \tfrac{1}{2} f''(a) h^{2}$ matches the
function in value, slope, and curvature at $a$, and its error is of
order $h^{3}$. The improvement over the linear model is dramatic
where the function curves: for $f(x) = \sqrt{x}$ at $a = 1$ with
$h = 0.1$, the linear model gives $1.05$ (error $0.0012$), while the
quadratic model adds $\tfrac{1}{2} f''(1) h^{2} = \tfrac{1}{2}
(-1/4)(0.01) = -0.00125$, giving $1.04875$ against the true
$1.04881$, an error of $6 \times 10^{-5}$, twenty times smaller.
Each retained order multiplies the working range by a factor for
the same accuracy, or the accuracy by a factor for the same range.

The engineer chooses the order by the accuracy budget. A
control-loop linearisation about an operating point keeps first
order, because the loop is designed to hold the system near the
point where the linear model is good. A trajectory propagator that
must run for many steps keeps second or higher order, because the
cubic error of each step accumulates and a smaller per-step error
buys a longer run before the accumulated error exceeds tolerance.
The bias-variance arithmetic of section~5.7 is the same arithmetic
seen from the numerical-differentiation side.

\subsection{Higher-order Taylor with an explicit remainder bound}

The pattern continues to all orders. Taylor's theorem with the
Lagrange remainder states that for $f$ with $n + 1$ continuous
derivatives on an interval containing $a$,

\[
f(a + h) = \sum_{k=0}^{n} \frac{f^{(k)}(a)}{k!} h^{k} +
\frac{f^{(n+1)}(\xi)}{(n+1)!} h^{n+1},
\]

for some $\xi$ between $a$ and $a + h$. The sum is the degree-$n$
Taylor polynomial; the last term is the remainder $R_{n}$. The
engineering value of the Lagrange form is that the remainder is a
single term of the same shape as the next series term, with the
derivative evaluated at an unknown intermediate point. If the
engineer can bound $|f^{(n+1)}|$ over the interval, the remainder
is bounded, and the approximation carries a certified error bar.

\engterm{Worked instance: $\cos x$ to fourth order, with a
guaranteed bound}. We approximate $\cos(0.4)$ using the
fourth-order Taylor polynomial of $\cos x$ at $a = 0$ and bound the
error rigorously. The derivatives of $\cos$ cycle through $\cos x$,
$-\sin x$, $-\cos x$, $\sin x$ with period four, so at $a = 0$ the
nonzero coefficients are $f(0) = 1$, $f''(0) = -1$, $f^{(4)}(0) =
1$. The degree-four Taylor polynomial is

\[
P_{4}(h) = 1 - \frac{h^{2}}{2} + \frac{h^{4}}{24}.
\]

At $h = 0.4$ this gives $P_{4}(0.4) = 1 - 0.08 + 0.001067 =
0.921067$. The remainder is $R_{4} = f^{(5)}(\xi) h^{5}/120$ with
$f^{(5)} = -\sin$. Since $|\sin\xi| \leq 1$ for every $\xi$, the
remainder is bounded by $|R_{4}| \leq h^{5}/120 = (0.4)^{5}/120 =
0.01024/120 = 8.5 \times 10^{-5}$. The true value is
$\cos(0.4) = 0.921061$, so the actual error is $6 \times 10^{-6}$,
comfortably inside the certified bound. The bound is conservative
because $|\sin\xi|$ is well below $1$ for $\xi$ near zero; a sharper
bound using $|\sin\xi| \leq \xi \leq 0.4$ on the interval gives
$|R_{4}| \leq 0.4 \cdot (0.4)^{5}/120 = 3.4 \times 10^{-5}$, still
conservative but tighter. The engineer who needs a guaranteed
digit-count chooses the polynomial order so the worst-case
remainder falls below the tolerance, then reports the polynomial
value with the bound attached.

\engterm{Choosing the order from the tolerance}. The remainder
bound inverts into an order selection. Suppose an embedded
controller must evaluate $\sin x$ for $|x| \leq 0.5$ to an
absolute accuracy of $10^{-6}$, with no library call, using only a
short polynomial. The Taylor remainder after the degree-$(2m-1)$
polynomial of $\sin$ is bounded by $|x|^{2m+1}/(2m+1)!$ because
every derivative of $\sin$ is bounded by $1$. For $|x| \leq 0.5$,
the bound is $(0.5)^{2m+1}/(2m+1)!$. At $m = 2$ (keeping terms
through $x^{3}$) the bound is $(0.5)^{5}/120 = 2.6 \times 10^{-4}$,
too coarse. At $m = 3$ (through $x^{5}$) it is $(0.5)^{7}/5040 =
1.6 \times 10^{-6}$, just above tolerance. At $m = 4$ (through
$x^{7}$) it is $(0.5)^{9}/362880 = 5.4 \times 10^{-9}$, well
inside. The controller therefore evaluates $\sin x \approx x -
x^{3}/6 + x^{5}/120 - x^{7}/5040$, a four-term polynomial whose
worst-case error over the operating range is certified below
$10^{-8}$. The remainder bound turned a vague accuracy goal into a
specific term count, and the term count into a specific number of
multiply-add operations the hardware must perform per evaluation.

\subsection{Differentials and the working notation}

The differential notation $\dd f = f'(x) \, \dd x$ packages the
linear approximation as a statement about infinitesimal changes.
Read operationally, it says: a change $\dd x$ in the input produces
a change $\dd f = f'(x) \, \dd x$ in the output, to first order. The
engineer uses the differential as a quick error-propagation device.
For $f(x) = x^{3}$, $\dd f = 3 x^{2} \, \dd x$, so a $1\,\%$ change
in $x$ near $x = 2$ produces $\dd f / f = 3 \, \dd x / x = 3\,\%$
change in $f$: the cube amplifies relative error by three. The
differential reproduces the power-law propagation rule of Volume~I
directly from the derivative, and it generalises in chapter~11 to
the total differential $\dd f = \sum_{i} (\partial f / \partial
x_{i}) \, \dd x_{i}$ for functions of several variables, which is
the engine of the GUM formula.

\subsection{Linearisation about an operating point}

A nonlinear device characterised by $y = f(x)$ is often operated in
a narrow band around a fixed working point $(x_{0}, y_{0})$. The
small-signal model replaces the device by its tangent: $\delta y
\approx f'(x_{0}) \, \delta x$, where $\delta x = x - x_{0}$ and
$\delta y = y - y_{0}$ are the deviations from the operating point.
The coefficient $f'(x_{0})$ is the small-signal gain. A diode with
exponential characteristic $I = I_{s}(e^{V/V_{T}} - 1)$ has
small-signal conductance $\dv{I}{V} = (I_{s}/V_{T}) e^{V/V_{T}} =
(I + I_{s})/V_{T} \approx I/V_{T}$ at a bias current $I$. The
reciprocal $V_{T}/I$ is the diode's small-signal resistance, the
quantity a circuit designer uses to analyse the device in its
linear regime. The entire discipline of small-signal analysis,
which Volume~IX develops for amplifiers and control systems, is the
local linearisation of this chapter applied at a bias point.

\subsection{Linearisation in action: Newton's iteration for roots}

The most heavily used application of the linear approximation in
working engineering is Newton's iteration for the root of a
nonlinear equation. Given a differentiable function $g$ and an
initial guess $x_{0}$ near a root of $g$, replace $g$ by its
linear approximation at $x_{0}$,
$L_{x_{0}}(x) = g(x_{0}) + g'(x_{0}) (x - x_{0})$, and solve the
linearised equation $L_{x_{0}}(x) = 0$ for an improved estimate
$x_{1} = x_{0} - g(x_{0})/g'(x_{0})$. Repeat from $x_{1}$ to
generate the iteration
\[
x_{n+1} = x_{n} - \frac{g(x_{n})}{g'(x_{n})}.
\]
\begin{figure}[ht]
\centering
\begin{tikzpicture}[x=1cm, y=1cm, font=\small]

% Newton's iteration on g(x) = x^3 - 2x - 5, root near x* = 2.0946.
% We picture two iterates: x0 = 3 (far) producing x1, then x1 -> x2.
% g and g' are large here, so we rescale g vertically by 1/8 for display.
% Display map: plotx = (x - 1.5)*2.4 ; ploty = g(x)/8 + 0.0 (vertical units cm).

\def\sx{2.4}      % horizontal scale
\def\ox{1.5}      % x-origin offset
\def\sy{0.18}     % vertical scale for g

% Axes
\draw[->, thick] ({(1.5-\ox)*\sx}, 0) -- ({(3.6-\ox)*\sx}, 0) node[right] {$x$};
\draw[->, thick] (0, -3.2) -- (0, 4.2) node[above] {$g(x)$};

% zero line already the x-axis. Draw curve g(x) = x^3 - 2x - 5 scaled.
\draw[thick, engblue] plot[domain=1.55:3.4, samples=80]
  ({(\x-\ox)*\sx}, {(\x*\x*\x - 2*\x - 5)*\sy});
\node[engblue, anchor=south west, font=\scriptsize] at ({(3.25-\ox)*\sx}, {(3.25*3.25*3.25 - 2*3.25 - 5)*\sy - 0.1}) {$g(x)=x^{3}-2x-5$};

% x0 = 3
\def\xa{3.0}
\pgfmathsetmacro{\ga}{(\xa*\xa*\xa - 2*\xa - 5)}
\pgfmathsetmacro{\gpa}{(3*\xa*\xa - 2)}
\pgfmathsetmacro{\xb}{\xa - \ga/\gpa}
\fill[engred] ({(\xa-\ox)*\sx}, {\ga*\sy}) circle (1.8pt);
\draw[dashed, enggray] ({(\xa-\ox)*\sx}, 0) -- ({(\xa-\ox)*\sx}, {\ga*\sy});
\node[anchor=north, font=\scriptsize] at ({(\xa-\ox)*\sx}, -0.05) {$x_{0}$};
% tangent at x0: y = g(xa) + gp(xa)(x - xa); hits zero at xb.
\draw[thick, engamber]
  ({(\xb-\ox)*\sx}, 0) -- ({(\xa-\ox)*\sx + 0.6}, {(\ga + \gpa*(\xa + 0.6/\sx - \xa))*\sy});
\fill[engblack] ({(\xb-\ox)*\sx}, 0) circle (1.6pt);
\node[anchor=north, font=\scriptsize] at ({(\xb-\ox)*\sx}, -0.05) {$x_{1}$};

% x1 -> x2
\pgfmathsetmacro{\gb}{(\xb*\xb*\xb - 2*\xb - 5)}
\pgfmathsetmacro{\gpb}{(3*\xb*\xb - 2)}
\pgfmathsetmacro{\xc}{\xb - \gb/\gpb}
\fill[engred] ({(\xb-\ox)*\sx}, {\gb*\sy}) circle (1.8pt);
\draw[dashed, enggray] ({(\xb-\ox)*\sx}, 0) -- ({(\xb-\ox)*\sx}, {\gb*\sy});
\draw[thick, engamber]
  ({(\xc-\ox)*\sx}, 0) -- ({(\xb-\ox)*\sx}, {\gb*\sy});
\fill[englime] ({(\xc-\ox)*\sx}, 0) circle (1.8pt);
\node[anchor=north east, font=\scriptsize] at ({(\xc-\ox)*\sx}, -0.05) {$x_{2}\!\approx\! x^{*}$};

\node[engamber, anchor=west, font=\scriptsize] at ({(3.05-\ox)*\sx}, 2.2) {tangents};

\end{tikzpicture}
\caption[Newton's iteration as repeated linearisation]{Newton's
iteration for the root of $g(x) = x^{3} - 2x - 5$. At each iterate
$x_{n}$ the curve is replaced by its tangent line; the tangent's
zero crossing is the next iterate $x_{n+1} = x_{n} -
g(x_{n})/g'(x_{n})$. Starting at $x_{0} = 3$ the iteration reaches
the root $x^{*} \approx 2.0946$ in three steps, the error roughly
squaring at each step. The picture is the linear approximation of
section~5.4 used in a loop: each tangent is a local model and each
zero crossing is the model's prediction of the root.}
\label{fig:vol02:ch05:newton-iteration}
\end{figure}


\Cref{fig:vol02:ch05:newton-iteration} draws the iteration as
repeated linearisation: each tangent is the local model of section
5.4, and each tangent's zero crossing is the model's prediction of
the root. Starting from $x_{0} = 3$, the iteration on $g(x) = x^{3}
- 2x - 5$ reaches the root $x^{*} \approx 2.0946$ in three steps.

The convergence theorem makes the quadratic claim precise.
Expanding $g$ to second order about the root $x^{*}$ and using
$g(x^{*}) = 0$, the iteration error $e_{n} = x_{n} - x^{*}$ obeys

\[
e_{n+1} = e_{n} - \frac{g(x_{n})}{g'(x_{n})} =
\frac{g''(\xi_{n})}{2 g'(x_{n})} \, e_{n}^{2},
\]

for some $\xi_{n}$ between $x_{n}$ and $x^{*}$. When $g'(x^{*}) \neq
0$ and $g''$ is bounded near the root, the prefactor approaches the
constant $C = |g''(x^{*})| / (2 |g'(x^{*})|)$, and the error squares
at each step: $|e_{n+1}| \approx C \, e_{n}^{2}$. Each iteration
roughly doubles the number of correct significant figures. The
quadratic law is the reason Newton's method is the default
root-finder for smooth problems: a starting error of $10^{-1}$
becomes $10^{-2}$, then $10^{-4}$, then $10^{-8}$, then $10^{-16}$
in four steps, reaching machine precision faster than any
linearly-convergent method.

The quadratic convergence is local. It holds only once the iterate
is inside a neighbourhood of the root where the error model
applies. Starting far from the root, or near a point where $g'$ is
small, the iteration can overshoot wildly, cycle, or diverge.

\begin{figure}[ht]
\centering
\begin{tikzpicture}[x=1cm, y=1cm, font=\small]

% Newton failure: iterate lands near a near-zero derivative and is
% flung far. Picture f(x) = arctan(x): the root is at 0, but starting
% outside the basin of attraction the iterate diverges. We show the
% overshoot x0 -> x1 with |x1| > |x0|.

\def\sx{1.0}
\def\sy{1.6}

% Axes
\draw[->, thick] (-5.0, 0) -- (5.2, 0) node[right] {$x$};
\draw[->, thick] (0, {-1.7*\sy}) -- (0, {1.7*\sy}) node[above] {$g(x)$};

% g(x) = arctan(x), scaled vertically
\draw[thick, engblue] plot[domain=-5:5, samples=120]
  ({\x*\sx}, {rad(atan(\x))*\sy});
\node[engblue, anchor=south west, font=\scriptsize] at (3.0, {rad(atan(3))*\sy + 0.1}) {$g(x)=\arctan x$};

% Root at origin
\fill[englime] (0, 0) circle (1.8pt);
\node[anchor=north west, font=\scriptsize] at (0.1, -0.1) {$x^{*}=0$};

% x0 = 2 (just outside the basin for arctan, which is ~1.39)
\def\xa{2.0}
\pgfmathsetmacro{\ga}{rad(atan(\xa))}
\pgfmathsetmacro{\gpa}{1/(1 + \xa*\xa)}
\pgfmathsetmacro{\xb}{\xa - \ga/\gpa}
\fill[engred] ({\xa*\sx}, {\ga*\sy}) circle (1.8pt);
\draw[dashed, enggray] ({\xa*\sx}, 0) -- ({\xa*\sx}, {\ga*\sy});
\node[anchor=south, font=\scriptsize] at ({\xa*\sx}, {\ga*\sy + 0.05}) {$x_{0}$};

% tangent at x0 to its zero crossing xb (which overshoots to the left)
\draw[thick, engamber]
  ({\xb*\sx}, 0) -- ({\xa*\sx}, {\ga*\sy});
\fill[engblack] ({\xb*\sx}, 0) circle (1.6pt);
\node[anchor=north, font=\scriptsize] at ({\xb*\sx}, -0.08) {$x_{1}$};

% show x1 on curve and next tangent overshooting further right
\pgfmathsetmacro{\gb}{rad(atan(\xb))}
\pgfmathsetmacro{\gpb}{1/(1 + \xb*\xb)}
\pgfmathsetmacro{\xc}{\xb - \gb/\gpb}
\fill[engred] ({\xb*\sx}, {\gb*\sy}) circle (1.8pt);
\draw[dashed, enggray] ({\xb*\sx}, 0) -- ({\xb*\sx}, {\gb*\sy});
\draw[thick, engamber]
  ({\xc*\sx}, 0) -- ({\xb*\sx}, {\gb*\sy});
\node[anchor=south, font=\scriptsize] at ({\xc*\sx}, 0.05) {$x_{2}$ (further out)};

\node[engred, anchor=west, font=\scriptsize] at (-4.8, {1.3*\sy}) {$|x_{2}| > |x_{1}| > |x_{0}|$: the iterates run away};

\end{tikzpicture}
\caption[Newton's iteration diverging]{Newton's iteration on
$g(x) = \arctan x$ started outside the basin of attraction. The
true root is $x^{*} = 0$, but the curve flattens away from the
origin, so $g'$ is small where $|x|$ is large. A tangent taken at a
flat point crosses zero far from the current iterate; each step
overshoots the root by more than the previous step, and the
iterates diverge. The same mechanism produces the wild behaviour of
Newton's method near a critical point of $g$ or far from a simple
root. The defensive disciplines of section~5.4 (bracket the root,
verify the iterate stays inside the bracket) guard against this
failure.}
\label{fig:vol02:ch05:newton-divergence}
\end{figure}


\Cref{fig:vol02:ch05:newton-divergence} shows the divergent mode on
$g(x) = \arctan x$, whose only root is at the origin. The function
flattens away from the origin, so $g'$ is small where $|x|$ is
large. A tangent taken at a flat point crosses zero far from the
current iterate, and each step overshoots by more than the
previous: the iterates run away. For $g(x) = \arctan x$ the basin
of attraction is roughly $|x_{0}| < 1.39$; outside it, the method
diverges. The defensive discipline is to bracket the root by a sign
change, take a coarse first guess inside the bracket, and verify
the iterate stays inside; when it leaves, fall back to a bracketing
method (bisection) for a few steps to re-enter the basin. The
accompanying scripts \texttt{code/newton\_root.py} (quadratic
convergence on $x^{3} - 2x - 5$) and \texttt{code/newton\_basins.py}
(convergence-versus-divergence sweep over starting points)
illustrate both regimes numerically.

\subsection{The secant method: Newton without the derivative}

Newton's iteration needs the analytic derivative $g'(x_{n})$ at
every step. When $g$ is a black box, a measurement, or a function
whose derivative is expensive or unavailable, the engineer replaces
the exact derivative by a finite difference across the last two
iterates. The result is the \engterm{secant method},

\[
x_{n+1} = x_{n} - g(x_{n}) \,
\frac{x_{n} - x_{n-1}}{g(x_{n}) - g(x_{n-1})}.
\]

The geometric reading is that each step draws the secant line
through the last two points $(x_{n-1}, g(x_{n-1}))$ and
$(x_{n}, g(x_{n}))$ and takes its zero crossing as the next
iterate, in place of Newton's tangent line. Two starting points
are needed rather than one, and no derivative evaluation occurs.

The convergence rate is intermediate between linear and quadratic.
A careful expansion of the error recursion shows that near a simple
root the secant error obeys $|e_{n+1}| \approx C \, |e_{n}| \,
|e_{n-1}|$, and seeking a self-consistent power law $|e_{n+1}| \sim
|e_{n}|^{p}$ gives $p = (1 + \sqrt{5})/2 = \varphi \approx 1.618$,
the golden ratio. The order is superlinear but below Newton's $2$.
The tradeoff is that each secant step costs one function evaluation,
while each Newton step costs one function evaluation plus one
derivative evaluation. When the derivative costs about as much as
the function, two secant steps (order $\varphi^{2} \approx 2.62$ in
combined progress, two evaluations) outrun one Newton step (order
$2$, two evaluations), and the secant method is the better trade
even though its per-step order is lower.

\begin{figure}[ht]
\centering
\begin{tikzpicture}[x=1cm, y=1cm, font=\small]

% Convergence comparison: log10 of the absolute error versus iteration
% index for three root-finders on a smooth root. Newton (quadratic),
% secant (order ~1.618), bisection (linear, slope ~ -log10(2) per step).
% Axes: iteration 0..8 mapped to x in [0, 9.6] (1.2 per step);
% log10 error from 0 down to -16 mapped to y: y = 0 at top (log err 0),
% y = 4.8 at bottom is log err -16, so y = 4.8 + 0.3*logerr.

% Frame
\draw[->, thick] (-0.3, 0) -- (10.2, 0) node[right] {iteration $n$};
\draw[->, thick] (0, -0.3) -- (0, 5.4)
  node[above, align=center] {$\log_{10}|e_{n}|$};

% y ticks (log error values 0, -4, -8, -12, -16)
\foreach \le in {0, -4, -8, -12, -16} {
  \pgfmathsetmacro\yy{4.8 + 0.3*\le}
  \draw (0, \yy) -- (-0.12, \yy);
  \node[anchor=east, font=\scriptsize] at (-0.14, \yy) {\le};
}
% x ticks
\foreach \n in {0,1,2,3,4,5,6,7,8} {
  \pgfmathsetmacro\xx{1.2*\n}
  \draw (\xx, 0) -- (\xx, -0.12);
  \node[anchor=north, font=\scriptsize] at (\xx, -0.12) {\n};
}

% Newton: log err doubles in magnitude each step (quadratic).
% logerr_n approx -1 * 2^n, clipped at -16.
% n: 0->-0.5, 1->-1, 2->-2, 3->-4, 4->-8, 5->-16
\draw[very thick, engblue]
  plot coordinates {
    ({1.2*0}, {4.8 + 0.3*(-0.5)})
    ({1.2*1}, {4.8 + 0.3*(-1)})
    ({1.2*2}, {4.8 + 0.3*(-2)})
    ({1.2*3}, {4.8 + 0.3*(-4)})
    ({1.2*4}, {4.8 + 0.3*(-8)})
    ({1.2*5}, {4.8 + 0.3*(-16)})
  };
\foreach \pt in {{0/-0.5},{1/-1},{2/-2},{3/-4},{4/-8},{5/-16}} {}
\node[engblue, anchor=west, font=\scriptsize] at (6.4, 4.6)
  {Newton (order 2)};

% Secant: order phi approx 1.618. logerr_n approx -0.5 * 1.618^n.
% n: 0->-0.5, 1->-0.81, 2->-1.31, 3->-2.12, 4->-3.43, 5->-5.55,
% 6->-8.98, 7->-14.5, 8->-16(clip)
\draw[very thick, engorange, dashed]
  plot coordinates {
    ({1.2*0}, {4.8 + 0.3*(-0.5)})
    ({1.2*1}, {4.8 + 0.3*(-0.81)})
    ({1.2*2}, {4.8 + 0.3*(-1.31)})
    ({1.2*3}, {4.8 + 0.3*(-2.12)})
    ({1.2*4}, {4.8 + 0.3*(-3.43)})
    ({1.2*5}, {4.8 + 0.3*(-5.55)})
    ({1.2*6}, {4.8 + 0.3*(-8.98)})
    ({1.2*7}, {4.8 + 0.3*(-14.5)})
    ({1.2*8}, {4.8 + 0.3*(-16)})
  };
\node[engorange, anchor=west, font=\scriptsize] at (6.4, 4.2)
  {secant (order $\varphi$)};

% Bisection: linear, logerr decreases by log10(2)=0.301 per step.
% n: 0->-0.5, then -0.5 - 0.301*n
\draw[very thick, engred, dotted]
  plot coordinates {
    ({1.2*0}, {4.8 + 0.3*(-0.5)})
    ({1.2*1}, {4.8 + 0.3*(-0.80)})
    ({1.2*2}, {4.8 + 0.3*(-1.10)})
    ({1.2*3}, {4.8 + 0.3*(-1.40)})
    ({1.2*4}, {4.8 + 0.3*(-1.70)})
    ({1.2*5}, {4.8 + 0.3*(-2.00)})
    ({1.2*6}, {4.8 + 0.3*(-2.31)})
    ({1.2*7}, {4.8 + 0.3*(-2.61)})
    ({1.2*8}, {4.8 + 0.3*(-2.91)})
  };
\node[engred, anchor=west, font=\scriptsize] at (6.4, 3.8)
  {bisection (linear)};

\end{tikzpicture}
\caption[Convergence rates of Newton, secant, and bisection]{Absolute
error against iteration index on a smooth root, plotted on a
logarithmic error axis. Newton's method doubles the number of correct
digits per step: the error curve bends downward, each segment twice as
steep as the last, reaching machine precision in about five steps from
a starting error of $10^{-0.5}$. The secant method, which replaces the
analytic derivative with a finite difference across the last two
iterates, converges with order $\varphi = (1 + \sqrt{5})/2 \approx
1.618$, slightly slower per step but free of a derivative evaluation.
Bisection halves the error every step, a straight line of slope
$-\log_{10} 2 \approx -0.30$, and reaches the same precision only after
roughly fifty steps. The superlinear methods win decisively once the
iterate is inside the basin where their error models hold.}
\label{fig:vol02:ch05:secant-convergence}
\end{figure}


\Cref{fig:vol02:ch05:secant-convergence} compares the three
root-finders on a logarithmic error axis. Newton bends downward
fastest; the secant trails it by a constant factor in the number of
steps; bisection, the linearly-convergent bracketing fallback,
descends as a straight line of shallow slope. The reader who
implements all three sees the qualitative ordering at once: the
superlinear methods reach machine precision in a handful of steps,
the linear method in dozens.

\engterm{Worked convergence comparison}. We solve $g(x) = x^{3} -
2x - 5 = 0$, whose root is $x^{*} \approx 2.0945514815$. Newton
from $x_{0} = 2$ produces errors $|e_{0}| = 9.5 \times 10^{-2}$,
$|e_{1}| = 5.4 \times 10^{-3}$, $|e_{2}| = 1.8 \times 10^{-5}$,
$|e_{3}| = 2.0 \times 10^{-10}$: each error is roughly the square
of the previous, scaled by the constant $C = |g''(x^{*})|/(2
|g'(x^{*})|) \approx 0.6$. The secant method from $x_{0} = 2$,
$x_{1} = 3$ produces a slower but still superlinear descent, with
the error exponent settling toward $\varphi$ once the iterates are
near the root. Bisection on the bracket $[2, 3]$ removes one bit of
uncertainty per step and needs about $50$ steps to match the
precision Newton reaches in four. The companion script
\texttt{code/newton\_root.py} tabulates the Newton run; a reader
extending it to the secant and bisection recursions reproduces the
three curves of the figure directly.

\begin{mastery}
Newton's quadratic convergence assumes a \engterm{simple} root,
one where $g(x^{*}) = 0$ but $g'(x^{*}) \neq 0$. At a root of
multiplicity $m > 1$, where $g$ and its first $m - 1$ derivatives
vanish, the convergence degrades to linear with rate $1 - 1/m$. The
iteration still converges, but it gains only a fixed fraction of a
digit per step rather than doubling the digit count. Worse, near a
multiple root the derivative $g'(x_{n})$ in the denominator is
small, so the step $g(x_{n})/g'(x_{n})$ is a ratio of two small,
noisy quantities, and floating-point round-off in each can throw
the iterate far from the root. This is the root-finding face of the
catastrophic cancellation that section~5.7 derives for numerical
differentiation: a small denominator amplifies representation
error. Two defences exist. The \engterm{modified Newton iteration}
$x_{n+1} = x_{n} - m\, g(x_{n})/g'(x_{n})$ restores quadratic
convergence when the multiplicity $m$ is known. Alternatively,
applying Newton to $u(x) = g(x)/g'(x)$, whose roots are all simple
regardless of the multiplicity of $g$, recovers quadratic
convergence without knowing $m$, at the cost of needing $g''$. The
stiffness of Newton near a multiple root is a recurring trap in
solving polynomial systems and in continuation methods that track a
solution to the edge of its existence, where a simple root and a
nearby root coalesce into a double root and $g'$ collapses.
\end{mastery}

\section{Mean value theorem and L'H\^opital}\pathtag{core}

Two classical results connect the derivative to the function's
behaviour over an interval. The first, the mean value theorem,
tells us that the average rate of change over an interval is
realised at some interior point. The second, L'H\^opital's rule,
lets us evaluate certain limits by replacing each side of a ratio
by its derivative.

Volume I Chapter 4 deferred the general partial-derivative form of
the GUM uncertainty propagation formula to this chapter on the
grounds that calculus had not yet arrived. We retroactively
legitimise it here. The
working forms of section~5.4, applied to a function of several
variables, give precisely the GUM general formula
$u^{2}(y) = \sum_{i} (\partial f / \partial x_{i})^{2} u^{2}(x_{i})$,
with the understanding that the partial derivatives are introduced
formally in chapter~11 of this volume on multivariable calculus.
The reader who returns to that Volume I mastery box now has the
calculus to read the formula as a first-order Taylor expansion in
the input deviations, with the partial derivatives playing the
role of one-variable derivatives in each input direction.

\subsection{The mean value theorem}

\begin{theorem}[Mean value theorem]
If $f$ is continuous on the closed interval $[a, b]$ and
differentiable on the open interval $(a, b)$, then there exists at
least one point $c \in (a, b)$ such that

\[
f'(c) = \frac{f(b) - f(a)}{b - a}.
\]
\end{theorem}

The right side is the average rate of change of $f$ over $[a, b]$.
The theorem says: the average rate of change over an interval is
the instantaneous rate of change at some point inside the interval.

A geometric reading: there is a point on the graph where the
tangent is parallel to the secant connecting the endpoints. A
physical reading: a car that travels from $A$ to $B$ in a stated
time has, at some moment between $A$ and $B$, an instantaneous
velocity equal to the average velocity over the trip. The theorem
does not say where that moment is; it says only that it exists.

\begin{figure}[ht]
\centering
\begin{tikzpicture}[x=1cm, y=1cm, font=\small]

% Mean value theorem: secant from (a,f(a)) to (b,f(b)) is parallel to
% the tangent at some interior c. f(x) = 0.5 + 1.6*sin(0.5*x) on [0.5, 7].

\def\f(#1){(2.0 + 1.3*sin(deg(0.45*(#1))))}

% Axes
\draw[->, thick] (-0.3, 0) -- (8.2, 0) node[right] {$x$};
\draw[->, thick] (0, 0) -- (0, 4.4) node[above] {$y = f(x)$};

% Curve
\draw[thick, engblue] plot[domain=0.4:7.6, samples=100] (\x, {\f(\x)});

% Endpoints a = 1, b = 7
\def\xa{1.0}\def\xb{7.0}
\pgfmathsetmacro{\fa}{2.0 + 1.3*sin(deg(0.45*\xa))}
\pgfmathsetmacro{\fb}{2.0 + 1.3*sin(deg(0.45*\xb))}
\fill[engred] (\xa, {\fa}) circle (1.8pt);
\fill[engred] (\xb, {\fb}) circle (1.8pt);
\node[anchor=north, font=\scriptsize] at (\xa, -0.05) {$a$};
\node[anchor=north, font=\scriptsize] at (\xb, -0.05) {$b$};

% Secant
\pgfmathsetmacro{\slope}{(\fb - \fa)/(\xb - \xa)}
\draw[thick, engamber] (\xa, {\fa}) -- (\xb, {\fb});
\node[engamber, anchor=south, font=\scriptsize] at (4.0, {\fa + \slope*3.0 + 0.1}) {secant, slope $\frac{f(b)-f(a)}{b-a}$};

% c where f'(c) = slope. f'(x) = 1.3*0.45*cos(0.45 x) = slope =>
% cos(0.45 c) = slope/(0.585). We solve numerically: pick c near 2.2.
% slope = (fb-fa)/6. For these values f'(c)=slope has a solution; mark c approx 2.15.
\def\xc{2.15}
\pgfmathsetmacro{\fc}{2.0 + 1.3*sin(deg(0.45*\xc))}
\fill[englime] (\xc, {\fc}) circle (1.8pt);
\draw[dashed, enggray] (\xc, 0) -- (\xc, {\fc});
\node[anchor=north, font=\scriptsize] at (\xc, -0.05) {$c$};
% tangent at c parallel to secant
\draw[thick, englime] (\xc-1.2, {\fc - \slope*1.2}) -- (\xc+1.2, {\fc + \slope*1.2});
\node[englime, anchor=west, font=\scriptsize] at (\xc+1.0, {\fc + \slope*1.0 + 0.2}) {tangent at $c$ (parallel)};

\end{tikzpicture}
\caption[The mean value theorem]{For $f$ continuous on $[a,b]$ and
differentiable on $(a,b)$, there is an interior point $c$ where the
instantaneous rate of change $f'(c)$ equals the average rate of
change $(f(b)-f(a))/(b-a)$ over the interval. Geometrically the
tangent at $c$ is parallel to the secant joining the endpoints. The
theorem asserts existence, not location; a function may have several
such $c$.}
\label{fig:vol02:ch05:mean-value}
\end{figure}


\Cref{fig:vol02:ch05:mean-value} draws the parallel tangent. The
existence claim has a working consequence for measurement: an
average-speed enforcement camera that records a vehicle's entry and
exit times over a known distance computes the average speed, and
the mean value theorem certifies that the vehicle reached at least
that speed at some instant inside the interval. The legal logic and
the calculus coincide.

The proof is short. Define $g(x) = f(x) - L(x)$ where $L(x)$ is
the linear function passing through $(a, f(a))$ and $(b, f(b))$.
Then $g(a) = g(b) = 0$. By Rolle's theorem (which is the special
case $f(a) = f(b)$ of the mean value theorem, proved from the
extreme value theorem), there exists $c \in (a, b)$ with $g'(c) =
0$. Computing $g'$ from $f$ and $L$ gives $f'(c) = L'(c) = (f(b)
- f(a))/(b - a)$.

The mean value theorem has consequences the engineer uses without
re-deriving them. If $f'(x) = 0$ for all $x$ in an interval, then
$f$ is constant on the interval. If $f'(x) > 0$ for all $x$ in an
interval, then $f$ is strictly increasing on the interval. If
$|f'(x)| \leq M$ for all $x$ in $[a, b]$, then $|f(b) - f(a)| \leq
M (b - a)$. Each is a direct consequence of the mean value
theorem, applied to a difference of function values at two
points.

\subsection{The bounded-derivative bound and the Lipschitz constant}

The last consequence of the mean value theorem, that a derivative
bounded by $M$ forces $|f(b) - f(a)| \leq M\,|b - a|$, is worth a name
because the engineer relies on it constantly. A function satisfying
\[
|f(x) - f(y)| \leq L\,|x - y| \quad \text{for all } x, y
\]
is called \engterm{Lipschitz continuous} with \engterm{Lipschitz
constant} $L$. The bound says the function cannot change faster than a
fixed rate: its graph is trapped between two lines of slope $\pm L$
through any of its points. The mean value theorem shows that for a
differentiable function the smallest valid $L$ is the supremum of
$|f'|$ over the interval, so the Lipschitz constant is the worst-case
slope. A Lipschitz function need not be differentiable; $f(x) = |x|$ has
no derivative at the origin yet is Lipschitz with $L = 1$, because its
slope never exceeds one in magnitude where it exists. The Lipschitz
constant is the working replacement for the derivative bound when the
function has kinks, because it survives the corners where the derivative
does not exist.

The constant earns its keep in three places the engineer meets early.
First, it certifies that a numerical scheme is stable: the error
amplification of a one-step integrator over a time step $\Delta t$ is
controlled by $L\,\Delta t$ where $L$ is the Lipschitz constant of the
right-hand side, and the step is taken small enough that $L\,\Delta t$
stays below a stability threshold. Chapter~7 makes this precise for
ordinary differential equations; the Lipschitz constant is the quantity
that the existence-and-uniqueness theorem for those equations requires,
because two solution curves of a Lipschitz field cannot cross or merge.
Second, it bounds the sensitivity of an output to an input over a finite
range, not merely at a point: a controller whose plant has Lipschitz
constant $L$ cannot see its output move by more than $L$ times the input
move, which sets a worst-case actuator authority. Third, it governs the
convergence of fixed-point iteration. A map $g$ with Lipschitz constant
$L < 1$ on an interval it maps into itself is a \engterm{contraction},
and the iteration $x_{n+1} = g(x_{n})$ converges to a unique fixed point
$x^{*} = g(x^{*})$ at the geometric rate $|x_{n} - x^{*}| \leq L^{n}
|x_{0} - x^{*}|$. The contraction condition is checked by bounding
$|g'|$: if $|g'(x)| \leq L < 1$ throughout, the mean value theorem
delivers the contraction directly.

A worked instance. The fixed-point form of a root-finding problem
rewrites $h(x) = 0$ as $x = g(x)$ and iterates. For the equation
$x = \cos x$ (the root of $h(x) = x - \cos x$), the iteration map is
$g(x) = \cos x$ with $g'(x) = -\sin x$, so $|g'(x)| = |\sin x| \leq
\sin(1) \approx 0.84 < 1$ on the interval around the root near
$x^{*} \approx 0.739$. The map is a contraction with $L \approx 0.84$,
so the naive iteration $x_{n+1} = \cos x_{n}$ converges from any
starting point in the interval, gaining a fixed fraction of a digit per
step, $\log_{10}(1/L) \approx 0.075$ digits per iteration. The slow
geometric rate is the price of not using the derivative in the step the
way Newton does. A reader who compares this contraction iteration to
Newton on the same root sees the contrast between linear convergence
governed by a Lipschitz constant and the quadratic convergence governed
by the second derivative, the two regimes section~5.4 set side by side.

\subsection{Convexity detection on a real cost function}

The second-derivative convexity reading of section~5.3 becomes a working
test when the engineer must certify that a cost function has a single
global minimum before trusting a solver to find it. A twice-differentiable
cost is convex on an interval exactly when $f''(x) \geq 0$ throughout, and
on a convex cost the first-order condition $f'(x) = 0$ locates the global
minimum with no search over competing local minima. The test is
worth running before optimisation, because a solver launched on a
non-convex cost can converge to a local minimum that is not the answer.

A worked instance from regression. The least-squares cost for fitting a
single parameter $\theta$ to data, $J(\theta) = \sum_{i} (y_{i} -
\theta x_{i})^{2}$, has $J''(\theta) = 2\sum_{i} x_{i}^{2} \geq 0$, so it
is convex (strictly, whenever any $x_{i} \neq 0$), and the normal
equation $J'(\theta) = 0$ returns the global optimum directly. This is
the structural reason least-squares fitting is reliable: its cost is
convex by construction. Contrast a cost with a sigmoid saturation, such
as $C(\theta) = (y - \sigma(\theta))^{2}$ with $\sigma$ the logistic
function. Its second derivative changes sign, so the cost is non-convex,
and gradient descent on it can stall at a flat shoulder far from the
optimum. The engineer who checks the sign of $C''$ before optimising
learns which regime they are in and chooses the method accordingly:
a single Newton step for a convex cost, a multi-start or globalised
search for a non-convex one.

Convexity also fixes the step size of gradient descent through the
Lipschitz constant of the derivative. If the gradient $f'$ is itself
Lipschitz with constant $\beta$ (equivalently $f'' \leq \beta$ where the
second derivative exists), then gradient descent $x_{n+1} = x_{n} -
\eta\,f'(x_{n})$ converges for any fixed step $\eta < 2/\beta$, and the
fastest safe choice is $\eta = 1/\beta$
\cite{text:boyd-vandenberghe2004}. The bound is read off the curvature:
a sharply curving cost (large $\beta$) demands a small step, a gently
curving cost tolerates a large one. The convexity of the cost and the
Lipschitz constant of its gradient are the two numbers a reader needs to
set up a first-order optimiser, and both are second-derivative readings
of the chapter.

\subsection{L'H\^opital's rule}

A common practical question: what is the limit of a ratio
$f(x)/g(x)$ as $x \to a$, when $f(a) = g(a) = 0$? The ratio is
the indeterminate form $0/0$ at $x = a$, and the limit is not
recovered by direct substitution.

\begin{theorem}[L'H\^opital's rule]
Let $f$ and $g$ be differentiable in a neighbourhood of $a$, with
$g'(x) \neq 0$ near $a$ (except possibly at $a$). Suppose
$\lim_{x \to a} f(x) = \lim_{x \to a} g(x) = 0$ (or both limits
are $\pm \infty$). If $\lim_{x \to a} f'(x)/g'(x)$ exists (or is
$\pm \infty$), then $\lim_{x \to a} f(x)/g(x)$ exists and equals
the limit of the ratio of derivatives.
\end{theorem}

The rule has the geometric reading: near $a$, both $f$ and $g$
are approximately linear, so $f(x) \approx f'(a)(x - a)$ and
$g(x) \approx g'(a)(x - a)$. The ratio is approximately
$f'(a)/g'(a)$, exact in the limit if that quotient is finite.

A worked use. $\lim_{x \to 0} \sin x / x$. Direct substitution
gives $0/0$. By L'H\^opital, the limit equals $\lim_{x \to 0}
\cos x / 1 = 1$. The reader who wants to confirm this from first
principles can use the Taylor expansion $\sin x = x - x^{3}/6 +
\cdots$, so $\sin x / x = 1 - x^{2}/6 + \cdots \to 1$.

Another. $\lim_{x \to 0} (e^{x} - 1)/x = \lim_{x \to 0} e^{x}/1
= 1$. And $\lim_{x \to 0} (1 - \cos x)/x^{2} = \lim_{x \to 0}
\sin x / (2 x) = \tfrac{1}{2}$, where the second equality uses
L'H\^opital again because the intermediate ratio is still $0/0$
at $x = 0$.

The rule extends, with a certain amount of care, to other
indeterminate forms ($\infty/\infty$, $0 \cdot \infty$,
$\infty - \infty$, $0^{0}$, $1^{\infty}$, $\infty^{0}$). Each
non-quotient form is converted to a $0/0$ or $\infty/\infty$ ratio
by algebraic manipulation, often using $\ln$ to handle exponential
forms.

A worked exponential form. The limit $\lim_{x \to 0^{+}} x^{x}$ is
the indeterminate $0^{0}$. Take the logarithm: $\ln(x^{x}) = x \ln
x$, an $0 \cdot (-\infty)$ form, which rewrites as $\ln x / (1/x)$,
an $\infty/\infty$ ratio. L'H\^opital gives $\lim (1/x)/(-1/x^{2})
= \lim (-x) = 0$, so $\ln(x^{x}) \to 0$ and $x^{x} \to e^{0} = 1$.
The compound-interest limit $\lim_{n \to \infty} (1 + r/n)^{n}$ is
the indeterminate $1^{\infty}$; its logarithm is $n \ln(1 + r/n)$,
which by the same route tends to $r$, so the limit is $e^{r}$, the
continuous-compounding formula.

\subsection{Sensitivity propagation, several inputs}

The first-order linearisation of section~5.4, applied to a function
of several inputs, is the working form of error propagation. If a
derived quantity $y = f(x_{1}, \ldots, x_{m})$ depends on inputs
each known with a small uncertainty, then the deviation in $y$ is
the sum of the deviations in each input weighted by its
sensitivity,

\[
\delta y \approx \sum_{i=1}^{m} \frac{\partial f}{\partial x_{i}}
\, \delta x_{i},
\]

where the partial derivative $\partial f / \partial x_{i}$ is the
ordinary one-variable derivative taken with all other inputs held
fixed (chapter~11 introduces the partial derivative formally). For
independent inputs whose uncertainties combine in quadrature, the
deviations add as variances rather than linearly, giving the GUM
general formula

\[
u^{2}(y) = \sum_{i=1}^{m} \left(\frac{\partial f}{\partial x_{i}}
\right)^{2} u^{2}(x_{i}),
\]

the promise Volume~I deferred to this chapter. The partial
derivatives are the sensitivity coefficients: they convert each
input's uncertainty, in its own units, into a contribution to the
output's uncertainty, in the output's units.

A worked instance. The electrical power dissipated in a resistor is
$P = V^{2}/R$. The sensitivities are $\partial P / \partial V =
2V/R$ and $\partial P / \partial R = -V^{2}/R^{2}$. Dividing the
GUM formula through by $P^{2} = V^{4}/R^{2}$ collapses it to the
relative form $(u(P)/P)^{2} = 4 (u(V)/V)^{2} + (u(R)/R)^{2}$: the
voltage enters with its relative uncertainty doubled (the square in
$V^{2}$), the resistance with its relative uncertainty unamplified.
A reader measuring power across a $1\,\%$-uncertain voltage and a
$0.5\,\%$-uncertain resistance gets $(u(P)/P)^{2} = 4(0.01)^{2} +
(0.005)^{2} = 4.25 \times 10^{-4}$, so $u(P)/P \approx 2.1\,\%$.
The voltage dominates the budget, and the engineer who wants a
tighter power figure improves the voltage measurement first.

\subsection{The derivative becomes a matrix: a Jacobian preview}

The single-input single-output derivative of this chapter has a direct
generalisation that chapter~11 develops in full and that the engineer
meets the moment a calculation has more than one output. When a map
sends a vector of inputs $\mathbf{x} = (x_{1}, \ldots, x_{n})$ to a
vector of outputs $\mathbf{f} = (f_{1}, \ldots, f_{m})$, the role of the
scalar $f'(x)$ is played by the \engterm{Jacobian matrix}, the
rectangular array of all first partial derivatives,
\[
J(\mathbf{x}) = \begin{pmatrix}
\dfrac{\partial f_{1}}{\partial x_{1}} & \cdots &
\dfrac{\partial f_{1}}{\partial x_{n}} \\[1.2ex]
\vdots & & \vdots \\[0.6ex]
\dfrac{\partial f_{m}}{\partial x_{1}} & \cdots &
\dfrac{\partial f_{m}}{\partial x_{n}}
\end{pmatrix},
\]
an $m \times n$ matrix whose entry in row $i$, column $j$ is the
sensitivity of the $i$-th output to the $j$-th input. The local
linearisation $f(a + h) \approx f(a) + f'(a) h$ of section~5.4 carries
over verbatim with the scalar derivative promoted to the matrix and the
scalar product promoted to a matrix-vector product,
\[
\mathbf{f}(\mathbf{x} + \mathbf{h}) \approx \mathbf{f}(\mathbf{x}) +
J(\mathbf{x})\,\mathbf{h},
\]
the best linear approximation of a vector map near a point. The structure
is the same value-plus-linear-correction the chapter has carried since
section~5.1; the Jacobian is the linear correction written as a matrix.

A worked instance reads the Jacobian of the polar-to-Cartesian map
$x = r \cos\theta$, $y = r \sin\theta$, the change of variables that
recurs in every transport calculation in cylindrical geometry. The
inputs are $(r, \theta)$ and the outputs are $(x, y)$, so the Jacobian is
the $2 \times 2$ matrix
\[
J = \begin{pmatrix}
\partial x/\partial r & \partial x/\partial\theta \\
\partial y/\partial r & \partial y/\partial\theta
\end{pmatrix}
= \begin{pmatrix}
\cos\theta & -r \sin\theta \\
\sin\theta & r \cos\theta
\end{pmatrix}.
\]
The determinant of this matrix is $r \cos^{2}\theta + r \sin^{2}\theta =
r$, the factor that appears as the $r\,\dd r\,\dd\theta$ area element in
polar integration. The determinant of the Jacobian measures how the map
stretches or shrinks volume locally, the several-variable analogue of
$|f'(x)|$ measuring how the one-dimensional map stretches length; where
the determinant vanishes, the map collapses a region to lower dimension
and is locally non-invertible, the several-variable face of the
$f'(x) = 0$ condition that defeats the inverse-function formula of
section~5.2. Chapter~11 develops the partial derivative, the gradient,
and the Jacobian as first-class objects; the present preview names the
matrix and reads its determinant so the reader meets the object as a
generalisation of the scalar derivative rather than as a new invention.

\subsection{The multivariable chain rule, previewed}

The chain rule of section~5.2 has a several-variable form that the
engineer uses whenever a quantity depends on time through several
intermediate variables. If $z = f(x, y)$ and both $x$ and $y$ depend on
a single parameter $t$, the total rate of change of $z$ along the path is
the sum of the two channels through which $t$ reaches $z$,
\[
\dv{z}{t} = \frac{\partial f}{\partial x}\dv{x}{t} +
\frac{\partial f}{\partial y}\dv{y}{t},
\]
each partial derivative weighting the rate at which its input moves. The
single-variable chain rule $\dv{f}{t} = \dv{f}{u}\dv{u}{t}$ is the case
of one intermediate variable; the multivariable form adds one term per
intermediate variable, and the structure is exactly the matrix-vector
product $\dv{z}{t} = \nabla f \cdot \dv{\mathbf{x}}{t}$ with the gradient
$\nabla f = (\partial f/\partial x, \partial f/\partial y)$ playing the
role of the row of sensitivities. A worked instance: the temperature
felt by a sensor moving through a spatial temperature field $T(x, y)$
along a trajectory $(x(t), y(t))$ changes at the rate $\dv{T}{t} =
T_{x}\dot{x} + T_{y}\dot{y}$, the dot product of the temperature gradient
with the velocity, which is why a probe crossing a steep thermal gradient
quickly reads a fast temperature change even when the field is static.
Chapter~11 derives this rule; the present preview connects it to the
one-variable chain rule so the reader sees the addition of one term per
path as the natural generalisation.

\subsection{Differentiating an integral: the Leibniz rule}

A quantity is often defined as an integral whose limits or integrand
depend on a parameter, and the engineer needs the derivative of that
quantity with respect to the parameter. The \engterm{Leibniz integral
rule} provides it. For an integral with parameter-dependent limits and a
parameter-dependent integrand,
\[
\dv{}{t}\int_{a(t)}^{b(t)} g(x, t)\,\dd x =
g(b(t), t)\,b'(t) - g(a(t), t)\,a'(t) +
\int_{a(t)}^{b(t)}\frac{\partial g}{\partial t}(x, t)\,\dd x.
\]
The first two terms are the chain rule applied to the moving limits: as
the upper limit advances at rate $b'(t)$ it sweeps area at the local
integrand value $g(b(t), t)$, and the lower limit subtracts by the same
logic. The third term is the integral of the integrand's own rate of
change, the contribution from the integrand shifting while the limits
hold still. The rule is the product rule and chain rule of section~5.2
applied to the two ways a parametrised area can change.

A worked instance the engineer meets in heat transfer. The total thermal
energy stored in a growing layer of thickness $L(t)$ with temperature
profile $T(x, t)$ is $E(t) = \int_{0}^{L(t)} \rho c\,T(x, t)\,\dd x$,
with $\rho$ the density and $c$ the specific heat. The Leibniz rule gives
the rate of energy accumulation
\[
\dv{E}{t} = \rho c\,T(L(t), t)\,L'(t) +
\int_{0}^{L(t)} \rho c\,\frac{\partial T}{\partial t}\,\dd x,
\]
the first term the energy carried in by the advancing boundary (a
moving-boundary or Stefan-problem term, the rate at which the growing
layer engulfs fresh material at the boundary temperature) and the second
term the energy added by the existing material warming in place. A
reader who differentiates an integral whose limits move and keeps only
the integrand term, forgetting the boundary terms, undercounts the
accumulation by the moving-boundary flux, a recurring error in
melting-front, charging-front, and reaction-front calculations. The
Leibniz rule is the bookkeeping that keeps both contributions.

\begin{mastery}
The mean value theorem extends in several directions. The Cauchy
mean value theorem replaces the linear function $L$ with a more
flexible parametric statement. The integral mean value theorem
expresses the integral of a continuous function over an interval
as the function's value at some interior point times the interval
length. Both extensions appear in chapter~6 (integration) and in
the foundational results of numerical analysis. The Cauchy form
is also the engine of L'H\^opital's rule's proof; we omit the
formal derivation and refer the reader to a standard analysis
text.
\end{mastery}

\section{Optimisation of one variable}\pathtag{core}

The engineer's first formal use of the derivative is to find the
maximum or minimum of a function. The chapter has been preparing
this subject from its opening: the derivative locates the rate, the
second derivative locates the curvature, and the optimisation
archetype is announced in section~5.1. We now formalise the
procedure.

\begin{archetype}[Optimisation in one variable]
Given a function $f$ defined on a domain (an interval, the whole
real line, a finite set of allowed points), the optimisation
problem is to find the input that maximises or minimises $f$. In
one continuous variable, the candidate inputs are the points where
$f'(x) = 0$ (interior critical points) and the endpoints of the
domain. The second derivative test selects among the interior
critical points. The endpoint test compares values at the
boundaries.
\end{archetype}

\begin{figure}[ht]
\centering
\begin{tikzpicture}[x=1cm, y=1cm, font=\small]

% Curve illustrating local min, local max, inflection, endpoint optimum
% Use f(x) = 0.05 x (x-3)(x-7)(x-11) + 1.5, shifted/scaled into [0, 12] x [0, 4.5]
% Actually use a hand-tuned cubic-like polynomial showing two interior crit pts + endpoints

% Axes
\draw[->, thick] (-0.3, 0) -- (12.5, 0) node[right] {$x$};
\draw[->, thick] (0, -0.2) -- (0, 5.0) node[above] {$f(x)$};

% Domain markers
\draw[engblack, dashed, thin] (1, 0) -- (1, 5.0);
\draw[engblack, dashed, thin] (11, 0) -- (11, 5.0);
\node[engblack, anchor=north, font=\scriptsize] at (1, 0) {$x = a$};
\node[engblack, anchor=north, font=\scriptsize] at (11, 0) {$x = b$};

% A smooth curve: f(x) = 1.5 + sin((x - 1)*pi/5) + 0.4*sin((x-1)*pi/2.5)
\draw[very thick, engblue]
  plot[domain=1:11, samples=120]
    (\x, {2.2 + 1.0*sin((\x - 1)*36) + 0.4*sin((\x - 1)*72)});

% Compute features by inspection at sampled points
% At x = 3.4 (local max), x = 6 (inflection w/ horizontal tangent? not really, just inflection), x = 8.6 (local min)
% Just place markers approximately
\fill[engred] (3.4, {2.2 + 1.0*sin(2.4*36) + 0.4*sin(2.4*72)}) circle (2.5pt);
\node[engred, anchor=south, font=\scriptsize, align=center]
  at (3.4, {2.2 + 1.0*sin(2.4*36) + 0.4*sin(2.4*72) + 0.05})
  {local max\\ $f'=0,\,f''<0$};

\fill[engamber] (8.6, {2.2 + 1.0*sin(7.6*36) + 0.4*sin(7.6*72)}) circle (2.5pt);
\node[engamber, anchor=north, font=\scriptsize, align=center]
  at (8.6, {2.2 + 1.0*sin(7.6*36) + 0.4*sin(7.6*72) - 0.1})
  {local min\\ $f'=0,\,f''>0$};

% Endpoint markers
\fill[engblack] (1, {2.2}) circle (2pt);
\node[engblack, anchor=south east, font=\scriptsize, align=right]
  at (0.9, {2.2 + 0.05}) {left\\ endpoint};

\fill[engblack] (11, {2.2 + 1.0*sin(10*36) + 0.4*sin(10*72)}) circle (2pt);
\node[engblack, anchor=south west, font=\scriptsize, align=left]
  at (11.1, {2.2 + 1.0*sin(10*36) + 0.4*sin(10*72) + 0.05}) {right\\ endpoint};

% Horizontal tangents at the critical points
\draw[engred, thick] (2.6, {2.2 + 1.0*sin(2.4*36) + 0.4*sin(2.4*72)})
                  -- (4.2, {2.2 + 1.0*sin(2.4*36) + 0.4*sin(2.4*72)});
\draw[engamber, thick] (7.8, {2.2 + 1.0*sin(7.6*36) + 0.4*sin(7.6*72)})
                   -- (9.4, {2.2 + 1.0*sin(7.6*36) + 0.4*sin(7.6*72)});

% Caption legend
\node[anchor=north, font=\scriptsize, align=left]
  at (6, -0.7)
  {On $[a, b]$, the extrema lie at interior critical points (Fermat: $f'(x) = 0$, classify by $f''$) or at the endpoints.\\
   The exhaustive procedure: list all candidates, evaluate $f$ at each, pick the largest and smallest.};

\end{tikzpicture}
\caption[Critical points, second-derivative classification, and
endpoint candidates]{One-variable optimisation on a closed interval
$[a, b]$. Interior local extrema satisfy Fermat's necessary condition
$f'(x) = 0$, with the second-derivative test sorting maxima ($f'' < 0$)
from minima ($f'' > 0$). Endpoints are candidates regardless of the
derivative. The global maximum and minimum lie at one of the
candidates; the procedure evaluates $f$ at each and picks the
extremes. Points where $f'(x) = 0$ but $f''(x) = 0$ require higher
derivatives to classify, and points where $f$ is not differentiable
(kinks, cusps) must be added to the candidate list directly.}
\label{fig:vol02:ch05:optimisation-extrema}
\end{figure}


\Cref{fig:vol02:ch05:optimisation-extrema} draws the catalogue: a
single closed-interval optimisation contains every kind of
candidate the chapter discusses. The interior critical points
(Fermat's necessary condition) are classified by the second
derivative; the endpoints enter without a derivative test. The
multivariable analogue of the figure, with critical points
replaced by gradient-zero conditions and the curvature
information held in the Hessian matrix, appears in chapter~15.

\subsection{First-order necessary condition}

\begin{theorem}[Fermat's stationary-point theorem]
If $f$ has a local maximum or minimum at an interior point $c$ of
its domain, and $f$ is differentiable at $c$, then $f'(c) = 0$.
\end{theorem}

The proof is a sign-of-difference-quotient argument. If $f$ has a
local maximum at $c$, then for $h$ small, $f(c + h) \leq f(c)$,
so $(f(c + h) - f(c))/h \leq 0$ for $h > 0$ and
$(f(c + h) - f(c))/h \geq 0$ for $h < 0$. Taking the limit from
each side gives $f'(c) \leq 0$ and $f'(c) \geq 0$, so $f'(c) =
0$. The argument for a minimum is symmetric.

A point where $f'(x) = 0$ is called a \engterm{stationary point}
or \engterm{critical point} of $f$. The theorem says: every
interior local maximum or minimum of a differentiable function is
a critical point. The converse is false. The function $f(x) =
x^{3}$ has $f'(0) = 0$ but is neither a maximum nor a minimum at
$x = 0$; it is an inflection point with horizontal tangent.

\subsection{Second-order sufficient condition}

\begin{theorem}[Second-derivative test]
Let $f$ be twice continuously differentiable on an interval
containing $c$, with $f'(c) = 0$. If $f''(c) > 0$, then $c$ is a
local minimum. If $f''(c) < 0$, then $c$ is a local maximum. If
$f''(c) = 0$, the test is inconclusive.
\end{theorem}

The proof reads the second derivative as concavity. If $f''(c) >
0$ and $f''$ is continuous, then $f'' > 0$ in a neighbourhood of
$c$, so $f'$ is increasing through zero at $c$. The function $f$
is therefore decreasing on a left neighbourhood of $c$ (where
$f' < 0$) and increasing on a right neighbourhood (where $f' >
0$); the value $f(c)$ is a local minimum.

When $f''(c) = 0$, the test fails and the engineer needs further
information. Higher-order terms determine the answer: if $f'''(c)
\neq 0$, the point is an inflection (neither a maximum nor a
minimum); if $f'''(c) = 0$ and $f''''(c) > 0$, the point is a
local minimum; and so on. The fourth-order case is rare in
practical engineering optimisation; the second-derivative test
covers most working problems.

\subsection{Worked examples}

\engterm{Aluminium can}. A cylindrical can of fixed volume $V$ has
radius $r$ and height $h$ with $\pi r^{2} h = V$. The surface
area, including top and bottom, is $A(r) = 2\pi r^{2} + 2\pi r h
= 2\pi r^{2} + 2 V / r$. The first derivative is $A'(r) = 4\pi r
- 2 V / r^{2}$, vanishing at $r^{*} = (V / (2\pi))^{1/3}$. The
second derivative is $A''(r) = 4\pi + 4 V / r^{3} > 0$, so the
critical point is a local minimum. At the optimum, $h^{*} = V/
(\pi (r^{*})^{2}) = 2 r^{*}$: the optimal can is twice as tall as
it is wide. (Real soup cans are taller than this, because labels,
shipping, and grip preferences enter the cost function alongside
material.)

\engterm{Power transfer to a resistor}. A battery with electromotive
force $\mathcal{E}$ and internal resistance $R_{i}$ delivers power
$P(R) = \mathcal{E}^{2} R / (R + R_{i})^{2}$ to an external load
$R$. We maximise $P$ over $R > 0$. Using the quotient rule,
$P'(R) = \mathcal{E}^{2} ((R + R_{i})^{2} - R \cdot 2(R + R_{i}))
/ (R + R_{i})^{4} = \mathcal{E}^{2} (R_{i} - R) / (R + R_{i})^{3}$.
This vanishes at $R^{*} = R_{i}$. The second derivative is
negative there, so the critical point is a maximum. Maximum power
transfer occurs when the load resistance equals the internal
resistance. The maximum power is $P^{*} = \mathcal{E}^{2} / (4
R_{i})$.

\engterm{Drag minimisation for a fixed cruising speed}. The drag
on a body in air at low speed has two components: a viscous drag
proportional to velocity, $F_{\text{visc}} = \alpha v$, and an
induced drag (for a lifting body) inversely proportional to the
square of velocity, $F_{\text{ind}} = \beta / v^{2}$. The total
drag is $F(v) = \alpha v + \beta / v^{2}$. Setting $F'(v) =
\alpha - 2 \beta / v^{3} = 0$ gives $v^{*} = (2 \beta /
\alpha)^{1/3}$. The second derivative $F''(v) = 6 \beta / v^{4}$
is positive, so $v^{*}$ minimises drag. The cruise speed of an
aircraft, current as of 2026 in commercial design practice, is
chosen near this optimum, modified by engine efficiency and
flight-time considerations.

\engterm{Economic order quantity}. A warehouse orders a part in
batches of size $Q$. Ordering costs $K$ per order, and holding the
inventory costs $h$ per unit per year; the annual demand is $D$.
The annual total cost is $C(Q) = K D / Q + h Q / 2$: the first term
falls with larger batches (fewer orders), the second rises (more
held stock). Setting $C'(Q) = -K D / Q^{2} + h/2 = 0$ gives the
economic order quantity $Q^{*} = \sqrt{2 K D / h}$. The second
derivative $C''(Q) = 2 K D / Q^{3} > 0$ confirms a minimum. The
square-root form is the signature of a two-term cost with one term
rising and one falling; the same algebra recurs in batch-process
scheduling, server-provisioning, and any setting trading a fixed
per-event cost against a per-unit holding cost. At the optimum the
two cost terms are equal, $K D / Q^{*} = h Q^{*}/2$, a balance the
engineer can check by inspection.

\engterm{Snell's law from least time}. Light travels from a point
in one medium to a point in a second medium, crossing a flat
interface. The speed is $v_{1}$ in the first medium and $v_{2}$ in
the second. Fermat's principle states that the path taken minimises
the travel time. Parametrise the crossing point on the interface by
its horizontal coordinate $x$; the travel time is $T(x) =
\sqrt{a^{2} + x^{2}}/v_{1} + \sqrt{b^{2} + (d - x)^{2}}/v_{2}$,
where $a$ and $b$ are the perpendicular distances of the two points
from the interface and $d$ is their horizontal separation.
Differentiating and setting $T'(x) = 0$ gives

\[
\frac{x}{v_{1}\sqrt{a^{2} + x^{2}}} = \frac{d - x}{v_{2}
\sqrt{b^{2} + (d - x)^{2}}},
\]

which is $\sin\theta_{1}/v_{1} = \sin\theta_{2}/v_{2}$, Snell's law
of refraction, with $\theta_{1}$ and $\theta_{2}$ the angles from
the interface normal. A physical law falls out of a one-variable
optimisation; the derivative locates the path nature actually
takes. The same least-time argument, applied to a lifeguard
running on sand and swimming in water, gives the lifeguard's
optimal entry point, the standard pedagogical mirror of the optical
problem.

\begin{estimation}
\textbf{Question.} A delivery company runs warehouses that each
order a fast-moving part with annual demand $D = 50{,}000$ units.
Each order costs about $K = \$80$ to place and process, and holding
one unit in stock for a year costs about $h = \$4$. Estimate the
economic order quantity and the number of orders placed per year,
then estimate how sensitive the total cost is to ordering, say,
twice the optimal batch.

\textbf{Estimate.} The economic order quantity is $Q^{*} =
\sqrt{2 K D / h} = \sqrt{2 \cdot 80 \cdot 50{,}000 / 4} =
\sqrt{2 \times 10^{6}} \approx 1{,}400$ units, so the company
places about $D/Q^{*} \approx 36$ orders per year, roughly three
per month. The minimum annual cost is $C(Q^{*}) = \sqrt{2 K D h} =
\sqrt{2 \cdot 80 \cdot 50{,}000 \cdot 4} \approx \$5{,}700$.
Ordering twice the optimum, $Q = 2 Q^{*}$, gives $C = K D /(2 Q^{*})
+ h (2 Q^{*})/2 = \tfrac{1}{2} C_{\text{order}}^{*} + 2
C_{\text{hold}}^{*}$. Since the two cost halves are equal at the
optimum, this is $\tfrac{1}{2}(C^{*}/2) + 2(C^{*}/2) = \tfrac{5}{4}
C^{*}$, a $25\,\%$ cost penalty for a factor-of-two batch error.
The cost minimum is flat: near the optimum the total cost is
quadratic in the relative batch error, so modest mis-sizing costs
little. This flatness is why simple round-number batch policies
survive in practice, and it is the optimisation analogue of the
small remainder near a smooth minimum.
\end{estimation}

\subsection{Constrained optimisation in one variable}

When the domain is a closed interval $[a, b]$, the optimum may
occur at an interior critical point, at $x = a$, or at $x = b$.
The procedure is to compute $f$ at every interior critical point
(found via $f'(x) = 0$) and at the two endpoints, and pick the
largest value (for a maximum) or the smallest value (for a
minimum). The procedure is exhaustive: every candidate is on the
list.

\engterm{A worked constrained problem}. Maximise $f(x) = x e^{-x}$
on $[0, 3]$. The first derivative is $f'(x) = e^{-x} (1 - x)$,
vanishing at $x = 1$ (the only interior critical point in the
domain). Compute $f(0) = 0$, $f(1) = 1/e \approx 0.368$, $f(3) =
3 e^{-3} \approx 0.149$. The maximum on $[0, 3]$ is at $x = 1$
with value $1/e$, and the minimum is at $x = 0$ with value $0$.

The same procedure applies to chapter~15's multivariable
constrained optimisation, with critical points replaced by
gradient-zero conditions and Lagrange multipliers handling
equality constraints. We introduce the multivariable case
formally there.

\engterm{A worked design optimisation with an active constraint}.
A rectangular open-topped channel of cross-sectional area $A =
2\,\si{\meter\squared}$ is to carry water with the least wetted
perimeter, because the perimeter sets the frictional drag on the
flow. The channel has width $w$ and depth $d$ with $w d = A$, and
the wetted perimeter (bottom plus two sides, open top) is $P = w +
2 d$. Eliminating $w = A/d$ gives $P(d) = A/d + 2 d$, a function of
the single variable $d > 0$. The first derivative is $P'(d) = -A/
d^{2} + 2$, vanishing at $d^{*} = \sqrt{A/2} = 1\,\si{\meter}$. The
second derivative $P''(d) = 2 A/d^{3} > 0$ confirms a minimum. At
the optimum $w^{*} = A/d^{*} = 2\,\si{\meter}$, so the best
open channel is twice as wide as it is deep, $w^{*} = 2 d^{*}$, the
half-square cross-section that minimises wetted perimeter for a
rectangular open channel. The minimum perimeter is $P^{*} =
2 + 2 = 4\,\si{\meter}$.

Now impose a fabrication constraint: the available sheet metal
limits the depth to $d \leq 0.6\,\si{\meter}$. The unconstrained
optimum $d^{*} = 1\,\si{\meter}$ lies outside the feasible interval
$[d_{\min}, 0.6]$, so the constraint is \engterm{active}: the
optimum sits at the boundary $d = 0.6\,\si{\meter}$, not at the
interior critical point. The procedure catches this automatically.
On $[0.3, 0.6]$ (taking a lower fabrication limit $d_{\min} =
0.3\,\si{\meter}$), $P'(d) = -2/d^{2} + 2 < 0$ throughout, because
$d < d^{*}$ makes the $-A/d^{2}$ term dominate; the perimeter is
strictly decreasing on the feasible interval, so the minimum is at
the right endpoint $d = 0.6\,\si{\meter}$, where $w = A/d = 3.33\,
\si{\meter}$ and $P = 3.33 + 1.2 = 4.53\,\si{\meter}$. The
constrained optimum is worse than the unconstrained one ($4.53$
against $4.00\,\si{\meter}$), and the difference $0.53\,\si{\meter}$
is the price the fabrication limit imposes on the hydraulic design.
A reader who solved $P'(d) = 0$ and reported $d = 1\,\si{\meter}$
without checking feasibility would have specified a channel the
shop cannot build. The discipline is to find the interior critical
points, discard those outside the feasible set, and compare the
surviving candidates against the endpoints.

\begin{mastery}
The second-derivative test is a local statement. A function with
$f''(c) > 0$ at a critical point has a local minimum there, but
local minima need not be global. Convexity upgrades the local
statement to a global one. A function is \engterm{convex} on an
interval if $f''(x) \geq 0$ throughout, equivalently if every chord
lies on or above the graph, equivalently if $f(\lambda a + (1 -
\lambda) b) \leq \lambda f(a) + (1 - \lambda) f(b)$ for all $a, b$
in the interval and $\lambda \in [0, 1]$. For a convex function, any
critical point is a global minimum, and the minimum is unique if the
convexity is strict ($f'' > 0$). This is the property that makes
convex optimisation tractable: the first-order condition $f'(x) = 0$
locates the global optimum directly, with no search over multiple
local minima and no need for the second-order test as a tie-breaker.
The three convexity criteria are equivalent for twice-differentiable
functions, and an engineer checks whichever is cheapest: the
second-derivative sign when a formula is in hand, the chord test
when only sample points are available, the definition when the
function is given abstractly. The economic-order-quantity cost
$C(Q) = KD/Q + hQ/2$ of this section is convex on $Q > 0$ (its
second derivative $2KD/Q^{3}$ is positive), which is why its single
critical point is the global minimum and why the square-root formula
is the answer with no further checking. Chapter~15 builds the
multivariable theory, where convexity is the condition that the
Hessian matrix is positive semidefinite and the global-optimum
guarantee carries over to many design variables at once.
\end{mastery}

\section{Numerical differentiation, why it is hard}\pathtag{core}

Up to this point we have computed derivatives analytically. In
many engineering settings the function $f$ is not given by a
formula but by a sequence of measurements, a black-box simulation,
or a tabulated lookup. The derivative must be approximated from
function values alone. \engterm{Numerical differentiation} is the
discipline of doing this on a finite-precision machine.

The discipline is harder than its analytical counterpart, for a
reason we will derive: numerical differentiation amplifies noise
and floating-point error in a way that integration, in
chapter~6, does not. The engineer who does it badly produces a
derivative whose error is many times the function's error; the
engineer who does it well produces a derivative whose error is at
the floor set by the function's measurement uncertainty.

\subsection{Forward, backward, and central differences}

The simplest approximation to $f'(x)$ uses two function values.
The forward difference is

\[
D_{+}(x; h) = \frac{f(x + h) - f(x)}{h}.
\]

This is the difference quotient itself, evaluated at finite $h$
rather than in the limit. The backward difference is

\[
D_{-}(x; h) = \frac{f(x) - f(x - h)}{h},
\]

which uses the function value to the left of $x$. The central
difference is

\[
D_{0}(x; h) = \frac{f(x + h) - f(x - h)}{2 h}.
\]

\subsection{Truncation error: the bias}

Each formula has a truncation error: the difference between the
formula's output and the exact $f'(x)$, computed assuming infinite
precision arithmetic. We compute it from Taylor's theorem.

For the forward difference, expand $f(x + h)$ around $x$:

\[
f(x + h) = f(x) + f'(x) h + \tfrac{1}{2} f''(x) h^{2} +
\tfrac{1}{6} f'''(\xi) h^{3},
\]

for some $\xi$ between $x$ and $x + h$. Substituting and dividing
by $h$,

\[
D_{+}(x; h) = f'(x) + \tfrac{1}{2} f''(x) h + \tfrac{1}{6}
f'''(\xi) h^{2}.
\]

The truncation error of the forward difference is approximately
$\tfrac{1}{2} f''(x) h$ for $h$ small, which scales as $\oom{h}$.
The forward difference is a first-order method.

The backward difference has the same first-order truncation error,
with $f''$ evaluated at $x$ minus a small correction.

For the central difference, expand both $f(x + h)$ and $f(x - h)$
around $x$:

\[
f(x \pm h) = f(x) \pm f'(x) h + \tfrac{1}{2} f''(x) h^{2} \pm
\tfrac{1}{6} f'''(x) h^{3} + \tfrac{1}{24} f''''(\xi^{\pm}) h^{4}.
\]

The difference $f(x + h) - f(x - h)$ has the even-order terms
cancel and the odd-order terms double:

\[
f(x + h) - f(x - h) = 2 f'(x) h + \tfrac{1}{3} f'''(x) h^{3} +
\oom{h^{5}}.
\]

Dividing by $2h$,

\[
D_{0}(x; h) = f'(x) + \tfrac{1}{6} f'''(x) h^{2} + \oom{h^{4}}.
\]

The truncation error of the central difference scales as $\oom
{h^{2}}$, one order better than the forward and backward
differences. Wherever the function is smooth and the values can
be computed at $x \pm h$, the central difference is the engineer's
default.

\subsection{Round-off error: the variance}

Truncation error is not the only error in numerical
differentiation. The function values $f(x \pm h)$ are computed in
finite precision: a 64-bit floating-point number stores $f$ to
about $10^{-16}$ relative accuracy. If $f(x) \sim O(1)$, the
absolute error in each function value is $\varepsilon \sim
10^{-16}$. When we form the central difference

\[
D_{0}(x; h) = \frac{f(x + h) - f(x - h)}{2 h},
\]

the numerator is the difference of two numbers that are very close
to each other when $h$ is small. The function values
$f(x \pm h)$ each carry an error of order $\varepsilon$; the
difference carries an error of order $\varepsilon$ (in the worst
case, the errors do not cancel); the divided difference carries
an error of order $\varepsilon / h$.

The total error in $D_{0}(x; h)$ as an estimate of $f'(x)$ is then

\[
|D_{0}(x; h) - f'(x)| \lesssim \tfrac{1}{6} |f'''(x)| h^{2} +
\frac{C \varepsilon}{h},
\]

with $C$ a small constant depending on how the function value is
computed. The truncation term decreases with $h$; the round-off
term increases with $h$. The total has a minimum.

Setting the derivative of the total error with respect to $h$ to
zero, the optimal step is $h^{*} \sim (\varepsilon / |f'''|)^{1/3}$
for the central difference (and $h^{*} \sim (\varepsilon /
|f''|)^{1/2}$ for the forward difference). The optimal error scales
as $\varepsilon^{2/3}$ for the central difference and as
$\varepsilon^{1/2}$ for the forward difference. With double-
precision $\varepsilon \approx 10^{-16}$, the central-difference
optimal error is about $10^{-11}$ and the forward-difference
optimal error is about $10^{-8}$. A naive engineer who picks $h =
10^{-8}$ for the central difference, or $h = 10^{-10}$ for the
forward difference, will lose precision rapidly to round-off.

This is \engterm{catastrophic cancellation}: the loss of significant
digits when subtracting two nearly equal numbers. It is the
operative noise mechanism of numerical differentiation, and it is
the reason numerical differentiation is harder than numerical
integration.

\begin{figure}[ht]
\centering
\begin{tikzpicture}[x=1cm, y=1cm, font=\small]

% Log-log error vs h, central difference, double precision
% h axis: 10^-15 to 10^0, mapped to plot x from 0 to 12 (each decade = 0.8)
% error axis: 10^-15 to 10^0, mapped to plot y from 0 to 9 (each decade = 0.6)

% Axes
\draw[->, thick] (0, 0) -- (12.5, 0) node[right] {$h$};
\draw[->, thick] (0, 0) -- (0, 9.5) node[above] {absolute error $|D_{0}(x;h) - f'(x)|$};

% x-axis log decade ticks (h = 10^-15 ... 10^0, plotx = (15 + log10(h)) * 0.8)
\foreach \k in {-15, -12, -9, -6, -3, 0} {
  \pgfmathsetmacro{\px}{(15 + \k)*0.8}
  \draw (\px, 0) -- (\px, -0.15);
  \node[anchor=north, font=\scriptsize] at (\px, -0.15) {$10^{\k}$};
}

% y-axis log decade ticks
\foreach \k in {-15, -12, -9, -6, -3, 0} {
  \pgfmathsetmacro{\py}{(15 + \k)*0.6}
  \draw (0, \py) -- (-0.15, \py);
  \node[anchor=east, font=\scriptsize] at (-0.15, \py) {$10^{\k}$};
}

% Round-off branch: error_round ~ eps / h, eps = 1e-16
% log10(error) = log10(eps) - log10(h) = -16 - log10(h)
% h goes from 10^-15 to 10^-5 for visible round-off branch
% plotx = (15 + log10(h))*0.8; ploty = (15 + (-16 - log10(h)))*0.6 = (-1 - log10(h))*0.6
% For log10(h) = -15: plotx = 0, ploty = 14*0.6 = 8.4
% For log10(h) = -5: plotx = 8, ploty = 4*0.6 = 2.4
\draw[thick, engred]
  plot[domain=-15:-5, samples=50] ({(15 + \x)*0.8}, {(-1 - \x)*0.6});

% Truncation branch: error_trunc ~ (1/6) h^2 |f'''|, take |f'''| ~ 1
% log10(error) = -0.78 + 2 log10(h)
% h goes from 10^-6 to 10^0
% For log10(h) = -6: plotx = 9*0.8 = 7.2, ploty = (15 + (-0.78 - 12))*0.6 = 2.22*0.6 = 1.33
% For log10(h) = 0:  plotx = 15*0.8 = 12, ploty = (15 - 0.78)*0.6 = 14.22*0.6 = 8.5
\draw[thick, engblue]
  plot[domain=-6:0, samples=20] ({(15 + \x)*0.8}, {(15 + (-0.78 + 2*\x))*0.6});

% Optimal point: h* = (eps / |f'''|)^(1/3) ~ 10^(-16/3) ~ 10^-5.33
% error at optimum ~ eps^(2/3) ~ 10^(-10.67)
\fill[engblack] ({(15 - 5.33)*0.8}, {(15 - 10.67)*0.6}) circle (2.5pt);
\node[engblack, anchor=south west, font=\scriptsize, align=left]
  at ({(15 - 5.33)*0.8 + 0.2}, {(15 - 10.67)*0.6 + 0.1})
  {$h^{*} \sim \varepsilon^{1/3}$\\ error $\sim \varepsilon^{2/3}$};

% Label the two branches
\node[engred, anchor=west, font=\scriptsize] at (1.0, 7.5)
  {round-off $\sim \varepsilon/h$};
\node[engblue, anchor=east, font=\scriptsize] at (11.0, 5.5)
  {truncation $\sim h^{2} |f'''|/6$};

% Floor annotation
\draw[thin, gray, dashed]
  (0, {(15 - 10.67)*0.6}) -- (12, {(15 - 10.67)*0.6});
\node[gray, anchor=west, font=\tiny] at (10, {(15 - 10.67)*0.6 - 0.3})
  {optimum};

\end{tikzpicture}
\caption[Bias-variance tradeoff in numerical differentiation]{Total
error of the central-difference estimate $D_{0}(x; h)$ for $f'(x)$
in double precision ($\varepsilon \approx 10^{-16}$). The round-off
branch (red) scales as $\varepsilon / h$ and dominates for small $h$;
the truncation branch (blue) scales as $h^{2} |f'''|/6$ and dominates
for large $h$. The total error is minimised at $h^{*} \sim
\varepsilon^{1/3} \approx 10^{-5}$ (for $|f'''| \sim 1$), with
optimum error of order $\varepsilon^{2/3} \approx 10^{-11}$. The
plot collapses the entire tradeoff onto a single V-shaped curve;
the engineer reads $h^{*}$ off the bottom and the achievable
precision off the left axis.}
\label{fig:vol02:ch05:bias-variance}
\end{figure}


\Cref{fig:vol02:ch05:bias-variance} collapses the two error
contributions onto a single log-log plot. The V-shape is universal:
every numerical-differentiation rule has a truncation branch that
falls as a power of $h$ and a round-off branch that rises as
$\varepsilon / h$. The position of the V's apex slides with the
order of the rule and with the precision of the arithmetic; the
shape itself is invariant. The companion script
\texttt{code/finite\_diff\_accuracy.py} produces the table that
underlies the figure for $f(x) = e^{x}$ at $x = 1$.

\subsection{The bias-variance tradeoff}

The structure of the total error reveals a tradeoff that recurs
throughout numerical analysis and statistics: bias against
variance, or in this language, truncation error against round-off
error. Reducing $h$ reduces the bias (the systematic, function-
dependent component of the error) but increases the variance (the
noise-driven component). A larger $h$ reduces variance but
increases bias.

The same tradeoff reappears, with statistics and noisy data
replacing floating-point arithmetic, when the engineer differentiates
a measured time series. Smaller windows and shorter smoothing
filters reduce bias against the underlying signal but amplify the
measurement noise. Larger windows reduce noise but smear
genuine signal features. The chapter's project (section~5.10)
walks the reader through this tradeoff on a real signal.

\subsection{Higher-order finite differences}

For greater accuracy at fixed $h$, the engineer can use more
function values. A symmetric four-point stencil gives

\[
D_{4}(x; h) = \frac{-f(x + 2h) + 8 f(x + h) - 8 f(x - h) +
f(x - 2h)}{12 h},
\]

with truncation error $\oom{h^{4}}$. Each additional function
value buys two additional orders of accuracy (because the central-
symmetry of the stencil cancels another pair of even-order Taylor
terms). The cost is that the round-off problem is similar: the
numerator is the difference of nearly equal terms, so the noise
amplification by $1/h$ remains. The optimal step now scales as
$\varepsilon^{1/5}$ and the optimal error as $\varepsilon^{4/5}$,
better than the two-point central difference but at the cost of
more function evaluations and a wider stencil.

For a function that can be evaluated cheaply and is smooth, the
four-point stencil is a good default. For a function that is
expensive to evaluate (a long simulation, a slow measurement), the
two-point central difference is the better trade. For functions
with derivative discontinuities, all stencils break and section~5.8
applies.

\begin{mastery}
Finite differences and the complex-step method both approximate the
derivative by evaluating the function at nearby points.
\engterm{Automatic differentiation} computes the derivative exactly,
to machine precision, by propagating derivative information through
the arithmetic operations of the program that evaluates the
function, with no step size at all. The forward mode carries a
\engterm{dual number} $x + \dot{x}\,\epsilon$ with $\epsilon^{2} =
0$ through every operation: the rules $\,(a + \dot{a}\epsilon)(b +
\dot{b}\epsilon) = ab + (a\dot{b} + \dot{a}b)\epsilon$ and
$\sin(a + \dot{a}\epsilon) = \sin a + \dot{a}\cos a\,\epsilon$ are
the product rule and chain rule of this chapter, executed
automatically by overloaded arithmetic. The coefficient of
$\epsilon$ in the result is the exact derivative. Seeding $\dot{x}
= 1$ at the input and reading the $\epsilon$-coefficient at the
output evaluates $f$ and $f'$ in a single pass, with the chain rule
applied operation by operation rather than symbolically. The
reverse mode, the basis of backpropagation in machine learning,
runs the chain rule from the output backward and computes the
derivative of one output with respect to many inputs in one sweep,
which is why training a network with millions of parameters is
feasible. Automatic differentiation has neither the truncation bias
of finite differences nor the round-off floor; its only error is
the floating-point rounding of the exact arithmetic it performs.
The complex-step derivative of the previous box is the poor cousin
of forward-mode automatic differentiation: the imaginary unit plays
the role of $\epsilon$, except that $i^{2} = -1$ rather than $0$, so
the complex-step method carries a small $h^{2}$ truncation that true
dual numbers do not. Chapter~16 develops automatic differentiation
as the modern derivative engine for machine-evaluated functions; the
rules it automates are the product rule and chain rule derived in
section~5.2.
\end{mastery}

\begin{mastery}
The bias-variance V has a clean escape when the function is
real-analytic and can be evaluated in complex arithmetic. The
\engterm{complex-step derivative} evaluates the function at a
complex argument $x + i h$ and reads the derivative off the
imaginary part:

\[
f'(x) \approx \frac{\operatorname{Im} f(x + i h)}{h}.
\]

The expansion $f(x + i h) = f(x) + i h f'(x) - \tfrac{1}{2} h^{2}
f''(x) - \cdots$ shows the imaginary part is $h f'(x) - \tfrac{1}{6}
h^{3} f'''(x) + \cdots$, so dividing by $h$ gives $f'(x)$ with
truncation error $\oom{h^{2}}$, the accuracy of a central
difference. The advantage is that no subtraction of nearly equal
numbers occurs: the formula has no difference in the numerator, so
catastrophic cancellation is absent, and $h$ may be taken as small
as $10^{-200}$ with the derivative accurate to machine precision.
The cost is that the function must be coded to accept complex
arguments and must be analytic (no $\operatorname{abs}$, no
branch on the real part, no comparison operators that break
analyticity). Where it applies, the complex-step derivative removes
the step-size tuning problem entirely. The accompanying script
\texttt{code/complex\_step.py} tabulates the complex-step error
against the central-difference error for $f(x) = e^{x}$ and shows
the complex-step error falling monotonically as $h^{2}$ with no
round-off floor. The technique is the simplest member of the
automatic-differentiation family, which chapter~16 develops as the
modern alternative to finite differences for machine-evaluated
functions.
\end{mastery}

\subsection{A worked accuracy study: finite difference against
complex step}

The contrast between finite differencing and the complex-step
derivative is worth a side-by-side study, because it shows the
round-off floor appearing and disappearing under a controlled
change of method. We differentiate $f(x) = e^{x}$ at $x = 1$, where
the exact answer is $f'(1) = e = 2.718281828\ldots$, and sweep the
step size $h$ from $10^{-1}$ down to $10^{-16}$ for three schemes:
the forward difference, the central difference, and the complex
step.

\engterm{Forward difference}. The truncation error is
$\tfrac{1}{2} f''(1) h = \tfrac{1}{2} e h$, falling linearly with
$h$; the round-off error is of order $\varepsilon/h$ with
$\varepsilon \approx 2 \times 10^{-16}$. The two cross at $h^{*}
\approx \sqrt{2\varepsilon/|f''|} = \sqrt{2 \times 10^{-16}/e}
\approx 10^{-8}$, where the total error bottoms out near $10^{-8}$.
At $h = 10^{-2}$ the error is dominated by truncation and reads
about $1.4 \times 10^{-2}$; at $h = 10^{-14}$ it is dominated by
round-off and reads about $10^{-2}$ again; at the optimum near
$h = 10^{-8}$ it dips to roughly $10^{-8}$.

\engterm{Central difference}. The truncation error is $\tfrac{1}{6}
f'''(1) h^{2} = \tfrac{1}{6} e h^{2}$, falling quadratically; the
round-off error is again of order $\varepsilon/h$. The crossover is
at $h^{*} \approx (\varepsilon/|f'''|)^{1/3} \approx 10^{-5.5}$,
where the total error reaches about $10^{-11}$. The central scheme
buys three extra decades of accuracy at its optimum over the
forward scheme, for one extra function evaluation per derivative.

\engterm{Complex step}. Evaluating $f(1 + ih) = e^{1+ih} = e
(\cos h + i \sin h)$, the imaginary part is $e \sin h$, and
$\operatorname{Im} f(1 + ih)/h = e \sin h / h$. The truncation
error is $e(\sin h - h)/h = -e h^{2}/6 + \cdots$, of order $h^{2}$,
matching the central difference. The decisive difference is that
the numerator $\operatorname{Im} f(1 + ih)$ contains no subtraction
of nearly equal real numbers; the imaginary part is computed
directly. There is no round-off arm. As $h$ shrinks, the error
falls as $h^{2}$ until it meets the machine-precision floor near
$10^{-16}$, and then it stays there for every smaller $h$,
including $h = 10^{-50}$ or $h = 10^{-200}$. The step size simply
needs to be small enough; there is no optimum to tune and no
penalty for overshooting.

\begin{figure}[ht]
\centering
\begin{tikzpicture}[x=1cm, y=1cm, font=\small]

% Error versus step size on log-log axes for three derivative schemes
% on a smooth analytic function: forward difference (V apex high),
% central difference (V apex lower), complex-step (monotone fall to a
% machine-precision floor with no round-off arm).
% x-axis: log10(h) from -16 (left) to 0 (right). Map: x = 0.55*(log h +16).
%   log h = -16 -> x = 0; log h = 0 -> x = 8.8.
% y-axis: log10(error) from 0 (top) to -16 (bottom).
%   Map: y = 5.0 + 0.3*log err. log err = 0 -> y = 5.0; -16 -> y = 0.2.

\draw[->, thick] (-0.3, 0) -- (9.4, 0)
  node[right] {$\log_{10} h$};
\draw[->, thick] (0, -0.3) -- (0, 5.6)
  node[above, align=center] {$\log_{10}(\text{error})$};

% x ticks at log h = -16, -12, -8, -4, 0
\foreach \lh in {-16,-12,-8,-4,0} {
  \pgfmathsetmacro\xx{0.55*(\lh + 16)}
  \draw (\xx, 0) -- (\xx, -0.12);
  \node[anchor=north, font=\scriptsize] at (\xx, -0.12) {\lh};
}
% y ticks at log err = 0, -4, -8, -12, -16
\foreach \le in {0,-4,-8,-12,-16} {
  \pgfmathsetmacro\yy{5.0 + 0.3*\le}
  \draw (0, \yy) -- (-0.12, \yy);
  \node[anchor=east, font=\scriptsize] at (-0.14, \yy) {\le};
}

% Forward difference: trunc ~ h (slope +1 on right), round-off ~ eps/h
% (slope -1 on left), apex near log h = -8, error ~ -8.
% Build from sample points (log h, log err).
\draw[very thick, engred, dotted]
  plot coordinates {
    ({0.55*(-16+16)}, {5.0 + 0.3*(0)})
    ({0.55*(-12+16)}, {5.0 + 0.3*(-4)})
    ({0.55*(-8+16)},  {5.0 + 0.3*(-8)})
    ({0.55*(-6+16)},  {5.0 + 0.3*(-6)})
    ({0.55*(-2+16)},  {5.0 + 0.3*(-2)})
    ({0.55*(0+16)},   {5.0 + 0.3*(0)})
  };
\node[engred, anchor=west, font=\scriptsize] at (5.2, 4.7)
  {forward diff.};

% Central difference: trunc ~ h^2 (slope +2 on right), round-off ~ eps/h
% (slope -1), apex near log h = -5.3, error ~ -11.
\draw[very thick, engamber, dashed]
  plot coordinates {
    ({0.55*(-16+16)}, {5.0 + 0.3*(-1)})
    ({0.55*(-12+16)}, {5.0 + 0.3*(-5)})
    ({0.55*(-8+16)},  {5.0 + 0.3*(-9)})
    ({0.55*(-5.3+16)},{5.0 + 0.3*(-11)})
    ({0.55*(-2+16)},  {5.0 + 0.3*(-4.6)})
    ({0.55*(0+16)},   {5.0 + 0.3*(-0.6)})
  };
\node[engamber, anchor=west, font=\scriptsize] at (5.2, 4.3)
  {central diff.};

% Complex step: trunc ~ h^2 (slope +2 on right), no round-off arm; the
% error falls until it hits the machine floor ~ -16 and stays there.
\draw[very thick, engblue]
  plot coordinates {
    ({0.55*(-16+16)}, {5.0 + 0.3*(-16)})
    ({0.55*(-12+16)}, {5.0 + 0.3*(-16)})
    ({0.55*(-8+16)},  {5.0 + 0.3*(-16)})
    ({0.55*(-6+16)},  {5.0 + 0.3*(-12.6)})
    ({0.55*(-2+16)},  {5.0 + 0.3*(-4.6)})
    ({0.55*(0+16)},   {5.0 + 0.3*(-0.6)})
  };
\node[engblue, anchor=west, font=\scriptsize] at (5.2, 0.55)
  {complex step};

% machine floor guide line
\draw[thin, enggray, densely dotted] (0, {5.0 + 0.3*(-16)})
  -- (8.8, {5.0 + 0.3*(-16)});

\end{tikzpicture}
\caption[Finite-difference and complex-step accuracy against step
size]{Absolute error of three first-derivative schemes on a smooth
analytic function, as the step size $h$ sweeps from $10^{0}$ down to
$10^{-16}$. The forward and central differences each show the
characteristic V: a truncation arm that falls as a power of $h$ on the
right and a round-off arm that rises as $\varepsilon/h$ on the left,
with the central scheme's apex one order lower because its truncation
falls as $h^{2}$ rather than $h$. The complex-step derivative
$\operatorname{Im} f(x + ih)/h$ has no subtraction of nearly equal
numbers, so it carries no round-off arm: the error falls as $h^{2}$
and then sits on the machine-precision floor for every smaller $h$.
The step-size tuning problem that dominates finite differencing
disappears for the complex-step method wherever the function can be
evaluated in complex arithmetic.}
\label{fig:vol02:ch05:complex-step}
\end{figure}


\Cref{fig:vol02:ch05:complex-step} plots the three error curves on
log-log axes. The forward and central differences show the V; the
complex-step curve has only the falling arm and the floor. The
practical reading: when the function can be evaluated in complex
arithmetic and is analytic over the path (no \texttt{abs}, no
branch on the real part, no comparison that breaks analyticity),
the complex-step derivative removes the step-size tuning problem
entirely and delivers a derivative accurate to machine precision.
When the function cannot be so evaluated, the central difference at
its optimum is the working default, and the engineer pays the
attention the V demands. The companion script
\texttt{code/complex\_step.py} produces the table that underlies
the figure; \texttt{code/finite\_diff\_accuracy.py} produces the
forward and central arms.

\begin{mastery}
The earlier box introduced forward-mode automatic differentiation, which
carries one derivative through the program alongside each value. Its cost
structure decides when it is the wrong tool. Forward mode computes the
derivative of all outputs with respect to one input in a single sweep, at
a cost of a small constant factor over one function evaluation. To obtain
the derivative with respect to $n$ inputs, forward mode runs $n$ sweeps,
so its cost scales as $n$ times the function cost. The
\engterm{reverse mode} inverts this trade. It records the operations of
the function on a tape during a forward pass, then walks the tape
backward applying the chain rule from the output toward the inputs, and
in that one backward pass it produces the derivative of one output with
respect to all $n$ inputs at once. Its cost is a small constant factor
over a single function evaluation, independent of $n$. The structural
asymmetry is decisive for optimisation: a cost function maps many
parameters to one scalar, so reverse mode delivers the entire gradient
for the price of a few function evaluations, while forward mode would
pay $n$ times over. This is why training a model with millions of
parameters uses reverse-mode differentiation, called backpropagation in
that setting, and why a Jacobian of an $m$-output, $n$-input map is
assembled column by column with forward mode when $n < m$ and row by row
with reverse mode when $m < n$. The cost of reverse mode is memory: the
tape holds every intermediate value of the forward pass so the backward
pass can reuse it, so a long computation can exhaust memory before it
exhausts time, and checkpointing schemes trade recomputation for storage
to fit the tape. The chain rule of section~5.2 is the same in both modes;
the modes differ only in the order they traverse the composition, and
that order sets the cost.

A different deepening answers a different question: not how accurate the
derivative is on average, but how to bound it with certainty.
\engterm{Interval} or \engterm{validated} differentiation carries an
interval $[\underline{x}, \overline{x}]$ rather than a point through the
arithmetic, with every operation rounded outward so the true result is
guaranteed to lie in the returned interval. Combined with the dual-number
rules of automatic differentiation, interval arithmetic produces an
interval that is guaranteed to contain the exact derivative over the
entire input interval, not merely an estimate at one point. The
guarantee is what some applications require: a global optimiser that must
prove no minimum was missed uses interval derivatives to certify that a
region contains no stationary point (the derivative interval does not
contain zero) and can be discarded, and a safety argument that must bound
the worst-case sensitivity over a range uses the interval derivative as a
rigorous envelope. The cost is conservatism: interval arithmetic ignores
the correlations between an expression and itself, so the returned
interval can be far wider than the true range, and the technique repays
careful formulation to keep the bounds tight. Where a point estimate of
the derivative suffices, finite differences or point-mode automatic
differentiation are the working tools; where a certified enclosure is
required, the interval derivative is the one that carries a proof.
\end{mastery}

\subsection{Differentiating noisy data}

When the function is a measured time series, the noise in each
sample sets the floor for $h$. If the standard deviation of the
measurement noise is $\sigma$, the round-off-equivalent error
floor in the central difference is $\sigma / h$, and the optimal
$h$ scales as $(\sigma / |f'''|)^{1/3}$. For a signal with
$|f'''| \sim 1$ and $\sigma \sim 10^{-3}$, the optimal $h$ is
roughly $0.1$, and the optimal error in the derivative estimate
is roughly $10^{-2}$. The optimum is conservative: the engineer
should take a small range of $h$ values around the predicted
optimum, compute the derivative at each, and look for a plateau.
The width of the plateau and the value of the plateau set the
defensible derivative estimate.

The alternative, when the engineer knows something about the
underlying signal, is to fit a smooth model to the data first
(polynomial, spline, Savitzky-Golay filter) and differentiate the
model. The model's smoothness controls the variance, and the
goodness of the fit controls the bias. We return to this in the
project of section~5.10.

\section{Failure: discontinuities, kinks, derivatives that do not exist}\pathtag{core}

The chapter's apparatus assumes the function is smooth. We now
catalogue the cases where the assumption fails and the apparatus
breaks. The catalogue is the engineer's guard against using a
calculus tool on a function that does not admit it.

\begin{figure}[ht]
\centering
\begin{tikzpicture}[x=1cm, y=1cm, font=\small]

% Four small panels: jump, kink, cusp, oscillation. Each 3.2 wide.
\def\pw{3.4}

% --- Panel 1: jump ---
\begin{scope}[shift={(0,0)}]
  \draw[->, thick] (-1.4, 0) -- (1.6, 0) node[right, font=\scriptsize] {$x$};
  \draw[->, thick] (0, -1.4) -- (0, 1.6);
  \draw[thick, engblue] (-1.3, -0.9) -- (0, -0.3);
  \draw[thick, engblue] (0, 0.5) -- (1.3, 1.1);
  \fill[engblue] (0, -0.3) circle (1.4pt);
  \draw[engblue] (0, 0.5) circle (1.4pt);
  \node[font=\scriptsize, anchor=north] at (0.1, -1.5) {jump: $f'$ undefined};
\end{scope}

% --- Panel 2: kink |x| ---
\begin{scope}[shift={(\pw,0)}]
  \draw[->, thick] (-1.4, 0) -- (1.6, 0) node[right, font=\scriptsize] {$x$};
  \draw[->, thick] (0, -0.3) -- (0, 1.7);
  \draw[thick, engred] (-1.3, 1.3) -- (0, 0) -- (1.3, 1.3);
  \fill[engred] (0,0) circle (1.4pt);
  \node[font=\scriptsize, anchor=north] at (0.1, -0.4) {kink: $f'(0^{-})\neq f'(0^{+})$};
\end{scope}

% --- Panel 3: cusp x^{2/3} ---
\begin{scope}[shift={(2*\pw,0)}]
  \draw[->, thick] (-1.4, 0) -- (1.6, 0) node[right, font=\scriptsize] {$x$};
  \draw[->, thick] (0, -0.3) -- (0, 1.7);
  \draw[thick, engamber] plot[domain=0.0008:1.3, samples=60] (\x, {1.2*pow(\x,0.6667)});
  \draw[thick, engamber] plot[domain=0.0008:1.3, samples=60] (-\x, {1.2*pow(\x,0.6667)});
  \fill[engamber] (0,0) circle (1.4pt);
  \node[font=\scriptsize, anchor=north] at (0.1, -0.4) {cusp: $f'\to\pm\infty$};
\end{scope}

% --- Panel 4: oscillation x sin(1/x) ---
\begin{scope}[shift={(3*\pw,0)}]
  \draw[->, thick] (-1.4, 0) -- (1.6, 0) node[right, font=\scriptsize] {$x$};
  \draw[->, thick] (0, -1.4) -- (0, 1.6);
  \draw[thick, enggray] plot[domain=0.03:1.3, samples=400] (\x, {1.0*\x*sin(deg(1/\x))});
  \draw[thick, enggray] plot[domain=0.03:1.3, samples=400] (-\x, {-1.0*\x*sin(deg(1/\x))});
  \fill[enggray] (0,0) circle (1.4pt);
  \node[font=\scriptsize, anchor=north] at (0.1, -1.5) {oscillation: quotient never settles};
\end{scope}

\end{tikzpicture}
\caption[Four ways differentiability fails]{The four canonical
non-differentiable points. A jump (left) makes the difference
quotient unbounded as $h\to 0$. A kink (centre-left, here $|x|$) has
one-sided derivatives that disagree. A cusp (centre-right, here
$x^{2/3}$) has a vertical tangent and an unbounded one-sided
derivative. A wild oscillation (right, here $x\sin(1/x)$) is
continuous yet the difference quotient never settles to a limit at
the origin. Each has engineering instances catalogued in
section~5.8; in every case the calculus tool must be replaced by a
tool matched to the local structure.}
\label{fig:vol02:ch05:nondifferentiable}
\end{figure}


\Cref{fig:vol02:ch05:nondifferentiable} sets the four canonical
failures side by side: jump, kink, cusp, and wild oscillation. Each
panel is a function for which the difference quotient fails to
settle to a single finite limit at the marked point, and each has
the engineering instances developed below.

\subsection{Jump discontinuities}

A function with a jump at $x = a$ has $\lim_{h \to 0^{+}} f(a +
h) \neq \lim_{h \to 0^{-}} f(a + h)$. The forward and backward
difference quotients at $x = a$ both diverge as $h \to 0$, and
the derivative does not exist at $a$. The function $f(x) = x$ for
$x \geq 0$, $f(x) = x - 1$ for $x < 0$, has a jump of magnitude
$1$ at $x = 0$.

Engineering examples include phase-change boundaries (a material's
heat capacity is discontinuous across a phase transition), saturating
amplifiers (the output is one value below saturation and another
above), digital control signals (a step input from a controller),
and accounting boundaries (tax brackets, regulatory thresholds).
At a jump, the derivative is not defined; the engineer must work
with one-sided derivatives or with a generalised derivative
concept (the Dirac delta at the jump).

\subsection{Kinks: continuous but not differentiable}

A function with a kink at $a$ is continuous (no jump) but has
$f'(a^{-}) \neq f'(a^{+})$: the left and right derivatives exist
but are unequal, and the derivative at $a$ is not defined.

The canonical example is $f(x) = |x|$. The function is continuous
at $x = 0$, with $f(0) = 0$. The right derivative is $\lim_{h \to
0^{+}} (|h| - 0)/h = 1$, the left derivative is $\lim_{h \to
0^{-}} (|h| - 0)/h = -1$. The two are unequal; the derivative at
$0$ does not exist.

Engineering examples: rectifier outputs ($v_{\text{out}}(v_{\text
{in}}) = \max(0, v_{\text{in}})$), tension cables (cable force
is zero in compression, positive in tension, with a kink at zero
elongation), surface contact stresses (Hertzian contact has a
square-root singularity in the contact-pressure derivative at the
contact-zone edge), elastic-perfectly-plastic material models (the
stress-strain curve has a kink at yield).

For a kink, optimisation methods that use $f'(x) = 0$ may miss
the optimum if the optimum lies at the kink. The function
$f(x) = |x|$ has its global minimum at $x = 0$, where the
derivative is undefined. The minimum is found not by $f'(x) = 0$
but by inspecting the kink directly. A reader who solves
$f'(x) = 0$ on a piecewise-linear function will get the wrong
answer if the optimum is at a corner; the right procedure is to
compute the function at every kink point and at every endpoint
of every smooth piece, and pick the smallest (or largest) value.

\begin{mastery}
A kink defeats the condition $f'(x) = 0$ only because the derivative is a
single number that does not exist at the corner. Convex analysis repairs
the defect by replacing the derivative with a set. The
\engterm{subdifferential} of a convex function $f$ at a point $x$ is the
set of all slopes $s$ for which the line through $(x, f(x))$ with slope
$s$ stays on or below the graph,
\[
\partial f(x) = \{\, s : f(y) \geq f(x) + s\,(y - x) \text{ for all } y
\,\},
\]
the set of \engterm{subgradients}. Where $f$ is differentiable the
subdifferential is the single point $\{f'(x)\}$; at a kink it is the
whole closed interval between the left and right derivatives. For
$f(x) = |x|$ the subdifferential at the origin is the interval
$[-1, 1]$, every slope from the left derivative $-1$ to the right
derivative $+1$. The optimality condition generalises cleanly: a convex
function attains its global minimum at $x^{*}$ exactly when
$0 \in \partial f(x^{*})$, the zero slope belongs to the set. For
$f(x) = |x|$ this reads $0 \in [-1, 1]$, which holds at $x = 0$ and
locates the minimum that the ordinary $f'(x) = 0$ condition missed. The
subgradient is not a curiosity. The $L^{1}$ penalties that drive sparse
regression and compressed sensing are sums of absolute values, nonsmooth
at every coordinate axis, and their minimisation rests entirely on the
subgradient condition. The proximal and subgradient methods that solve
these problems are the nonsmooth analogues of the gradient descent of
section~5.5, with the subgradient standing in for the gradient wherever
the kink lives. The engineer who meets an $L^{1}$ regulariser, a hinge
loss, a worst-case (minimax) objective, or any cost built from maxima and
absolute values is optimising a nonsmooth function, and the subdifferential
is the tool that keeps the first-order condition meaningful where the
ordinary derivative has dissolved.
\end{mastery}

\subsection{Cusps and vertical tangents}

A cusp at $a$ is a point where $f$ is continuous, the derivative
exists in a one-sided sense, and that one-sided derivative is
$\pm\infty$. The function $f(x) = x^{2/3}$ has a cusp at $x = 0$:
the derivative $f'(x) = (2/3) x^{-1/3}$ blows up to $+\infty$ from
both sides as $x \to 0$. The graph has a vertical tangent at $x =
0$, and the function is not differentiable there in the ordinary
sense.

Engineering examples: stress concentrations at sharp corners
(theoretical stress is infinite, physical reality includes a
notch radius, plastic yielding, or fracture), capillary rise at
a corner of a wetting wall, and certain power-law approximations
of phase-transition critical behaviour.

\subsection{Wildly oscillating functions}

A function may be continuous and bounded yet fail to be
differentiable because the difference quotient never settles. The
classic example is $f(x) = x^{2} \sin(1/x)$ for $x \neq 0$ and
$f(0) = 0$. This function is differentiable everywhere, including
at $0$, with $f'(0) = 0$, but $f'$ is not continuous at $0$: the
derivative oscillates between $\pm 1$ as $x \to 0$. A worse
example, $f(x) = x \sin(1/x)$ for $x \neq 0$ and $f(0) = 0$, is
continuous but not differentiable at $0$.

Engineering examples include high-frequency vibrations near
resonance, optical interference patterns near caustics, and the
tail of a fast-oscillating thermal-noise signal. Numerical
differentiation of such signals returns garbage; the right
treatment is to filter or downsample to a frequency band where
the signal admits a useful derivative.

\subsection{Derivatives that exist but are bad working tools}

A function may be differentiable in the technical sense and yet
have a derivative that fluctuates so much over working scales
that the derivative is useless. A noisy time series has a well-
defined derivative on every smooth approximation but the
approximation's choice changes the derivative wildly.

The engineering response is to admit that the derivative is not
the right summary of the signal and to look for a different
quantity: a moving average, a peak height, a frequency-domain
energy in a band, a counted number of events above a threshold.
The discipline is to choose the summary that is robust to the
noise mechanism, not to insist on the derivative when the
function does not support it.

\subsection{The Patriot case as a numerical-differentiation
analogue}

A clean engineering case where accumulated numerical error
defeated a differentiation-like calculation is the Patriot missile
system at Dhahran, Saudi Arabia, on 25 February 1991. We treat the
case in three parts: the technical mechanism, the organisational
context, and the lesson for differentiation and rates. The accident
is recorded in the named-cases registry at
\repopath{docs/research/accidents/patriot-dhahran-1991.md}; the
narrative below follows the registry's verified description, with
all empirical numbers traceable to the GAO report
\cite{acc:gao-patriot-1992}.

\paragraph{Technical mechanism.}
The Patriot's tracking software represented elapsed time as an
integer counter incremented every tenth of a second since the
system was started. When the tracking module needed the elapsed
time in real seconds, it multiplied the counter by a stored
constant intended to represent $1/10$. That constant was held in a
$24$-bit fixed-point register, and the exact decimal $0.1$ has no
exact binary representation; the $24$-bit truncation introduced a
per-conversion error of about $9.5 \times 10^{-8}$ in the
converted seconds value, accumulating linearly with the counter
\cite{acc:gao-patriot-1992}. After approximately $100 \,\si{\hour}$
of continuous operation the counter had reached approximately
$3.6 \times 10^{6}$, and the cumulative conversion error had
reached approximately $0.34 \,\si{\second}$.

The Patriot's range-gate prediction used the converted time to
forecast where an incoming target would be at the next radar
look: predicted position equals last-known position plus velocity
times elapsed time. An incoming Scud missile travelling at
approximately $1{,}676 \,\si{\meter\per\second}$ over a
$0.34 \,\si{\second}$ time error produces a predicted-position
offset of approximately $570 \,\si{\meter}$. The Patriot's range
gate, the radar search window around the predicted position, was
narrower than $570 \,\si{\meter}$; the system therefore failed to
find the target in the expected location, classified the
detection as "no track", and did not engage \cite{acc:gao-patriot-1992}.
A software patch correcting the conversion-truncation error had
been developed and was being distributed to Patriot batteries at
the time of the Dhahran incident. The patch arrived at the
Dhahran battery on 26 February 1991, the day after the strike;
an Iraqi Scud struck a US Army barracks at approximately 20:35
local time on 25 February, killing 28 service members and wounding
approximately 100 others.

\paragraph{Organisational context.}
The GAO report identified contributing institutional factors
beyond the arithmetic defect \cite{acc:gao-patriot-1992}. The
Patriot's original operational concept assumed that batteries
would be cycled through restarts as part of normal duty, which
would reset the counter and prevent the accumulated-error mode.
Gulf War deployment doctrine instructed field operators to keep
the systems running continuously to maintain readiness, an
operationally rational choice given the threat environment but
one that pushed the software outside its design envelope. The
accumulated-error effect had been observed in field tests and
reported through engineering channels, but the technical patch
and the operational guidance ("restart every N hours pending
patch deployment") moved through different distribution channels
and arrived at deployed batteries inconsistently. A battery could
hold the operational guidance without the patch, hold the patch
without the operational guidance, hold neither, or hold both;
the Dhahran battery on 25 February held neither.

The case became canonical in software-engineering training for
several reasons. It is one of the clearest documented examples of
a numerical-precision software error producing a directly
attributable loss of life. It illustrates the failure mode of
accumulated arithmetic error in long-running continuous operation,
a category that recurs in flight-control software, navigation
systems, and physical-process control. It also illustrates the
institutional pattern in which a technical patch and an operational
workaround are distributed through different channels with no
synchronisation, so that a deployed unit may have one without the
other and the gap is not visible from any single vantage in the
organisation.

\paragraph{Lesson for differentiation and rates.}
The chapter's apparatus turns on a small step $h$ in the input
and the response of the function to that step. The Patriot's
range-gate calculation is the same architecture: the relevant
quantity is velocity times an elapsed time, a difference between
two large position estimates divided implicitly by the time the
target took to traverse them. The arithmetic that produced the
elapsed-time figure carried a small per-conversion error that
accumulated linearly with the counter. The differenced quantity
(predicted minus measured position) inherited the full accumulated
error, and the system's downstream decision (engage or do not
engage) had no further input that could compensate. Numerical
differentiation has the same architecture in miniature: it
differences nearby function values and divides by a small step;
each step is a chance for catastrophic cancellation or for
accumulated representation error. The defensive disciplines are
the same: track the precision through every operation, choose the
step size against the noise floor, validate the calculation
against an independent computation, and bound the operating
regime against the assumptions of the model.

The Patriot lesson generalises: a measurement chain whose output
depends on a small difference of two large numbers is sensitive
to the precision with which each is represented and to the time
over which any per-operation error can accumulate. A reader who
writes a control loop that subtracts a stored setpoint from a
read sensor value, or a navigation system that integrates a
velocity estimate over a long mission, or a financial system that
discounts a small daily return over a long horizon, is in the
same numerical territory the Patriot occupied. The engineering
discipline is to compute the per-operation error explicitly, sum
it over the working duration, compare the total to the operating
tolerance, and choose either a higher-precision representation,
a periodic reset, or a model that is insensitive to the
accumulating quantity.

\subsection{The principle, restated}

\begin{principle}[The derivative serves the function it differentiates]
A derivative is a tool for analysing a function over a regime in
which the function is smooth enough to admit one. Outside that
regime (jumps, kinks, cusps, oscillations, accumulated round-off,
or measurement noise overwhelming the signal), the derivative is
either undefined or so noisy that it ceases to be informative.
The engineering action is to identify the regime before applying
the tool, and to choose a different tool when the regime fails.
\end{principle}

\section{Worked examples across domains}\pathtag{standard}

We close the analytical apparatus with worked examples from four
engineering domains. Each takes a working quantity, identifies the
derivative the engineer wants, computes it, and reads the answer.

\subsection{Kinematics: position to velocity to acceleration}

A small mass on a vertical spring oscillates with displacement
$y(t) = A \cos(\omega t + \phi)$, where $A$ is the amplitude,
$\omega$ is the angular frequency, $\phi$ is the phase. The
velocity is the derivative with respect to time:

\[
v(t) = \dv{y}{t} = -A \omega \sin(\omega t + \phi).
\]

The acceleration is the second derivative:

\[
a(t) = \dv{v}{t} = -A \omega^{2} \cos(\omega t + \phi) =
-\omega^{2} y(t).
\]

The acceleration is proportional to the displacement and opposite
in sign. This is the defining equation of simple harmonic motion:
$\ddot{y} + \omega^{2} y = 0$. The peak acceleration is $A
\omega^{2}$, occurring when the displacement is at its extremes
and the velocity is zero. The peak velocity is $A \omega$,
occurring when the displacement is zero. A reader designing a
mechanical system that must withstand vibration uses the peak
acceleration, not the peak displacement, to size the components.
Volume III takes up vibration formally; the calculus enters here.

\subsection{Electrical engineering: capacitor current}

A capacitor of capacitance $C$ holds a charge $q(t) = C v(t)$
where $v(t)$ is the voltage across it. The current into the
capacitor is the rate of change of charge:

\[
i(t) = \dv{q}{t} = C \dv{v}{t}.
\]

The capacitor current is proportional to the rate of change of
the voltage, not to the voltage itself. A capacitor with a
constant voltage carries zero current; a capacitor with a fast-
changing voltage carries a large current. This is the defining
linear-element relationship of the capacitor.

For a sinusoidal voltage $v(t) = V_{0} \cos(\omega t)$, the current
is $i(t) = -C V_{0} \omega \sin(\omega t)$. The current leads the
voltage by $\pi/2$ in phase, and its amplitude is $C V_{0} \omega$.
The capacitor's reactance is $1/(C \omega)$, derived from the
ratio of voltage amplitude to current amplitude. Volume II
chapter~3 (complex numbers and phasors) introduced the phasor
representation that makes this manipulation algebraic; the
calculus underlies it.

For a capacitor charging through a resistor from a constant
voltage source, the voltage approaches the source value
exponentially: $v(t) = V_{\text{source}} (1 - e^{-t/(RC)})$. The
current is $i(t) = (V_{\text{source}}/R) e^{-t/(RC)}$. At $t = 0$
the capacitor draws maximum current; as $t \to \infty$ the
current decays to zero. The exponential time constant is $\tau =
RC$, with units of seconds; after one time constant the voltage
has reached $1 - 1/e \approx 0.632$ of the source value, and after
five time constants it has reached about $0.993$. The exponential
solution is recovered formally in chapter~7 of this volume,
which takes up first-order ordinary differential equations.

\subsection{Thermodynamics: heat capacity and Newton's law of cooling}

The heat capacity of a substance is the derivative of the heat
required to raise its temperature: $C_{p} = \dv{Q}{T}$ at constant
pressure. For most substances, $C_{p}$ is a slowly-varying
function of temperature; engineering heat-transfer calculations
usually treat it as a constant over a working temperature range.

Newton's law of cooling states that the rate of heat loss from a
body to its surroundings is proportional to the temperature
difference between the body and the surroundings:

\[
\dv{T}{t} = -k (T - T_{\infty}),
\]

where $k$ is a heat-transfer coefficient with units of
$\si{\per\second}$ and $T_{\infty}$ is the ambient temperature. The
solution is $T(t) = T_{\infty} + (T_{0} - T_{\infty}) e^{-k t}$,
exponential approach to ambient with time constant $1/k$. The
calculus enters through the time derivative on the left and through
the exponential solution; chapter~7 again handles the formal ODE
solution.

A worked check. A mug of coffee at $80\,\si{\degreeCelsius}$ is
left in a $20\,\si{\degreeCelsius}$ room and reaches $40\,\si
{\degreeCelsius}$ after twenty minutes. The cooling constant
satisfies $40 = 20 + 60 e^{-k \cdot 20}$, so $e^{-20 k} = 1/3$ and
$k = \ln 3 / 20 \approx 0.055\,\si{\per\minute}$. The half-life
of the cooling is $\ln 2 / k \approx 12.6\,\si{\minute}$. The
mug's temperature is described by a single derivative-driven
relationship; the same calculus describes a building's thermal
response, an exhausting battery, and a transient electronic
component reaching equilibrium.

\subsection{Biology: logistic growth}

A bacterial population in a closed flask grows according to the
logistic equation

\[
\dv{N}{t} = r N \left(1 - \frac{N}{K}\right),
\]

where $N(t)$ is the population, $r$ is the intrinsic growth rate,
and $K$ is the carrying capacity. The right side is the rate at
which the population changes; it is $r N$ at small $N$ (where the
factor $(1 - N/K)$ is approximately $1$) and approaches zero as
$N \to K$. The growth rate per individual, $\dv{N}{t} \big/ N = r
(1 - N/K)$, is maximal at small $N$ and decreases linearly with
population.

The peak rate of population change occurs where the second
derivative $\frac{\dd^{2} N}{\dd t^{2}}$ is zero. Differentiating
the logistic equation with respect to time and using the chain
rule:

\[
\frac{\dd^{2} N}{\dd t^{2}} = r \dv{N}{t} \left(1 - \frac{2 N}
{K}\right).
\]

The right side vanishes when $\dv{N}{t} = 0$ (at $N = 0$ or $N =
K$, both extrema of the population) or when $N = K/2$. At
$N = K/2$, the population's rate of change is at its maximum:
$\dv{N}{t} = r K/4$. A microbiologist who wants the fastest
growth, in absolute terms, harvests at $N = K/2$. The same
calculation, with different $r$ and $K$, governs viral spread,
animal-population management, and many engineered fermentation
processes.

\subsection{Transport: the gradient as flux density}

\begin{archetype}[Transport, gradient as flux density]
A second pattern enters with the derivative: the gradient of a
field drives a flux. Heat flows down a temperature gradient, mass
diffuses down a concentration gradient, charge flows down a voltage
gradient, fluid flows down a pressure gradient. In one spatial
dimension each is a constitutive law of the form
\[
J = -\kappa \dv{\phi}{x},
\]
where $\phi$ is the driving field, $J$ is the flux (a quantity per
unit area per unit time), and $\kappa$ is a material transport
coefficient. The minus sign encodes that the flux runs from high
field to low. The derivative is the engine of the law; the
transport coefficient converts the dimensioned gradient into a
dimensioned flux. The archetype recurs across Volumes~IV through
VII, in Fourier's law of heat conduction, Fick's law of diffusion,
Ohm's law in its local form, and Darcy's law for porous flow. The
multivariable generalisation replaces the one-dimensional
derivative by the gradient vector of chapter~8 and the scalar
$\kappa$ by a conductivity tensor.
\end{archetype}

\engterm{Fourier's law in a wall}. Heat conducts through a plane
wall of thickness $L$ whose faces are held at temperatures $T_{1}$
and $T_{2} < T_{1}$. In steady state the temperature falls linearly
across the wall, $T(x) = T_{1} - (T_{1} - T_{2}) x / L$, so the
gradient is constant, $\dv{T}{x} = -(T_{1} - T_{2})/L$. Fourier's
law gives the heat flux $q = -k \dv{T}{x} = k (T_{1} - T_{2})/L$,
with $k$ the thermal conductivity in $\si{\watt\per\meter\per
\kelvin}$. For a $0.2 \,\si{\meter}$ brick wall with $k = 0.7
\,\si{\watt\per\meter\per\kelvin}$ and a $20 \,\si{\kelvin}$ face-
to-face difference, the flux is $q = 0.7 \cdot 20 / 0.2 = 70
\,\si{\watt\per\meter\squared}$. The derivative of the temperature
profile is the quantity the law multiplies; a steeper gradient (a
thinner wall or a larger temperature difference) drives a
proportionally larger flux.

\engterm{Fick's law for diffusion}. A membrane of thickness $L$
separates two solutions of a solute at concentrations $c_{1}$ and
$c_{2}$. In steady state the concentration falls linearly across
the membrane, and Fick's first law gives the molar flux $J = -D
\dv{c}{x} = D (c_{1} - c_{2})/L$, with $D$ the diffusion
coefficient in $\si{\meter\squared\per\second}$. The structure is
identical to Fourier's law with concentration in place of
temperature and the diffusion coefficient in place of the thermal
conductivity. An engineer who has internalised one transport law
has internalised them all; the derivative-driven flux is the
shared skeleton.

\subsection{Civil engineering: the deflection of a beam}

The shape of a loaded beam is described by its deflection $y(x)$
along its length. The successive derivatives carry the structural
quantities the engineer designs against. The first derivative
$\dv{y}{x}$ is the slope of the deflected beam. The second
derivative $\frac{\dd^{2} y}{\dd x^{2}}$ is proportional to the
bending moment $M(x)$ through the curvature relation $M = E I \,
\frac{\dd^{2} y}{\dd x^{2}}$, where $E$ is the elastic modulus and
$I$ is the second moment of area of the cross-section. The third
derivative is proportional to the shear force $V(x) = \dv{M}{x} =
E I \, \frac{\dd^{3} y}{\dd x^{3}}$, and the fourth derivative is
proportional to the distributed load $w(x) = -\dv{V}{x} = -E I \,
\frac{\dd^{4} y}{\dd x^{4}}$. Four derivatives of one deflection
function generate the entire family of structural quantities:
slope, moment, shear, and load. The Euler-Bernoulli beam equation
$E I \, \frac{\dd^{4} y}{\dd x^{4}} = w(x)$ that organises this
family is solved in chapter~7 of this volume; the derivatives that
populate it are the subject of the present chapter. A reader who
designs a beam reads the maximum bending stress off the second
derivative and the maximum shear off the third, each located where
the corresponding derivative reaches its extreme along the span.

\subsection{End-to-end case study: linearising a thermistor sensor}

We close the worked examples with a full engineering problem driven
from start to finish by the derivative. The task is to design and
characterise a temperature-measurement circuit built around a
negative-temperature-coefficient (NTC) thermistor, a resistor whose
resistance falls steeply with temperature. The problem exercises
every tool of the chapter in one connected calculation: the
derivative as sensitivity, the local linearisation about an
operating point, the second derivative as a bound on the
linearisation's validity, and the propagation of sensitivities into
an error budget. The case is representative of how a working
engineer turns a nonlinear device characteristic into a usable
instrument.

\paragraph{The device characteristic.}
An NTC thermistor's resistance follows the Steinhart family of
models; the simplest adequate form is the beta model

\[
R(T) = R_{0} \exp\!\left[B\left(\frac{1}{T} - \frac{1}{T_{0}}
\right)\right],
\]

where $T$ is the absolute temperature in kelvin, $R_{0}$ is the
resistance at the reference temperature $T_{0}$, and $B$ is the
material constant in kelvin. We take a common part: $R_{0} =
10\,\si{\kilo\ohm}$ at $T_{0} = 298.15\,\si{\kelvin}$ ($25\,
\si{\degreeCelsius}$) and $B = 3950\,\si{\kelvin}$. The
characteristic is strongly nonlinear; \Cref{fig:vol02:ch05:operating-point}
shows its convex, falling shape.

\begin{figure}[ht]
\centering
\begin{tikzpicture}[x=1cm, y=1cm, font=\small]

% Thermistor resistance-temperature curve R(T) = R0 exp(B(1/T - 1/T0)),
% strongly convex and falling. Show the curve, an operating point, and
% the tangent (small-signal) line through it. Illustrative shape only.
% T axis (deg C) from 0 to 80 mapped to x in [0, 9]; x = 9*(T)/80.
% R axis from 0 to 30 kOhm mapped to y in [0, 5]; y = 5*R/30.
% Curve sampled at chosen (T, R) points to show convex falling shape.

\draw[->, thick] (-0.3, 0) -- (9.6, 0)
  node[right] {$T\;(^{\circ}\mathrm{C})$};
\draw[->, thick] (0, -0.3) -- (0, 5.6)
  node[above] {$R\;(\mathrm{k}\Omega)$};

% x ticks
\foreach \T in {0,20,40,60,80} {
  \pgfmathsetmacro\xx{9*\T/80}
  \draw (\xx, 0) -- (\xx, -0.12);
  \node[anchor=north, font=\scriptsize] at (\xx, -0.12) {\T};
}
% y ticks
\foreach \R in {0,10,20,30} {
  \pgfmathsetmacro\yy{5*\R/30}
  \draw (0, \yy) -- (-0.12, \yy);
  \node[anchor=east, font=\scriptsize] at (-0.14, \yy) {\R};
}

% Curve: representative NTC points (T degC, R kOhm)
% (0,28) (10,20) (20,14.4) (25,12.0) (30,10.2) (40,7.4) (50,5.4)
% (60,4.0) (70,3.0) (80,2.3). Plot scaled.
\draw[very thick, engblue]
  plot coordinates {
    ({9*0/80},  {5*28.0/30})
    ({9*10/80}, {5*20.0/30})
    ({9*20/80}, {5*14.4/30})
    ({9*25/80}, {5*12.0/30})
    ({9*30/80}, {5*10.2/30})
    ({9*40/80}, {5*7.4/30})
    ({9*50/80}, {5*5.4/30})
    ({9*60/80}, {5*4.0/30})
    ({9*70/80}, {5*3.0/30})
    ({9*80/80}, {5*2.3/30})
  };

% Operating point at T0 = 25 degC, R0 = 12.0 kOhm.
% Local slope dR/dT approx (10.2 - 14.4)/(30-20) = -0.42 kOhm/degC.
% Tangent line over T in [10, 45]: R = 12.0 - 0.42*(T - 25).
% At T=10: R = 12.0 - 0.42*(-15) = 18.3; at T=45: R = 12.0 - 0.42*20 = 3.6.
\draw[thick, engred, dashed]
  plot coordinates {
    ({9*10/80}, {5*18.3/30})
    ({9*45/80}, {5*3.6/30})
  };

\fill[engblack] ({9*25/80}, {5*12.0/30}) circle (1.8pt);
\node[engblack, anchor=south west, font=\scriptsize]
  at ({9*25/80 + 0.1}, {5*12.0/30 + 0.05})
  {operating point $(T_{0}, R_{0})$};

% small-signal box near operating point
\draw[thin, enggray, densely dotted]
  ({9*20/80}, 0) -- ({9*20/80}, {5*14.4/30});
\draw[thin, enggray, densely dotted]
  ({9*30/80}, 0) -- ({9*30/80}, {5*10.2/30});
\node[enggray, anchor=north, font=\tiny]
  at ({9*25/80}, -0.45) {small-signal band};

\node[engblue, anchor=east, font=\scriptsize] at (9.4, 4.7)
  {$R(T)$ (NTC)};
\node[engred, anchor=east, font=\scriptsize] at (9.4, 4.3)
  {tangent at $T_{0}$};

\end{tikzpicture}
\caption[Operating-point linearisation of a thermistor]{The
resistance-temperature curve of a negative-temperature-coefficient
thermistor (blue) falls steeply and convexly with temperature. A
measurement circuit operated in a narrow band around the operating
point $(T_{0}, R_{0})$ replaces the curve by its tangent (red dashed),
the small-signal model $\delta R \approx (\mathrm{d}R/\mathrm{d}T)_{T_{0}}
\,\delta T$. The slope at the operating point sets the sensor's local
sensitivity in $\mathrm{k}\Omega$ per degree; the curvature, the
second derivative, sets the width of the band over which the linear
model stays inside a stated tolerance. Reading the design off the
derivative at one point is the entire content of small-signal sensor
calibration.}
\label{fig:vol02:ch05:operating-point}
\end{figure}


\paragraph{Step 1: the sensitivity at the operating point.}
The first job is the local sensitivity $\dv{R}{T}$, the quantity
that converts a temperature change into a resistance change the
circuit can read. Differentiating the beta model with the chain
rule, the inner derivative of $B/T$ is $-B/T^{2}$, so

\[
\dv{R}{T} = R(T) \cdot \left(-\frac{B}{T^{2}}\right) =
-\frac{B}{T^{2}} R(T).
\]

At the operating point $T = T_{0} = 298.15\,\si{\kelvin}$, where
$R = R_{0} = 10\,\si{\kilo\ohm}$, the sensitivity is

\[
\left.\dv{R}{T}\right|_{T_{0}} = -\frac{3950}{(298.15)^{2}} \cdot
10{,}000\,\si{\ohm} = -\frac{3950}{88{,}893} \cdot 10{,}000\,
\si{\ohm} \approx -444\,\si{\ohm\per\kelvin}.
\]

The sensor loses about $444\,\si{\ohm}$ of resistance per kelvin of
warming near room temperature. The minus sign is the defining
feature of an NTC device. The relative sensitivity is cleaner: $(1/
R)\dv{R}{T} = -B/T^{2} = -0.0444\,\si{\per\kelvin}$, so the
resistance changes by about $4.4\,\%$ per kelvin near $25\,
\si{\degreeCelsius}$, a large and easily measured fractional
change. This is why thermistors are the sensor of choice for narrow-
range, high-resolution temperature measurement.

\paragraph{Step 2: the small-signal linear model.}
A measurement circuit operated in a narrow band around $T_{0}$
replaces the exponential by its tangent line, the local
linearisation of section~5.4:

\[
R(T) \approx R_{0} + \left.\dv{R}{T}\right|_{T_{0}} (T - T_{0}) =
10{,}000 - 444\,(T - T_{0})\quad[\si{\ohm}],
\]

with $T - T_{0}$ in kelvin. Read backward, the model converts a
measured resistance into an inferred temperature: $T \approx T_{0}
+ (R - R_{0})/(\dv{R}{T}) = T_{0} - (R - 10{,}000)/444$. A circuit
that measures resistance to $\pm 1\,\si{\ohm}$ therefore resolves
temperature to about $1/444 \approx 2.3\,\si{\milli\kelvin}$ near
the operating point, a resolution the bare nonlinear curve does not
advertise but the derivative makes explicit.

\paragraph{Step 3: the second derivative bounds the linear range.}
The linear model is good only over a band where the curvature is
small enough that the quadratic Taylor remainder stays below
tolerance. The remainder of the first-order approximation is
$\tfrac{1}{2} R''(\xi)(T - T_{0})^{2}$, so we need the second
derivative. Differentiating $\dv{R}{T} = -(B/T^{2}) R(T)$ by the
product rule,

\[
\frac{\dd^{2} R}{\dd T^{2}} = \frac{2B}{T^{3}} R(T) +
\left(-\frac{B}{T^{2}}\right)\!\left(-\frac{B}{T^{2}} R(T)\right)
= R(T)\,\frac{B}{T^{2}}\left(\frac{B}{T^{2}} + \frac{2}{T}\right).
\]

At $T_{0}$ this evaluates to $R''(T_{0}) = 10{,}000 \cdot 0.0444
\cdot (0.0444 + 0.00671) \approx 10{,}000 \cdot 0.0444 \cdot
0.0511 \approx 22.7\,\si{\ohm\per\kelvin\squared}$. The linear
model's error at a displacement $\Delta T$ from the operating point
is about $\tfrac{1}{2} R''(T_{0}) \Delta T^{2} = 11.3\,\Delta T^{2}\,
\si{\ohm}$. Setting this equal to the $1\,\si{\ohm}$ resistance
resolution gives $\Delta T = \sqrt{1/11.3} \approx 0.30\,\si
{\kelvin}$: the linear model is good to within the measurement
resolution only over a band of about $\pm 0.3\,\si{\kelvin}$ around
the operating point. To cover a wider range at the same accuracy,
the engineer either re-linearises at several operating points (a
piecewise-linear lookup table, each segment with its own slope) or
keeps the full nonlinear model and inverts it numerically. The
second derivative made the design decision quantitative: it set the
width of each linear segment.

\paragraph{Step 4: the error budget by sensitivity propagation.}
The measured temperature inherits uncertainty from three sources:
the resistance measurement, the calibration of $R_{0}$, and the
material constant $B$. Treating the inferred temperature as a
function $T(R, R_{0}, B)$ and applying the first-order propagation
of section~5.5, each source contributes its uncertainty weighted by
the corresponding sensitivity. The resistance-measurement
contribution is $(\partial T/\partial R)\,u(R) = u(R)/|\dv{R}{T}|
= u(R)/444$. For $u(R) = 5\,\si{\ohm}$ (a typical mid-grade
ohmmeter at this range), this is $5/444 \approx 11\,\si{\milli
\kelvin}$. The material-constant contribution is the larger worry:
a $1\,\%$ uncertainty in $B$ (manufacturers often quote $B$ to
$\pm 1\,\%$) shifts the inferred temperature by an amount that
grows with the displacement from the calibration point. At the
operating point itself, where $T = T_{0}$, the $B$ term drops out
of the model, so the $B$ uncertainty contributes nothing at $T_{0}$
and grows linearly away from it; at $\Delta T = 10\,\si{\kelvin}$
from calibration the $B$ uncertainty contributes roughly
$|\Delta T| \cdot (T/T_{0}) \cdot u(B)/B \cdot$ a factor of order
one, on the order of $0.1\,\si{\kelvin}$, an order of magnitude
above the resistance-measurement floor. The budget tells the
engineer where to spend: near the calibration point the resistance
measurement dominates and a better ohmmeter helps; far from it the
material-constant uncertainty dominates and a multi-point
calibration that pins $B$ helps. The derivative-driven budget
ranks the interventions before any hardware is bought.

\paragraph{Interpretation.}
The thermistor problem is a microcosm of instrument design. The
first derivative set the sensitivity and hence the resolution. The
linear model turned the nonlinear device into a circuit the
engineer can read in one subtraction. The second derivative bounded
the range over which that linear reading is trustworthy and set the
segment width of a lookup table. The sensitivity propagation built
the error budget and ranked the uncertainty sources so the design
effort goes where it changes the answer. Every step was a
derivative of the device characteristic, evaluated at the operating
point, read as a sensitivity. A reader who internalises this
sequence has the working pattern for linearising any nonlinear
sensor, actuator, or component about its operating point, which is
the daily work of analogue and instrument engineering and the
content the small-signal models of Volume~IX assume.

\subsection{Second end-to-end case study: solving a cubic equation of
state with a globalised Newton iteration}\pathtag{standard}

The thermistor case ran the derivative forward, from a known model to a
linearisation and an error budget. We close with a case that runs the
derivative the other way, inside a loop, to solve an equation that has
no closed-form root. The task is to find the molar volume of a real gas
at a stated temperature and pressure from a cubic equation of state.
The calculation is the daily work of process and thermodynamic
engineering, and it exposes every pathology of Newton's iteration that
section~5.4 catalogued: a small derivative that throws a full step far
from the root, multiple roots of which only one is physical, and a
near-critical regime where two roots coalesce and the derivative
collapses. The defensive apparatus that survives these pathologies,
step damping and parameter continuation, is built entirely from the
linear approximation of this chapter.

\paragraph{The equation and its residual.}
The Peng-Robinson equation of state relates pressure $p$, absolute
temperature $T$, and molar volume $v$ through
\[
p = \frac{R T}{v - b} - \frac{a\,\alpha(T)}{v^{2} + 2 b v - b^{2}},
\]
with $R$ the gas constant, $a$ and $b$ substance constants fixed by the
critical point, and $\alpha(T)$ a dimensionless temperature function.
Multiplying out and grouping in the dimensionless compressibility
$Z = p v /(R T)$ converts the equation into a cubic in $Z$,
\[
\phi(Z) = Z^{3} - (1 - B) Z^{2} + \left(A - 2 B - 3 B^{2}\right) Z -
\left(A B - B^{2} - B^{3}\right) = 0,
\]
where $A = a\,\alpha(T)\, p /(R T)^{2}$ and $B = b\, p /(R T)$ are the
two dimensionless groups that carry the operating condition. The
engineer wants the root $Z$, then recovers the molar volume as
$v = Z R T / p$. The cubic has either one real root (a single phase) or
three real roots (the largest is the vapour-phase compressibility, the
smallest is the liquid-phase compressibility, and the middle root is
spurious, an artefact of the cubic with no physical meaning). The
derivative of the residual,
\[
\phi'(Z) = 3 Z^{2} - 2 (1 - B) Z + \left(A - 2 B - 3 B^{2}\right),
\]
is the slope Newton's iteration divides by at every step.

\paragraph{Step 1: a working operating point.}
We take carbon dioxide near its critical point, where the cubic is at
its most awkward. With $A = 0.45$ and $B = 0.078$ (representative of
$\mathrm{CO_{2}}$ at a reduced temperature just above critical and a
moderate pressure), the residual $\phi$ has three real roots close
enough together that a careless solver lands on the wrong one.
\Cref{fig:vol02:ch05:eos-newton} plots $\phi$ for an illustrative
near-critical condition: the liquid root, the spurious middle root, and
the vapour root sit within a narrow band of $Z$.

\begin{figure}[ht]
\centering
\begin{tikzpicture}[x=1cm, y=1cm, font=\small]

% Residual phi(Z) = Z^3 - (1 - B) Z^2 + (A - 2B - 3B^2) Z - (AB - B^2 - B^3)
% for a near-critical case with three real roots: liquid, spurious, vapour.
% We use illustrative coefficients giving roots near Z = 0.07, 0.30, 0.62.
% Display: plotx = (Z - 0)*12 ; ploty = phi(Z)*60 (vertical units cm).

\def\sx{12.0}     % horizontal scale (Z in [0,0.75])
\def\sy{55.0}     % vertical scale for phi

% phi(Z) = (Z-0.07)(Z-0.30)(Z-0.62), expanded only for plotting.
% Axes
\draw[->, thick] (0, 0) -- ({0.78*\sx}, 0) node[right] {$Z$};
\draw[->, thick] (0, -1.7) -- (0, 2.1) node[above] {$\phi(Z)$};

% Curve
\draw[thick, engblue] plot[domain=0.02:0.74, samples=120]
  ({\x*\sx}, {(\x-0.07)*(\x-0.30)*(\x-0.62)*\sy});

% Roots
\fill[engred] ({0.07*\sx}, 0) circle (1.8pt);
\node[anchor=north, font=\scriptsize] at ({0.07*\sx}, -0.08) {$Z_{\ell}$};
\fill[enggray] ({0.30*\sx}, 0) circle (1.6pt);
\node[anchor=south, font=\scriptsize] at ({0.30*\sx}, 0.06) {$Z_{\mathrm{spur}}$};
\fill[engred] ({0.62*\sx}, 0) circle (1.8pt);
\node[anchor=north, font=\scriptsize] at ({0.62*\sx}, -0.08) {$Z_{v}$};

% A naive Newton step from a bad guess x0 = 0.45 (near the local max of phi
% between the spurious and vapour roots) overshoots far to the left.
\def\xa{0.45}
\pgfmathsetmacro{\pa}{(\xa-0.07)*(\xa-0.30)*(\xa-0.62)*\sy}
% slope of plotted curve at xa (numerically, central difference on display)
\pgfmathsetmacro{\hh}{0.001}
\pgfmathsetmacro{\pap}{((\xa+\hh-0.07)*(\xa+\hh-0.30)*(\xa+\hh-0.62)
  - (\xa-\hh-0.07)*(\xa-\hh-0.30)*(\xa-\hh-0.62))*\sy/(2*\hh*\sx)}
\fill[engamber] ({\xa*\sx}, \pa) circle (1.8pt);
\node[engamber, anchor=south west, font=\scriptsize] at ({\xa*\sx + 0.05}, {\pa+0.05}) {$Z^{(0)}$};
% tangent: zero crossing at plotx_zero
\pgfmathsetmacro{\xzero}{\xa*\sx - \pa/\pap}
\draw[thick, engamber] ({\xzero}, 0) -- ({\xa*\sx + 0.9}, {\pa + \pap*0.9});
\fill[engblack] ({\xzero}, 0) circle (1.4pt);
\node[engamber, anchor=south, font=\scriptsize] at ({\xzero}, 0.06) {full step};

\node[engblue, anchor=north east, font=\scriptsize] at ({0.74*\sx}, {(0.74-0.07)*(0.74-0.30)*(0.74-0.62)*\sy - 0.1}) {$\phi(Z)$};

\end{tikzpicture}
\caption[Cubic equation-of-state residual and the Newton trap]{The
cubic-equation-of-state residual $\phi(Z)$ for a near-critical
condition with three real roots: the liquid root $Z_{\ell}$, a
spurious middle root $Z_{\mathrm{spur}}$ with no physical meaning, and
the vapour root $Z_{v}$. A Newton step taken from a guess $Z^{(0)}$
near the interior local maximum, where $\phi'$ is small, overshoots far
past the nearest root, the divergent mode of section~5.4. The
globalised solver of the worked case damps the step and converges to
the physically correct root for the requested phase.}
\label{fig:vol02:ch05:eos-newton}
\end{figure}


\paragraph{Step 2: the naive Newton iteration fails.}
A reader who applies the bare iteration $Z^{(k+1)} = Z^{(k)} -
\phi(Z^{(k)})/\phi'(Z^{(k)})$ from an arbitrary guess discovers the
failure mode of section~5.4 directly. Between the spurious root and the
vapour root the residual has an interior local maximum, where
$\phi'(Z) = 0$. A guess near that maximum has a tiny denominator, so
the step $\phi/\phi'$ is enormous and the iterate flies far past every
root, often to a negative $Z$ where the molar volume is unphysical. The
behaviour is the same divergent overshoot the chapter showed for
$\arctan x$, now embedded in a problem the engineer cannot avoid:
near-critical conditions are exactly where the equation of state is
needed and exactly where its residual is flattest.

\paragraph{Step 3: globalisation by step damping.}
The fix keeps Newton's direction but shortens the step. The Newton
direction $d^{(k)} = -\phi(Z^{(k)})/\phi'(Z^{(k)})$ is the displacement
that the linear model predicts will reach the root. When the linear
model is trustworthy the full step is taken; when it is not, a
\engterm{damped step} $Z^{(k+1)} = Z^{(k)} + \lambda\, d^{(k)}$ with a
fraction $\lambda \in (0, 1]$ is taken instead. The fraction $\lambda$
is chosen so that the magnitude of the residual actually decreases, the
\engterm{Armijo sufficient-decrease} test,
\[
|\phi(Z^{(k)} + \lambda d^{(k)})| \leq (1 - c\,\lambda)\,|\phi(Z^{(k)})|,
\]
with a small constant $c \approx 10^{-4}$. The engineer starts with
$\lambda = 1$ (the full Newton step) and halves $\lambda$ until the test
passes \cite{text:v2ch05-nocedal-wright-numerical-optimisation-2006}.
Near a root the full step always passes on the first try, so the
quadratic convergence of pure Newton is recovered once the iterate is
close; far from a root the backtracking shortens the step enough to make
guaranteed progress. The damped iteration is the working default for
every nonlinear solve in engineering practice, and its entire content is
the chapter's linear approximation, used as a direction rather than as a
final answer, with a one-dimensional search guarding the step length.

A worked damped step. From the trapping guess $Z^{(0)} = 0.45$ of
\Cref{fig:vol02:ch05:eos-newton}, the full Newton step overshoots and
the residual grows. The Armijo loop halves the step: $\lambda = 1$
fails, $\lambda = \tfrac{1}{2}$ fails, $\lambda = \tfrac{1}{4}$
reduces the residual and is accepted. The damped iterate lands inside
the basin of the vapour root, and from there the full step passes every
subsequent test, so the last three iterations converge quadratically to
$Z_{v}$ to machine precision. The molar volume follows as
$v = Z_{v} R T / p$. The same machinery, started from a small guess such
as $Z^{(0)} = B + 10^{-3}$ just above the liquid pole, converges to the
liquid root $Z_{\ell}$, which is how a flash calculation obtains both
phase volumes from one residual.

\paragraph{Step 4: selecting the physical root.}
The cubic returns up to three roots; the engineer keeps the one the
physics demands. The two outer roots are physical (vapour and liquid);
the middle root is discarded because the isothermal compressibility it
implies is negative, a thermodynamic impossibility, which the sign of
$\phi'$ at the root detects directly: at the spurious root $\phi'$ has
the opposite sign to that at the two physical roots. When the question
asks for the stable phase, the engineer compares the Gibbs energy of the
liquid and vapour roots and keeps the lower; the derivative reappears
there too, because the Gibbs energy difference is an integral of $v\,
\dd p$ whose integrand is read off the same equation of state. The root
selection is not a numerical detail. A solver that returns the spurious
root, or the vapour root when the liquid is stable, produces a molar
volume wrong by an order of magnitude, and a downstream sizing
calculation built on it sizes a vessel for the wrong phase.

\paragraph{Step 5: continuation through the critical point.}
The hardest condition is the critical point itself, where all three
roots coalesce into one triple root. There $\phi(Z) = 0$,
$\phi'(Z) = 0$, and $\phi''(Z) = 0$ simultaneously, and Newton's
quadratic convergence collapses to the linear crawl of a multiple root
that the section~5.4 mastery box described. The escape is
\engterm{continuation} (also called homotopy): solve a nearby easy
condition first, then track the root as the operating condition is moved
in small steps toward the hard target. The engineer embeds the problem
in a one-parameter family $\phi(Z; s) = 0$ with $s$ a continuation
parameter running from an easy value $s = 0$ (an ideal-gas condition far
from critical, where the single root $Z \approx 1$ is found in two
Newton steps) to the target $s = 1$ (the near-critical condition). At
each $s$, the converged root from the previous $s$ is the starting guess
for the next, and the implicit function theorem guarantees that the root
moves continuously with $s$ as long as $\phi'(Z; s) \neq 0$ along the
path. The continuation step size is shortened automatically wherever the
root path steepens, exactly the same step-control logic as the Armijo
backtracking, applied in the parameter direction rather than the
solution direction. Continuation turns an intractable cold start into a
sequence of easy warm starts, and the tangent that defines each warm
start is the derivative $\partial Z/\partial s = -(\partial \phi/
\partial s)/(\partial \phi/\partial Z)$, the implicit-differentiation
formula of section~5.2 used to predict where the root goes next.

\paragraph{Interpretation.}
The equation-of-state solve is the linear approximation used as an
engine rather than as an answer. Newton's step is the tangent-line
prediction of the root. The Armijo backtracking is a one-dimensional
search that trusts the tangent only as far as the residual rewards it.
The root selection reads the sign of the derivative to discard the
unphysical root. The continuation tracks the root through the critical
point on the implicit-function tangent, shortening its parameter step
where the derivative warns of trouble. Every safeguard is the chapter's
derivative, deployed defensively against the chapter's own failure
modes. The reader who internalises this sequence has the working pattern
for solving any smooth nonlinear equation that resists a closed-form
root, which is most of them, and the discipline that keeps the solve
from landing on the wrong answer in the regime where the answer matters
most.

\section{Project}\pathtag{core}

\begin{project}[Numerically differentiate a noisy signal; defend a smoothing strategy]

The reader takes a real or synthesised noisy time series of their
choice and computes its derivative numerically by at least three
methods. Each method has free parameters (step size, filter
length, polynomial order); the reader chooses each parameter,
documents the rationale, and reports the resulting derivative
estimate with its uncertainty.

\textbf{Track}. Simulation with analysis. \textbf{Hazard class}:
none. \textbf{Tools}. Paper, pencil, programming environment of
the reader's choice (Python, MATLAB, Julia, or equivalent).

\textbf{Source data}. The reader chooses one of:

\begin{itemize}[leftmargin=*]
  \item A time series from a Volume I project (the reader's
    pendulum swing, building-height shadow record, refrigerator-
    temperature log, or other dataset).
  \item A publicly available sensor record (smartphone
    accelerometer, weather-station temperature, public-utility
    power consumption, financial price series, biological tracking
    data).
  \item A synthesised signal $y(t) = f_{\text{true}}(t) +
    \varepsilon(t)$ with the reader's choice of $f_{\text{true}}$
    (a sinusoid, a polynomial, an exponential decay, a step plus
    ramp) and $\varepsilon(t) \sim \mathcal{N}(0, \sigma^{2})$ with
    a stated $\sigma$. For the synthesised case, the reader has the
    ground-truth derivative for comparison.
\end{itemize}

The data should contain at least $200$ samples; more is better.
The reader should record the sampling rate, the units of the
quantity being differentiated, and any pre-processing (mean-
removal, detrending) applied before the derivative is computed.

\textbf{Method 1: finite differences}. Compute the derivative by
the forward difference, the backward difference, and the central
difference at each interior sample. Sweep the effective step size
by sub-sampling the data: at sample spacing $\Delta t$, $2\Delta
t$, $4\Delta t$, $8\Delta t$. For each step size, compute the
mean and standard deviation of the derivative estimate over a
plateau region of the signal. Plot the derivative estimate and
its scatter as a function of step size on log-log axes. Identify
the optimal step size from the plot.

\textbf{Method 2: smoothing then differentiation}. Apply a moving-
average filter or a Savitzky-Golay filter of polynomial order $p$
and window width $w$ to the data, then differentiate the smoothed
signal by central differences. Sweep the window width and the
polynomial order. Compare the derivative estimates; identify a
window width and polynomial order that give the smoothest
derivative without removing genuine signal features. (The
Savitzky-Golay filter has a closed-form derivative that the
reader can use directly without re-differentiating after
smoothing.)

\textbf{Method 3: parametric fit then differentiation}. Fit a
parametric model to the data (a polynomial, a sum of sinusoids,
a logistic or exponential family, the reader's choice based on
the data's structure), compute the model's parameters by least-
squares, and differentiate the model analytically. Report the
parameter uncertainties and propagate them into the derivative
estimate.

\textbf{Comparison}. The three methods produce three derivative
estimates. The reader produces a table with one column per
method, one row per a chosen set of evaluation points, and one
additional column per method giving the estimated derivative
uncertainty at each point. For synthesised data, an extra column
gives the ground-truth derivative and the absolute error of each
method's estimate.

\textbf{Deliverables}. (1) The data record (with units and
sampling rate). (2) The derivative-estimate table. (3) Three
plots: derivative estimate versus step size for Method 1; effect
of filter parameters for Method 2; parameter uncertainties for
Method 3. (4) A 1500-word reflection on the bias-variance
tradeoff in the project's setting. The reflection addresses:
which method gave the smoothest defensible derivative, where the
methods agreed and where they disagreed, what the bias-variance
tradeoff looks like in this signal, and what the reader's
confidence in the derivative would change for a real engineering
deployment.

\textbf{Time}. Two to four hours of coding and analysis, plus two
to three hours of writing. Total approximately five to seven
hours, distributed across two or three days.

\textbf{Phases.} The work decomposes into five phases. (1)
Acquire and document the data: read the source, record sampling
rate and units, and plot the raw signal. Confirm by inspection
that the signal is not pathological (no constant zero windows, no
unexplained gaps, no quantisation visible at the working scale).
(2) Compute Method 1 (finite differences) at the native sample
rate. Observe how the derivative looks; record the visible noise
amplitude relative to any underlying signal feature. (3) Sweep
the effective step size by sub-sampling; produce the step-size
plot. Identify the truncation regime and the round-off (or
noise) regime. (4) Apply Method 2 (smoothing) and Method 3
(parametric fit). For each, record the parameter choice and the
rationale. (5) Compose the comparison table and the reflection.

\textbf{Success criteria.} The project is acceptable when each of
the following is true. (1) The raw signal, the three method
outputs, and the three diagnostic plots are present and labelled
with units and time. (2) Each method's parameters (step size,
filter window, polynomial order) are stated and the choice is
defended in one or two sentences with reference to the
bias-variance tradeoff. (3) The reflection identifies at least
one disagreement among the three methods and proposes a
defensible reading of the disagreement (the methods differ
because the underlying signal has structure on a timescale that
some methods are sensitive to and others smooth away, for
example). (4) For synthesised data, the comparison table includes
the absolute error of each method against the ground-truth
derivative; for measured data, the table reports the
inter-method spread as the working uncertainty floor of the
derivative estimate.

\textbf{Common failure modes.} (i) Picking the smallest available
step size (the native sample interval) and reporting the
finite-difference derivative without subsampling. The reader
loses the bias-variance plot and cannot defend the step size.
(ii) Choosing a Savitzky-Golay window that is much wider than
the signal's feature scale; the derivative becomes a smooth
nothing and disagrees catastrophically with the parametric fit.
(iii) Fitting a parametric model whose structure does not match
the signal (polynomial fit to oscillatory data, sinusoidal fit
to monotone data), then differentiating; the model's bad fit
shows up in residuals that the reader must inspect before
trusting the derivative. (iv) Reporting only one method's
derivative because "the methods all agreed" without computing
the disagreement. The point of the comparison is the size of the
inter-method spread; without that, the project has not done the
chapter's work. (v) Treating the ground-truth derivative as an
expected answer in the synthesised case rather than as a
verification target. Different methods can produce different
defensible derivatives even when the ground truth is known; the
reflection should report what each method is sensitive to, not
which method "got the right answer".

\end{project}

\section{Exercises}

The exercises train calculation, derivation, estimation,
simulation, design, diagnosis, failure analysis, and judgment.
Calculation and derivation problems have full solutions; open-
ended problems have rubrics. Exercises are tagged with their
reader-path tier.

\subsection*{Calculation}\pathtag{core}

\begin{exercise}
Compute the derivative of $f(x) = 3 x^{4} - 2 x^{3} + x - 7$.
\end{exercise}

\paragraph{Solution.}
Linearity lets the derivative be taken term by term, and the power
rule $\dv{}{x} x^{n} = n x^{n-1}$ handles each term:
$\dv{}{x}(3 x^{4}) = 12 x^{3}$, $\dv{}{x}(-2 x^{3}) = -6 x^{2}$,
$\dv{}{x}(x) = 1$, and $\dv{}{x}(-7) = 0$, since the derivative of
a constant is zero. Assembling,
\[
f'(x) = 12 x^{3} - 6 x^{2} + 1.
\]
A check at $x = 0$ reads $f'(0) = 1$, which must equal the
coefficient of the linear term of $f$ (the derivative at the origin
of a polynomial is its linear coefficient); it does. A second check
notes that the leading behaviour of $f'$ for large $x$ is $12 x^{3}$,
one degree below the $3 x^{4}$ that dominates $f$, as the power rule
requires.

\begin{exercise}
Compute the derivative of $f(x) = x^{2} \sin x$.
\end{exercise}

\paragraph{Solution.}
The expression is a product, so the product rule applies with
$u = x^{2}$ and $v = \sin x$. Then $u' = 2x$ (power rule) and
$v' = \cos x$ (elementary table), giving
\[
f'(x) = u' v + u v' = 2 x \sin x + x^{2} \cos x.
\]
At $x = \pi$, $\sin\pi = 0$ and $\cos\pi = -1$, so $f'(\pi) =
0 + \pi^{2}(-1) = -\pi^{2} \approx -9.87$. The result is negative
there, consistent with the graph of $x^{2}\sin x$ descending as it
crosses $x = \pi$. Near the origin, $\sin x \approx x$, so
$f(x) \approx x^{3}$ and $f'(x) \approx 3 x^{2}$; expanding the
exact derivative for small $x$ gives $2x \cdot x + x^{2} \cdot 1 =
3 x^{2}$, matching.

\begin{exercise}
Compute the derivative of $f(x) = e^{-x^{2}/2}$.
\end{exercise}

\paragraph{Solution.}
The function is a composition $f = e^{u}$ with $u(x) = -x^{2}/2$, so
the chain rule applies. The outer derivative is $\dv{}{u} e^{u} =
e^{u}$, and the inner derivative is $u'(x) = -x$. Multiplying,
\[
f'(x) = e^{-x^{2}/2} \cdot (-x) = -x \, e^{-x^{2}/2}.
\]
The derivative vanishes at $x = 0$, where the Gaussian peaks; it is
negative for $x > 0$ (the bell descending to the right) and positive
for $x < 0$ (ascending to the peak). The factor $-x$ in front of the
ever-positive exponential makes $f'$ an odd function, mirroring the
even symmetry of $f$. The second derivative, by the product rule on
$-x e^{-x^{2}/2}$, is $(x^{2} - 1) e^{-x^{2}/2}$, which vanishes at
$x = \pm 1$: the inflection points of the bell, one standard
deviation from the centre when the exponent is read as a unit-
variance Gaussian.

\begin{exercise}
Compute the derivative of $f(x) = (x^{2} + 1)/(x^{2} - 1)$ for
$x \neq \pm 1$.
\end{exercise}

\paragraph{Solution sketch.}
Quotient rule with $g(x) = x^{2} + 1$, $h(x) = x^{2} - 1$,
$g'(x) = h'(x) = 2 x$. Then
\[
f'(x) = \frac{2 x (x^{2} - 1) - (x^{2} + 1) \cdot 2 x}
{(x^{2} - 1)^{2}} = \frac{2 x ((x^{2} - 1) - (x^{2} + 1))}
{(x^{2} - 1)^{2}} = \frac{-4 x}{(x^{2} - 1)^{2}}.
\]
The derivative vanishes at $x = 0$ and diverges at $x = \pm 1$,
where $f$ itself has vertical asymptotes.

\begin{exercise}
Compute $\dv{y}{x}$ at $(x, y) = (1, 2)$ for the implicit
relation $x^{2} y + x y^{3} = 12$.
\end{exercise}

\paragraph{Solution sketch.}
Differentiate both sides with respect to $x$, treating $y$ as a
function of $x$ and applying the product rule and chain rule:
\[
2 x y + x^{2} \dv{y}{x} + y^{3} + 3 x y^{2} \dv{y}{x} = 0.
\]
Substitute $(x, y) = (1, 2)$: $2 \cdot 1 \cdot 2 + 1 \cdot
\dv{y}{x} + 8 + 3 \cdot 1 \cdot 4 \cdot \dv{y}{x} = 4 + \dv{y}{x}
+ 8 + 12 \dv{y}{x} = 12 + 13 \dv{y}{x} = 0$. Solving,
$\dv{y}{x}\big|_{(1,2)} = -12/13$. A consistency check: the
relation gives $12$ at the point, confirming the point is on the
curve.

\begin{exercise}
Compute the second derivative of $f(t) = t e^{-\alpha t}$ with
respect to $t$, where $\alpha > 0$ is a constant.
\end{exercise}

\paragraph{Solution sketch.}
Product rule once: $f'(t) = e^{-\alpha t} + t \cdot (-\alpha)
e^{-\alpha t} = (1 - \alpha t) e^{-\alpha t}$. Product rule
again: $f''(t) = (-\alpha) e^{-\alpha t} + (1 - \alpha t) \cdot
(-\alpha) e^{-\alpha t} = -\alpha e^{-\alpha t} (1 + 1 - \alpha t)
= -\alpha (2 - \alpha t) e^{-\alpha t} = \alpha (\alpha t - 2)
e^{-\alpha t}$. The second derivative is zero at $t = 2/\alpha$,
which is the inflection point of $f$; the first derivative is
zero at $t = 1/\alpha$, the location of the maximum.

\subsection*{Derivation}\pathtag{standard}

\begin{exercise}
Derive the derivative of $f(x) = \cos x$ from the limit
definition, using the trigonometric identity $\cos(x + h) =
\cos x \cos h - \sin x \sin h$ and the limits
$\lim_{h \to 0} \sin h / h = 1$, $\lim_{h \to 0} (1 - \cos h)/h
= 0$.
\end{exercise}

\paragraph{Solution sketch.}
The difference quotient is
\[
\frac{\cos(x + h) - \cos x}{h} = \frac{\cos x \cos h - \sin x
\sin h - \cos x}{h} = \cos x \cdot \frac{\cos h - 1}{h} -
\sin x \cdot \frac{\sin h}{h}.
\]
Take the limit as $h \to 0$. The first factor approaches
$\cos x \cdot 0 = 0$; the second factor approaches $\sin x \cdot
1 = \sin x$, so the difference quotient approaches $0 - \sin x =
-\sin x$. Thus $\dv{}{x} \cos x = -\sin x$.

\begin{exercise}
Derive the chain rule for $f(g(x))$ from the limit definition,
assuming $g$ is differentiable at $x$ and $f$ is differentiable
at $g(x)$. State the assumptions explicitly and identify where
the heuristic argument fails when $g$ is locally constant.
\end{exercise}

\paragraph{Solution sketch.}
Let $u = g(x)$ and $\Delta u = g(x + \Delta x) - g(x)$. Write the
difference quotient of the composite as
\[
\frac{f(g(x + \Delta x)) - f(g(x))}{\Delta x} =
\frac{f(u + \Delta u) - f(u)}{\Delta u} \cdot
\frac{\Delta u}{\Delta x},
\]
provided $\Delta u \neq 0$. As $\Delta x \to 0$, $\Delta u \to 0$
(continuity of $g$), the first factor approaches $f'(u) = f'(g(x))$,
and the second factor approaches $g'(x)$. The product approaches
$f'(g(x)) g'(x)$, the chain-rule formula. The heuristic fails when
$\Delta u = 0$ for some $\Delta x \neq 0$, which can happen if $g$
is locally constant. The standard fix replaces $f(u + \Delta u) -
f(u)$ by an expression $\phi(\Delta u) \Delta u$ where $\phi$ is
continuous at $0$ with $\phi(0) = f'(u)$; the product
$\phi(\Delta u) \Delta u / \Delta x$ then makes sense for all
$\Delta u$, including zero, and the limit is unchanged.

\begin{exercise}
Derive the second-derivative test for a local maximum from
Taylor's theorem to second order. State the assumption that
fails when $f''(c) = 0$.
\end{exercise}

\paragraph{Solution sketch.}
Suppose $f'(c) = 0$ and $f''$ is continuous near $c$. Taylor's
theorem to second order gives
$f(c + h) = f(c) + f'(c) h + \tfrac{1}{2} f''(\xi) h^{2} =
f(c) + \tfrac{1}{2} f''(\xi) h^{2}$ for some $\xi$ between $c$
and $c + h$. If $f''(c) < 0$, then by continuity $f''(\xi) < 0$
in a neighbourhood of $c$, so $f(c + h) - f(c) =
\tfrac{1}{2} f''(\xi) h^{2} < 0$ for all small $h \neq 0$; $c$
is a strict local maximum. The argument for $f''(c) > 0$ (local
minimum) is symmetric. When $f''(c) = 0$, the leading-order term
vanishes and the sign of $f(c + h) - f(c)$ is determined by the
next nonzero derivative; the theorem above is silent.

\begin{exercise}
Derive the truncation error of the central difference formula
$D_{0}(x; h) = (f(x + h) - f(x - h))/(2 h)$ from Taylor's theorem,
giving the leading-order term in $h$ and the conditions under
which the next term is negligible.
\end{exercise}

\paragraph{Solution sketch.}
Expand $f(x \pm h) = f(x) \pm f'(x) h + \tfrac{1}{2} f''(x) h^{2}
\pm \tfrac{1}{6} f'''(x) h^{3} + \tfrac{1}{24} f''''(\xi^{\pm})
h^{4}$. Subtract: $f(x + h) - f(x - h) = 2 f'(x) h +
\tfrac{1}{3} f'''(x) h^{3} + \tfrac{1}{12} (f''''(\xi^{+}) +
f''''(\xi^{-})) h^{4}/2$, with the even-order terms cancelling.
Divide by $2 h$ to obtain $D_{0}(x; h) = f'(x) + \tfrac{1}{6}
f'''(x) h^{2} + \oom{h^{4}}$. The leading-order truncation error
is $\tfrac{1}{6} f'''(x) h^{2}$, of order $h^{2}$. The next term
is negligible relative to the leading one when $h$ is small
enough that $|f''''| h^{4}$ is much less than $|f'''| h^{2}$,
that is, when $h^{2} \ll |f'''| / |f''''|$. For smooth functions
this holds for any reasonably small $h$.

\begin{exercise}
Derive the optimal step size $h^{*}$ for the forward difference
$D_{+}(x; h) = (f(x + h) - f(x))/h$ in the presence of round-off
error of magnitude $\varepsilon$ in each function value, by
balancing the truncation error against the round-off error.
Report $h^{*}$ and the optimal total error in terms of $|f''|$
and $\varepsilon$.
\end{exercise}

\paragraph{Solution sketch.}
The truncation error of $D_{+}$ is $\tfrac{1}{2} f''(x) h +
\oom{h^{2}}$, of order $h$. The round-off error is of order
$\varepsilon / h$ (the difference of two near-equal numbers each
known to absolute precision $\varepsilon$ divided by $h$). The
total error has magnitude $E(h) \sim \tfrac{1}{2} |f''| h +
\varepsilon / h$. Differentiate with respect to $h$ and set to
zero: $\tfrac{1}{2} |f''| - \varepsilon / h^{2} = 0$, giving
$h^{*} = \sqrt{2 \varepsilon / |f''|}$. Substituting,
$E(h^{*}) \sim \sqrt{2 \varepsilon |f''|}$, of order
$\varepsilon^{1/2}$. With double-precision $\varepsilon \approx
10^{-16}$ and $|f''| \sim 1$, $h^{*} \approx 10^{-8}$ and the
optimal error is about $10^{-8}$.

\begin{exercise}
Derive Newton's iteration for finding the root of a function $g$
by linearising $g$ at the current iterate $x_{n}$ and setting the
linearisation to zero. Report the iteration $x_{n + 1} = x_{n} -
g(x_{n})/g'(x_{n})$ and identify the conditions under which the
iteration converges.
\end{exercise}

\paragraph{Solution sketch.}
The linear approximation of $g$ at $x_{n}$ is $L_{x_{n}}(x) =
g(x_{n}) + g'(x_{n}) (x - x_{n})$. Setting $L_{x_{n}}(x) = 0$ and
solving for $x$ gives $x_{n + 1} = x_{n} - g(x_{n})/g'(x_{n})$,
provided $g'(x_{n}) \neq 0$. The iteration converges quadratically
when (i) $g'$ is continuous and nonzero in a neighbourhood of the
root $x^{*}$, and (ii) the initial guess $x_{0}$ is close enough
to $x^{*}$ that the iterates stay in the neighbourhood. The
quadratic-convergence statement is $|x_{n+1} - x^{*}| \leq C
|x_{n} - x^{*}|^{2}$ with $C$ proportional to $|g''(x^{*})| /
|g'(x^{*})|$. Failure modes include $g'(x_{n}) \approx 0$ (the
iterate flies far from the root), the initial guess being too
far from $x^{*}$ (the iteration can wander chaotically), and
multiple roots where $g'$ vanishes (the convergence becomes
linear or breaks).

\subsection*{Estimation}\pathtag{core}

\begin{exercise}
Estimate the rate of change of the kilometre-readout on a car's
odometer for a car cruising at $90\,\si{\kilo\meter\per\hour}$.
Report your estimate in odometer-units per second, then propagate
to the derivative the speedometer would read in metres per second.
\end{exercise}

\paragraph{Solution sketch.}
Convert: $90 \,\si{\kilo\meter\per\hour} = 90 / 3600 \,\si{\kilo
\meter\per\second} = 0.025 \,\si{\kilo\meter\per\second}$. The
odometer reads in kilometres; its time derivative is the speed in
kilometres per second, $0.025 \,\si{\kilo\meter\per\second}$. The
speedometer reads in metres per second; multiply by $1000$:
$25 \,\si{\meter\per\second}$. The two readings are the same
quantity in different units; the chain rule converts between
them.

\begin{exercise}
Estimate the maximum acceleration a person can comfortably tolerate
for a few seconds, in units of $g \approx 9.8\,\si{\meter\per
\second\squared}$. Use the result to estimate the minimum
stopping distance of a passenger car at $100\,\si{\kilo\meter\per
\hour}$.
\end{exercise}

\paragraph{Solution sketch.}
Comfortable deceleration for a few seconds is about $0.5 g \approx
5 \,\si{\meter\per\second\squared}$ (firm braking; a panic stop
in a passenger car can exceed $1 g$ but is uncomfortable). At
$100 \,\si{\kilo\meter\per\hour} \approx 27.8 \,\si{\meter\per
\second}$, deceleration at $5 \,\si{\meter\per\second\squared}$
brings the car to rest in $27.8 / 5 \approx 5.6 \,\si{\second}$.
The stopping distance is the integral $\tfrac{1}{2} v_{0}^{2} / a
= \tfrac{1}{2} (27.8)^{2} / 5 \approx 77 \,\si{\meter}$. Real
panic-stop distances (with a higher peak deceleration but a
reaction-time delay) are also in the $75$ to $100 \,\si{\meter}$
range.

\begin{exercise}
Estimate the rate at which the surface temperature of a hot cup
of tea decreases at room temperature in the first minute of
cooling. Use Newton's law of cooling with reasonable assumptions
about the heat-transfer coefficient and the surface area.
\end{exercise}

\paragraph{Solution sketch.}
Tea at $80 \,\si{\degreeCelsius}$ in a $20 \,\si{\degreeCelsius}$
room has $\Delta T = 60 \,\si{\degreeCelsius}$. The chapter's
worked cooling problem gave $k \approx 0.055 \,\si{\per\minute}$
for a similar mug. Newton's law gives $\dv{T}{t}\big|_{t=0} =
-k \Delta T = -0.055 \cdot 60 \approx -3.3 \,\si{\degreeCelsius
\per\minute}$, or about $-0.055 \,\si{\degreeCelsius\per\second}$.
Over the first minute the temperature drops about $3 \,\si{
\degreeCelsius}$ (the rate falls slightly as the temperature
difference narrows, but linearly to good approximation for the
first minute). A reader who has measured a mug confirms the
$3 \,\si{\degreeCelsius}$ per minute scale.

\begin{exercise}
Estimate the rate of change of the tide height at a coastal
location at the time the tide is rising fastest. Use the height
range, the period, and the assumption of approximately sinusoidal
tides. State whether your estimate is the peak rate or a typical
rate.
\end{exercise}

\paragraph{Solution sketch.}
Take a typical tidal range $\Delta H = 2 \,\si{\meter}$ and a
period $T = 12.4 \,\si{\hour}$ (semidiurnal). The sinusoidal
height $H(t) = (\Delta H / 2) \cos(2 \pi t / T) + H_{0}$ has
peak rate $|H'|_{\max} = (\Delta H / 2) \cdot 2 \pi / T = \pi
\Delta H / T = \pi \cdot 2 / 12.4 \approx 0.51 \,\si{\meter\per
\hour}$. This is the peak rate, achieved when the tide is at
half-amplitude (mid-flood and mid-ebb). The rate averaged over a
full cycle is zero; the rate averaged over the rising half is
$\Delta H / (T/2) = 4 / 12.4 \approx 0.32 \,\si{\meter\per\hour}$.
At extreme locations (Bay of Fundy, around $15 \,\si{\meter}$
range), the peak rate exceeds $3 \,\si{\meter\per\hour}$.

\begin{exercise}
Estimate the optimal step size for numerical differentiation of
a measured signal whose noise standard deviation is $10^{-2}$
times the signal amplitude, assuming the underlying signal has
$|f'''| \sim 1$ in the units of the problem. Compare your
estimate to the double-precision floating-point optimum.
\end{exercise}

\paragraph{Solution sketch.}
For a noisy signal with noise standard deviation $\sigma$, the
round-off-equivalent error floor in the central difference is
$\sigma / h$. Balancing against the truncation error $|f'''| h^{2}
/ 6$ gives $h^{*} \sim (\sigma / |f'''|)^{1/3}$. With
$\sigma = 10^{-2}$ and $|f'''| \sim 1$, $h^{*} \approx
(10^{-2})^{1/3} \approx 0.22$. The optimal derivative error is
of order $\sigma^{2/3} \approx 10^{-4/3} \approx 0.05$. The
double-precision floating-point optimum has $\varepsilon \sim
10^{-16}$ and gives $h^{*} \sim 10^{-5}$ with optimum error
about $10^{-11}$. Measurement noise dominates floating-point
round-off by roughly fourteen decades in this case; the engineer
should pick $h$ against the measurement noise, not against the
floating-point $\varepsilon$.

\subsection*{Simulation}\pathtag{mastery}

\begin{exercise}
In a language of the reader's choice, generate a sinusoidal signal
$y(t) = \sin(2\pi t)$ at $200$ samples per period over five
periods, contaminate the signal with Gaussian noise of standard
deviation $\sigma = 0.05$, and compute the derivative by central
differences at each interior sample. Compare to the true
derivative $y'(t) = 2\pi \cos(2\pi t)$, and report the root-mean-
square error of the numerical derivative.
\end{exercise}

\paragraph{Solution sketch.}
The companion script \texttt{code/noisy\_derivative.py} performs
the calculation. Sample spacing $\Delta t = 1/200 = 0.005$, so
the noise-induced floor of the central-difference derivative is
$\sigma / \Delta t = 0.05 / 0.005 = 10$, an order of magnitude
larger than the true derivative amplitude $2\pi \approx 6.3$. The
naive estimate is therefore dominated by noise. Sub-sampling to
stride $8$ gives effective $\Delta t = 0.04$ and noise floor
$\sigma / \Delta t \approx 1.25$, comparable to the signal.
Reported root-mean-square errors at strides $1, 2, 4, 8, 16, 32$
fall from order $10$ to order $1$ and then begin to rise as the
truncation bias becomes visible at the largest strides.

\begin{exercise}
For the same signal as above, compute the derivative by smoothing
with a Savitzky-Golay filter of polynomial order $3$ and window
widths $w = 5, 11, 21, 41$. Plot the root-mean-square error of the
derivative against window width, and identify the optimum.
\end{exercise}

\paragraph{Solution sketch.}
A Savitzky-Golay filter of polynomial order $p$ and window width
$w$ fits a degree-$p$ polynomial to a sliding window of $w$ samples
and reports the derivative of the polynomial at the window's
centre \cite{text:v2ch05-savitzky-golay-1964}. For order $3$ and
the given signal, narrow windows ($w = 5$) preserve the signal but
admit noise; wide windows ($w = 41$) suppress noise but smear the
$\sin(2\pi t)$ structure. The root-mean-square error against the
true $2\pi \cos(2\pi t)$ typically shows a minimum near $w = 11$
or $w = 21$, depending on $\sigma$. The reader should plot the
error curve and report the value at the minimum.

\begin{exercise}
Implement Newton's iteration for solving $g(x) = 0$ where
$g(x) = x^{3} - 2 x - 5$. Start from $x_{0} = 2$, run for ten
iterations, and report the convergence rate empirically.
Compare to the theoretical quadratic convergence.
\end{exercise}

\paragraph{Solution sketch.}
The companion script \texttt{code/newton\_root.py} runs the
iteration. The reference root is $x^{*} \approx 2.0945514815$.
From $x_{0} = 2$ the iteration produces $x_{1} = 2.1$,
$x_{2} \approx 2.0946$, $x_{3} \approx 2.09455148\ldots$, with
each successive error about the square of the previous error.
The number of correct decimal digits roughly doubles per step,
matching the theoretical quadratic-convergence statement
$|x_{n+1} - x^{*}| \leq C |x_{n} - x^{*}|^{2}$. The constant $C$
is approximately $|g''(x^{*})| / (2 |g'(x^{*})|) =
|6 x^{*}|/(2(3 (x^{*})^{2} - 2)) \approx 12.6 / 21.1 \approx 0.6$.

\begin{exercise}
Sweep the step size $h$ over the range $h = 10^{-1}, 10^{-2},
\ldots, 10^{-15}$ for the central-difference approximation to
$f'(1)$ where $f(x) = e^{x}$. Plot the absolute error against $h$
on log-log axes and identify the regime where truncation error
dominates and the regime where round-off dominates. Estimate the
optimal $h$ from the plot.
\end{exercise}

\paragraph{Solution sketch.}
The companion script \texttt{code/finite\_diff\_accuracy.py}
produces the table. The exact derivative is $e^{1} \approx 2.71828$.
The truncation term is $\tfrac{1}{6} f'''(1) h^{2} = \tfrac{1}{6}
e \cdot h^{2}$; the round-off term is of order $\varepsilon / h$
with $\varepsilon \approx 10^{-16}$. The two cross at $h^{*} \sim
(\varepsilon / |f'''|)^{1/3} \sim 10^{-5.5}$ with optimal error of
order $10^{-11}$. The plot has the characteristic V-shape of
\Cref{fig:vol02:ch05:bias-variance}, with the truncation arm
falling at slope $2$ on the right and the round-off arm rising
at slope $-1$ on the left. The empirical optimum is within a
factor of two of the predicted $h^{*}$.

\subsection*{Design}\pathtag{standard}

\begin{exercise}
Design a numerical-differentiation strategy for the velocity of a
vehicle measured by GPS at $1\,\si{\hertz}$ with horizontal
position uncertainty of $5\,\si{\meter}$ per sample. State the
filter or step size you would use, the resulting expected
uncertainty in the velocity estimate, and the regime (constant
velocity, accelerating, turning) in which your strategy is
defensible.
\end{exercise}

\paragraph{Discussion.}
With sample spacing $\Delta t = 1 \,\si{\second}$ and position
uncertainty $\sigma_{p} = 5 \,\si{\meter}$, a single-sample
central difference yields velocity uncertainty $\sigma_{p} /
\Delta t \cdot \sqrt{2} / (2) = 5 / 1 \cdot \sqrt{2}/2 \approx
3.5 \,\si{\meter\per\second}$ per sample, or about $13 \,\si
{\kilo\meter\per\hour}$. A stride-five central difference
($\Delta t_{\text{eff}} = 5 \,\si{\second}$) reduces this to
about $0.7 \,\si{\meter\per\second}$ but introduces a $5 \,\si
{\second}$ lag in the estimate. A defensible default for
quasi-constant-velocity cruising is a five-to-ten sample
central window or a Savitzky-Golay filter of polynomial order
$2$ with $w = 11$. The strategy is defensible for constant
velocity and gentle accelerations (the polynomial fits the
short-window trajectory cleanly); it fails for sharp turns or
abrupt accelerations, where the smoothing smears the manoeuvre.
A modern receiver typically uses Doppler-based velocity at
higher accuracy when available; the position-derivative path is
the fallback.

\begin{exercise}
Design a tank optimisation: a cylindrical hot-water storage tank
must hold a fixed volume $V$ and minimise heat loss through its
walls. Assume the heat loss per unit area is constant and that
the tank is open-topped (no top surface). Find the optimal radius
and height ratio. Compare your answer to the closed-cylinder case
in the chapter's worked example, and explain the qualitative
difference.
\end{exercise}

\paragraph{Discussion.}
The open-topped surface area is $A(r) = \pi r^{2} + 2 \pi r h$
with constraint $\pi r^{2} h = V$, so $h = V/(\pi r^{2})$ and
$A(r) = \pi r^{2} + 2 V / r$. Then $A'(r) = 2 \pi r - 2 V / r^{2}$,
vanishing at $r^{*} = (V/\pi)^{1/3}$. At the optimum, $h^{*} =
V/(\pi (r^{*})^{2}) = r^{*}$: the open tank is as tall as it is
wide. The closed-cylinder case (the chapter's aluminium can) had
$h^{*} = 2 r^{*}$. The qualitative difference: without a top, the
penalty for adding height is lower (no extra top-circle area), so
the optimum lengthens the cylinder relative to the closed case;
without a top, the penalty for adding radius is also lower
(the top circle no longer scales with $r^{2}$), so the optimum
shrinks the radius in the right way to keep the side area in
balance. The two opposing effects combine to give $h^{*} = r^{*}$
in the open case versus $h^{*} = 2 r^{*}$ in the closed case.

\begin{exercise}
Design a one-variable optimisation for the resistance of an
external load that maximises power transfer from a source with
internal resistance $R_{i}$ and an additional in-line wire
resistance $R_{w}$. Find the optimal load resistance and the
maximum power transferred.
\end{exercise}

\paragraph{Discussion.}
The source delivers current $I = \mathcal{E} / (R + R_{w} + R_{i})$
where $R$ is the load. Power dissipated in the load is
$P(R) = I^{2} R = \mathcal{E}^{2} R / (R + R_{w} + R_{i})^{2}$.
Setting $P'(R) = 0$ by the quotient rule gives $(R + R_{w} +
R_{i})^{2} = 2 R (R + R_{w} + R_{i})$, that is, $R^{*} = R_{w} +
R_{i}$: the optimal load matches the total of the source's
internal resistance plus the in-line wire resistance. The
maximum power transferred is $P^{*} = \mathcal{E}^{2} /
(4 (R_{w} + R_{i}))$, and the fraction of source power that
reaches the load is $R^{*} / (R^{*} + R_{w} + R_{i}) = 1/2$.
The qualitative lesson: load matching is to the entire upstream
impedance, not only to the source's internal impedance. The wire
loss is irrecoverable and shrinks the deliverable maximum.

\subsection*{Diagnosis and reverse engineering}\pathtag{core}

\begin{exercise}
A spreadsheet computes $\dv{y}{x}$ for a measured time series
using $D_{+}(x; h) = (y(x + h) - y(x))/h$ with the smallest
available $h$, and reports a derivative estimate that fluctuates
wildly between adjacent samples. Identify the most likely cause
of the fluctuation and propose a correction.
\end{exercise}

\paragraph{Solution sketch.}
The smallest available $h$ is the sample spacing, which places
the estimator on the noise-dominated branch of the bias-variance
tradeoff. The derivative carries amplified measurement noise
(noise standard deviation in $y$ becomes $\sigma_{y} / h$ in the
derivative). The fix is to sub-sample (use a larger effective
$h$ by skipping samples), to smooth the data before
differentiation, or to fit a parametric model and differentiate
it. Each option moves the estimator off the noise branch and
toward the truncation branch where the tradeoff is balanced.

\begin{exercise}
A control system computes the derivative of a process variable
by central differences and uses the derivative in a feedback
loop. The system is unstable. Identify three plausible numerical-
differentiation defects that would produce instability and
propose a diagnostic test for each.
\end{exercise}

\paragraph{Solution sketch.}
(1) Step size too small: the central difference is on the
round-off or noise branch, amplifying sensor noise into the
control loop. Test: report the derivative noise floor at several
step sizes; the floor should fall as $h$ grows until the
truncation branch takes over. (2) Group delay from a smoothing
filter applied before differencing: the filtered derivative is
lagged relative to the measurement, and the feedback adds the
lag to its phase margin. Test: bode-plot the input-to-output
phase versus frequency and confirm the delay is within the
loop's tolerance. (3) Quantisation of the measurement: a sensor
with a coarse least-significant bit produces a derivative that
is mostly zero with occasional large jumps; the loop alternates
between no correction and over-correction. Test: histogram the
sequential differences $y_{n+1} - y_{n}$; a discrete sensor
shows up as a discrete distribution.

\begin{exercise}
A chemistry laboratory computes a reaction rate by differencing
two concentration measurements taken at $t_{1}$ and $t_{2}$. The
reported rate has a standard deviation $50\,\%$ of its mean.
Identify the source of the noise and propose a measurement
strategy that reduces it.
\end{exercise}

\paragraph{Solution sketch.}
A two-point forward difference of two noisy concentration
measurements has noise standard deviation $\sigma_{r} =
\sigma_{C} \sqrt{2} / (t_{2} - t_{1})$. A $50\,\%$ relative
uncertainty in the rate means the measurement noise is
comparable in magnitude to the concentration change over the
interval $(t_{2} - t_{1})$. The fix is to widen the interval
(so the true change $\Delta C$ grows while the noise stays the
same), to repeat the measurement and average (so the noise per
measurement is reduced by $\sqrt{N}$), or to fit a kinetic model
to a longer time series and read the rate from the fit. The
right choice depends on whether the reaction is slow enough to
admit a longer interval without nonlinearity (in which case
widen the interval) or fast enough to require averaging
(in which case multiplexing or higher-precision instrumentation
is the route).

\subsection*{Failure analysis}\pathtag{core}

\begin{exercise}
The Patriot missile system at Dhahran (1991) failed to engage an
incoming Scud because of an accumulated software-clock error of
about $0.34\,\si{\second}$ over $100\,\si{\hour}$ of continuous
operation \cite{acc:gao-patriot-1992}. Argue, in 200 to 300 words,
how the per-conversion error of about $9.5 \times 10^{-8}$
accumulated to a tracking failure, and identify the engineering
discipline (numerical-precision audit, restart procedure,
operational guidance) that would have caught the defect before
deployment.
\end{exercise}

\paragraph{Solution sketch.}
The conversion of the integer tick counter to real seconds used a
$24$-bit fixed-point approximation of $1/10$ whose binary
expansion is non-terminating; the truncation introduced about
$9.5 \times 10^{-8}$ of error per conversion. After
$100 \,\si{\hour} = 3.6 \times 10^{5} \,\si{\second}$ the counter
had reached approximately $3.6 \times 10^{6}$ ticks, and the
accumulated error reached approximately $0.34 \,\si{\second}$.
The range-gate prediction used the converted time in a
velocity-times-time forecast; at $1{,}676 \,\si{\meter\per
\second}$ the time error mapped to a position error of about
$570 \,\si{\meter}$, larger than the range gate, and the system
declined to engage. Three disciplines would each have caught the
defect. A numerical-precision audit at design time, by
propagating the per-conversion error over the maximum intended
continuous-operation duration, would have flagged the
accumulated error and forced either a higher-precision
representation, periodic reset, or different clock-handling
algorithm. A restart procedure in operational doctrine, with a
hard interval well below the accumulated-error threshold (say,
restart every $24 \,\si{\hour}$), would have kept the system
inside its design envelope regardless of the software defect. An
operational guidance system that synchronised patches with
operational doctrine (the patch and the restart guidance moved
through different distribution channels) would have ensured
every deployed battery had at least one of the two protections
in place at any time.

\begin{exercise}
A consumer thermostat reports the rate of change of the
temperature it is controlling. The thermostat's sensor has a
quantisation step of $0.1\,\si{\degreeCelsius}$ and reports the
sensor value once per minute. Argue, in 200 to 300 words,
whether the thermostat's reported rate of change is meaningful at
the minute scale, and identify the engineering correction that
would make it meaningful.
\end{exercise}

\paragraph{Solution sketch.}
The minute-scale derivative computed by differencing two
consecutive sensor readings has a noise floor equal to the
quantisation step divided by the time interval, that is, $0.1 /
60 \approx 1.7 \times 10^{-3} \,\si{\degreeCelsius\per\second}$
or $0.1 \,\si{\degreeCelsius\per\minute}$. A real household
heating system changes temperature at perhaps
$0.5$ to $5 \,\si{\degreeCelsius\per\hour} = 0.008$ to $0.08
\,\si{\degreeCelsius\per\minute}$. The reported rate is therefore
dominated by quantisation noise at the minute scale; the
reported number alternates between $0$ and $\pm 0.1 \,\si{
\degreeCelsius\per\minute}$ depending on whether the sensor
flipped a bit during the minute. The fix is to widen the
differencing interval (report a ten-minute or hourly rate,
matching the system's actual time constant), to fit a slope to
a sliding window of (say) ten samples and report the slope as
the rate (Savitzky-Golay or weighted least squares), or to
increase the sensor's effective resolution by oversampling and
averaging (if the sensor admits faster reads). The minute-scale
reported rate is not informative without one of these
corrections; the temperature is changing slowly compared to the
quantisation noise floor.

\begin{exercise}
A small-amplitude pendulum is modelled by simple harmonic motion
with $\theta(t) = \theta_{0} \cos(\omega t)$, where $\omega =
\sqrt{g/L}$. The model breaks for large amplitudes. Identify the
amplitude scale at which the model's prediction of the period
deviates from reality by $5\,\%$ (chapter~4 of this volume gives
the leading-order correction). Identify the failure mode of the
linearised derivative-based analysis at this amplitude.
\end{exercise}

\paragraph{Solution sketch.}
The leading nonlinear correction to the pendulum period is
$T = T_{0} (1 + \theta_{0}^{2} / 16 + \cdots)$. A $5\,\%$
deviation corresponds to $\theta_{0}^{2}/16 = 0.05$, that is,
$\theta_{0}^{2} = 0.8$, $\theta_{0} \approx 0.89 \,\si{\radian}
\approx 51^{\circ}$. At this amplitude, the linear-approximation
$\sin\theta \approx \theta$ used to derive simple harmonic
motion fails by a margin that the period detects. The failure
mode of the derivative-based linearisation is that the
restoring-force differential equation $\ddot\theta = -(g/L)\sin
\theta$ is approximated by $\ddot\theta = -(g/L)\theta$,
discarding the higher-order Taylor terms of $\sin\theta$. The
discarded terms produce a period that depends on amplitude (the
isochronism of simple harmonic motion is itself a
linearisation artefact), and the working engineer who needs
sub-percent timing accuracy in a wide-swing pendulum must use
the elliptic-integral expression for $T$ or work with the
$5\,\%$-or-larger amplitude error in mind.

\subsection*{Judgment}\pathtag{standard}

\begin{exercise}
A colleague reports the derivative of a noisy measurement series
without specifying the differentiation method. Write a one-
paragraph response that takes the colleague seriously, identifies
the specific information missing to interpret the result, and
proposes a minimum acceptable form for the report.
\end{exercise}

\paragraph{Discussion.}
The colleague has reported a single number whose uncertainty band
is undefined. To interpret the derivative, the reader needs the
differencing scheme (forward, central, higher order), the
effective step size (raw sample spacing or a sub-sampled value),
any smoothing applied before differencing (filter type, window
width, polynomial order), and the noise standard deviation of
the underlying measurement. A minimum acceptable report names the
scheme, the effective step size, and the propagated noise floor
of the derivative; for measured data the report adds the spread
across at least two independent methods (a finite difference and
a smoothed-then-differentiated estimate, for example) as a
working uncertainty band. The response should be collegial: the
colleague's derivative is probably right; the gap is documentation
discipline.

\begin{exercise}
A regulatory agency requires a manufacturer to report the rate
of change of a quantity but specifies neither the time scale nor
the differentiation method. Argue, in 200 to 300 words, whether
the regulation is well-posed, and propose a specification that
makes the requirement testable.
\end{exercise}

\paragraph{Discussion.}
The regulation is not well-posed. A "rate of change" computed
over a millisecond is a different quantity from one computed over
an hour, even on the same underlying signal; the two can differ
by orders of magnitude when the signal carries high-frequency
components. Two manufacturers reporting their honest best
estimates can produce numbers that disagree by factors that
depend entirely on undisclosed implementation choices. A testable
specification names: the time window over which the rate is
computed (an averaging interval); the differentiation method (a
specific formula, including any smoothing filter and its
parameters); the sample rate of the underlying measurement; the
procedure for handling edge effects at the start and end of the
window; and the accepted uncertainty band of the reported rate.
With these, two laboratories using two manufacturers' instruments
produce comparable numbers, and the regulator's tolerance is a
real test rather than a contest between implementations. The
discipline mirrors the metrology tradition of Volume I: a
quantity that is not traceable to a stated method is not a
quantity that can be compared.

\begin{exercise}
A business publishes a chart showing the "rate of growth" of a
quantity over time, computed by a method the reader cannot
verify. Argue, in 200 to 300 words, what the chart cannot tell
the reader without disclosure of the differentiation method, and
identify the reader's defensive options.
\end{exercise}

\paragraph{Discussion.}
Without disclosure, the chart cannot tell the reader (i) whether
the displayed rate is an instantaneous derivative or an
averaged-over-some-window quantity, (ii) whether the smoothing
choice favours visual continuity over honest reporting of
volatility, (iii) whether the rate is computed on the raw data
or on a transformed quantity (logarithm, percent change), and
(iv) whether the visual scale (linear or log axis on the rate
itself) understates or overstates the swings the underlying
signal exhibits. Two competent analysts presented with the same
raw data can produce charts that tell visually different stories
by varying these choices. The reader's defensive options
include: ask for the raw data and reproduce the chart with
disclosed methods; ask for the differentiation method and an
uncertainty band on the displayed rate; treat the chart as
illustrative rather than authoritative pending disclosure;
compare the chart to a chart of the underlying quantity itself,
where the visual integration removes some implementation
artefacts. The general discipline: a rate is a derived quantity
that depends on choices, and a rate published without those
choices stated is a rate without a defensible interpretation.

\begin{exercise}
Argue, in one paragraph, whether engineering training should
emphasise analytical or numerical differentiation. Take a
position and defend it with a specific example from your own
experience or a documented case.
\end{exercise}

\paragraph{Discussion.}
The training has to develop both, in that order. Analytical
differentiation gives the engineer the structure of how a
quantity responds to its input: the chain rule reveals the
multiplicative composition of sensitivities; the product and
quotient rules reveal where small denominator errors blow up
into large derivative errors; Taylor expansion gives the
remainder bound that lets the engineer know when an
approximation is admissible. Numerical differentiation gives the
engineer the discipline of finite precision: the bias-variance
tradeoff, the noise floor, the choice of step size, the
catastrophic-cancellation failure mode. A training that provides
only the analytical side produces engineers who derive a
beautiful formula and apply it to floating-point data with no
sense of why the result is garbage; a training that provides
only the numerical side produces engineers who treat
differentiation as a black-box subroutine and miss the
structural information the analytical view delivers. The Patriot
case is the canonical illustration: the failure was numerical,
but the diagnosis required understanding that the differenced
quantity inherited the full accumulated representation error of
each input, which is a structural insight from analytical
calculus applied to a finite-precision computation.

\begin{exercise}
A sensor manufacturer states the noise floor of a measurement at
$\sigma = 10^{-3}$ in the measured units, with samples available
at $f_{s}$ samples per second. Write the differentiation strategy
for the smallest defensible derivative the manufacturer could
quote, including the optimal step size and the resulting
derivative noise floor.
\end{exercise}

\paragraph{Discussion.}
For a central difference with effective step size $h$ and noise
$\sigma$, the noise contribution to the derivative is $\sigma /
h$, and the truncation contribution is $|f'''| h^{2}/6$. The
optimal step size is $h^{*} \sim (\sigma / |f'''|)^{1/3}$ and
the optimal derivative noise floor is $\sigma^{2/3}
|f'''|^{1/3}$ times a constant of order unity. With $\sigma =
10^{-3}$ and $|f'''| \sim 1$ in the measurement units, $h^{*}
\approx 10^{-1}$ in the time units of the measurement: the
smallest defensible step uses on the order of $0.1 \cdot f_{s}$
samples. The derivative noise floor is approximately
$\sigma^{2/3} \approx 10^{-2}$ in (measured units per time
unit). The manufacturer should quote the signal-aware derivative
rather than the instrument-resolution-limited one and should
state the assumed $|f'''|$ scale, since the optimum depends on
it.

\begin{exercise}
Compute the derivative of $f(x) = \ln(\sqrt{x^{2} + 1})$ by
applying the chain rule and the logarithmic-derivative
manipulation. Compare the result to the answer obtained by first
simplifying $\ln(\sqrt{x^{2} + 1}) = \tfrac{1}{2}\ln(x^{2} + 1)$.
\end{exercise}

\paragraph{Solution sketch.}
By the chain rule with outer $u \mapsto \ln u$, inner $v \mapsto
\sqrt{v}$, innermost $x \mapsto x^{2} + 1$:
$f'(x) = (1/\sqrt{x^{2} + 1}) \cdot (1/(2\sqrt{x^{2} + 1})) \cdot
2 x = x / (x^{2} + 1)$. By the simplification, $f(x) =
\tfrac{1}{2} \ln(x^{2} + 1)$ and $f'(x) = \tfrac{1}{2} \cdot
(2 x)/(x^{2} + 1) = x/(x^{2} + 1)$. The two routes agree, as
they must. The simplification is shorter; the chain-rule route
is more general and survives when the logarithm-of-square-root
structure is not amenable to algebraic simplification.

\begin{exercise}
Compute, from the limit definition, the derivative of $f(x) =
\sqrt{x}$ at an arbitrary point $x > 0$. Use the conjugate-
multiplication trick to simplify the difference quotient
$(\sqrt{x + h} - \sqrt{x})/h$ before taking the limit.
\end{exercise}

\paragraph{Solution sketch.}
Multiply numerator and denominator by the conjugate $\sqrt{x + h}
+ \sqrt{x}$:
\[
\frac{\sqrt{x + h} - \sqrt{x}}{h} = \frac{(\sqrt{x + h} - \sqrt{x})
(\sqrt{x + h} + \sqrt{x})}{h (\sqrt{x + h} + \sqrt{x})} =
\frac{(x + h) - x}{h (\sqrt{x + h} + \sqrt{x})} =
\frac{1}{\sqrt{x + h} + \sqrt{x}}.
\]
As $h \to 0$, the limit is $1/(2 \sqrt{x})$. So $f'(x) = 1/(2
\sqrt{x})$ for $x > 0$, the same as the power rule applied to
$x^{1/2}$.

\begin{exercise}
A photochemical reaction has rate $r = k [A]^{2} [B]$ where $k$
is constant and $[A], [B]$ are concentrations measured with
relative uncertainties $u([A])/[A] = 0.05$ and $u([B])/[B] =
0.10$. Use the local-linearisation form of uncertainty
propagation (the GUM general formula, now legitimised by
this chapter) to compute the relative uncertainty in $r$ assuming
$[A]$ and $[B]$ are uncorrelated.
\end{exercise}

\paragraph{Solution sketch.}
For a product of powers $r = k [A]^{p} [B]^{q}$ the GUM formula
for uncorrelated inputs gives $(u(r)/r)^{2} = p^{2}
(u([A])/[A])^{2} + q^{2} (u([B])/[B])^{2}$. With $p = 2$,
$q = 1$, $u([A])/[A] = 0.05$, $u([B])/[B] = 0.10$:
$(u(r)/r)^{2} = 4 \cdot 0.0025 + 1 \cdot 0.01 = 0.01 + 0.01 =
0.02$, so $u(r)/r \approx 0.141$, about $14\,\%$. The squared
exponent on $[A]$ amplifies its relative uncertainty by a factor
of two; the larger relative uncertainty in $[B]$ enters
unamplified. The two contributions are comparable; reducing
either improves the result, with $[A]$ giving the larger leverage
because of the squared exponent.

\begin{exercise}
For the function $f(x) = \tan(x)$, identify all points in the
interval $(-2\pi, 2\pi)$ where the derivative does not exist.
For each such point, classify the failure (jump, kink, vertical
tangent, oscillation). Write a one-paragraph engineering
implication: a control system that uses $\tan$ as part of its
gain schedule must guard against operation near these points.
\end{exercise}

\paragraph{Solution sketch.}
The function $\tan(x)$ has vertical asymptotes at the odd
multiples of $\pi/2$ in the given interval: $x = -3\pi/2$,
$-\pi/2$, $\pi/2$, $3\pi/2$. At each such point $\tan$ itself is
undefined, and the derivative $\sec^{2} x$ diverges to
$+\infty$ as $x$ approaches an asymptote from either side. The
failure type is unbounded derivative (vertical-tangent character,
with the function value itself missing), distinct from a jump
(the function is unbounded near the asymptote) and from a kink
(the function is smooth where defined). A control system whose
gain schedule includes $\tan$ must clip the input away from
these asymptotes; the gain becomes arbitrarily large near the
asymptotes and the loop becomes unstable. A defensive design uses
$\arctan$, which is bounded and has bounded derivative
everywhere, or applies a saturating gain that caps the
amplification before it reaches the asymptote.

\begin{exercise}
Estimate the central-difference derivative of a smartphone
accelerometer signal sampled at $100\,\si{\hertz}$ with noise
standard deviation $\sigma = 0.01\,\si{\meter\per\second\squared}$
in each sample. State the sample spacing at which the central-
difference noise floor in the derivative is comparable to a
$0.5\,\si{\meter\per\second\cubed}$ jerk signal, and propose a
filtering strategy that recovers a usable jerk estimate.
\end{exercise}

\paragraph{Solution sketch.}
At native sample interval $\Delta t = 10 \,\si{\milli\second}$
the central-difference noise floor in the jerk estimate is
$\sigma / \Delta t = 0.01 / 0.01 = 1.0 \,\si{\meter\per
\second\cubed}$, above the target signal. To reach a noise floor
of $0.5 \,\si{\meter\per\second\cubed}$ comparable to the
signal, the effective sample spacing must grow to $\Delta t =
\sigma / 0.5 = 0.02 \,\si{\second} = 20 \,\si{\milli\second}$,
that is, a stride-two central difference. At this spacing the
truncation bias depends on the signal's third derivative and on
the size of the manoeuvre; for a smooth driving scenario it
remains small. A defensible strategy uses a Savitzky-Golay
filter of polynomial order $2$ with window width $11$ to $21$
samples, producing a derivative whose noise floor is roughly
$\sigma / (\sqrt{w/2} \cdot \Delta t) \approx 0.3 \,\si{
\meter\per\second\cubed}$, comparable to the target signal,
with a usable $100 \,\si{\milli\second}$ time resolution.
